% 本文件由 LaTeX 排版 Agent 根据有效 translation.json 自动生成。
% author=John of Damascus; work_id=3304; title=正统信仰阐详

\chapter{正统信仰阐详}

% structure: p0001 source=h1 target=omitted reason=page-title-already-rendered-by-parent

卷一\newline{}卷二\newline{}卷三\newline{}卷四\par

\SourceRecord{Translated by E.W. Watson and L. Pullan.；Nicene and Post-Nicene Fathers, Second Series；Vol. 9.；Edited by Philip Schaff and Henry Wace.；Buffalo, NY: Christian Literature Publishing Co.,；1899.；<http://www.newadvent.org/fathers/3304.htm>.}

% structure: p0001 source=h1 target=section reason=page-title

\section{\WorkTitle{正统信仰诠释（第一部）}}

% structure: p0002 source=h2 target=subsection reason=relative-heading-rank; base=h2

\subsection{第一章：论神性是不可理解的，我们不应探究和干预那未被圣先知、宗徒和圣史所传授给我们的事。}

\BibleQuote{从来没有人见过天主；只有在圣父怀里的独生子，祂彰显了祂。}因此，神性是无可言喻且不可理解的。\BibleQuote{因为除了子，没有人认识父；除了父，也没有人认识子。}\BibleRef{玛 11:27}同样，圣神也认识天主的事，就如人的灵认识人里面的事。\BibleRef{格前 2:11}此外，在第一个有福的本性之后，不仅没有人，连超世的能力，甚至革鲁宾和色辣芬自己，从未认识天主，除非祂向谁启示了自己。\par

然而，天主并没有让我们完全无知。因为对天主存在的知识，祂已藉本性植入众人心中。这受造界及其维持和管理，也宣扬了神圣本性的威严。\BibleRef{智 13:5}此外，在古时透过法律和先知，后来透过祂的独生子、我们的主、天主和救主耶稣基督，祂按我们可能的方式，向我们启示了对祂自己的知识。因此，凡藉法律、先知、宗徒和圣史传授给我们的一切，我们都接纳、认识并尊崇，不寻求超出这些的事物。因为天主是善的，是一切善的原因，既无嫉妒，也无任何情欲。因为嫉妒远离神圣本性，这本性是毫无情欲且唯独是善的。因此，祂既知悉一切，并为每个人的益处筹划，就启示了那对我们有益而可知的事；但我们不能承受的，祂就隐瞒了。让我们满足于此，并持守这些，不挪移永久的界限，也不越过神圣的传统。\BibleRef{箴 22:28}\par

% structure: p0005 source=h2 target=subsection reason=relative-heading-rank; base=h2

\subsection{第二章：论可言说与不可言说之事，以及可知与不可知之事。}

因此，凡愿谈论或聆听关于天主之事的人，必须清楚明白：无论是在关于神性的道理中，还是在关于降生成人的道理中，并非所有事都不可言说，也非所有事都可言说；并非所有事都不可知，也非所有事都可知。但可知的事属于一类，可言说的事属于另一类；正如说话是一回事，知道是另一回事。因此，许多关于天主的模糊可理解之事，无法用恰当的言辞表达，但关于超乎我们的事，我们只能按有限的能力来表达；例如，当我们谈论天主时，我们使用诸如\Emph{睡眠}、\Emph{愤怒}、\Emph{不顾，手}，以及\Emph{脚}等表达。\par

因此，我们既知道也承认：天主是无始、无终、永恒而永在的，非受造、不变、不易、单纯、无复合、无形体、不可见、不可触摸、无界限、无限、不可知、不可定义、不可理解、善、公义、万有受造之物的创造者、全能、统治万有、鉴察万有、监察万物、君主、审判者；并且天主是独一的，即一个本体；祂在三个位格中被认识并存在，即圣父、圣子、圣神；圣父、圣子、圣神在各方面都是一样，除了不受生、受生与发出之别；独生子、天主的圣言和天主，因祂的慈悲心怀，为了我们的救恩，凭天主的善意和圣神的合作，无种子而受孕，由圣神从圣童贞女、天主之母玛利亚无玷而生，并成了出自她的完人；同一位同时是完满的天主和完满的人，具有两个本性，即天主性和人性，并在两个本性中拥有理智、意志、动力和自由，总之，按各自本性的度量和比例是完满的，即对天主性和人性而言，但却是同一个复合位格；祂曾忍受饥饿、口渴和疲倦，被钉十字架，并三天经历死亡和埋葬，然后升天，祂也曾从那里来到我们中间，并将再次来临。圣经和全体圣人都为此作证。\par

但我们既不知道，也不能说出天主的本体是什么，或它如何存在于万有之中，或独生子、天主如何空虚自己，由童贞之血成了人，藉着与本性相反的另一种律法，或祂如何干脚行走在水面上。因此，我们无法说出或甚至思考关于天主的任何事，超出那已神圣地启示给我们的事，无论是藉着言语，或藉着显现，藉着旧约和新约的神圣神谕。\par

% structure: p0009 source=h2 target=subsection reason=relative-heading-rank; base=h2

\subsection{第三章：证明有天主存在。}

那么，对于接受圣经，即旧约与新约的人而言，天主的存在是毫无疑义的；对大多数希腊人也是如此。因为我们说过，天主存在的知识是自然铭刻在我们心中的。但恶者的恶行如此强有力地压倒人的本性，甚至使一些人否认天主的存在，那最愚昧、最可悲的毁灭深渊（揭示天主圣意的达味曾揭露其愚妄，说：\BibleQuote{愚人在心中说：没有天主}），因此，主的门徒和祂的宗徒们，由圣神所智慧化，并因祂的能力和恩宠行奇事，用奇迹之网将他们俘获，从无知的深渊拉出，带到认识天主的光明中。同样，他们在恩宠和德性上的继承者，即牧者和教师，既领受了圣神的光照恩宠，也常藉奇迹的能力和恩宠之言，光照在黑暗中行走的人，将迷途者领回正路。至于我们，既未领受奇迹的恩赐，也未领受教导的恩赐（因为我们因贪图享乐而使自己不配得到这些），那么，让我们就这一主题，讨论一些由恩宠的解释者们传给我们的道理，并呼求圣父、圣子、圣神。\par

一切存在的事物，要么是受造的，要么是非受造的。如果事物是受造的，那么它们也完全是可变的。因为其存在由变化而来，也必然服从于变化，无论是消亡，或是因意志行为而变成别样的。但如果事物是非受造的，那么它们也必须是全然不变的。因为彼此存在本质对立的事物，其存在模式也必然对立，即具有相反的特性：那么，谁会拒绝承认，一切存在事物，不仅是感官所及的，甚至天使本身，都服从于变化、转变和各种运动？因为属于理性世界的事物，我指天使、灵体和魔鬼，都服从于意志的变化，无论是向善进步或退步，是奋斗还是屈服；而其余事物则经历生成与毁灭、增减、质变与空间运动的变化。因此，可变的事物完全是受造的。但受造的事物必然是某个制造者的作品，而制造者不能是受造的。因为如果他是受造的，那么他也必定是由某个主体创造的，如此类推，直到我们达到某个非受造者。因此，创造者既然是非受造的，也完全是不变的。这不是天主又是什么呢？\par

甚至受造物的持续存在、保存和管理，也教导我们：确实有一位天主，祂支持、维持、保存并永远眷顾这个宇宙。因为相反的本性，如火与水、气与土，如何能彼此结合以形成一个完整的世界，并继续存在于不可分解的结合中，除非有一种全能的力量将它们联系在一起，并始终保存它们不致消散？\par

是什么赋予天上和地上的事物，以及所有在空中和水中运动的事物以秩序？或者更确切地说，是在这之前已经存在的事物，即天、地、空气以及火和水的元素？是什么将它们混合并分配？是什么使之运动并保持在其持续无阻的轨道上？难道不是这些事物的创造者，并将万物运行的法则植入万物的那位吗？那么，这些事物的创造者是谁？难道不是创造它们并使它们存在的那位吗？因为我们将不把这样的能力归于自发性。因为，假设它们的存在是由于自发性，那么使万物有序的力量又是什么？如果我们愿意，姑且承认这一点。那使它们保持并与其原初存在的规律和谐一致的力量又是什么？显然，这完全是不同于自发性的东西。这不是天主又是什么呢？\par

% structure: p0014 source=h2 target=subsection reason=relative-heading-rank; base=h2

\subsection{第四章　论天主的本性：它是不可理解的。}

那么，显然有一位天主。但祂在其本质和本性上是什么，完全是不可理解和不可知的。因为很明显，祂是非物质的。因为无限、无界限、无形、不可触摸、不可见，总之，纯一而非复合的，怎会具有形体呢？那受限制且受情欲影响的，怎会是不变的呢？那由元素组成并可分解为它们的，怎会是无情欲的呢？因为组合是冲突的开始，冲突是分离的开始，分离是解体的开始，而解体完全与天主无关。\par

再者，又怎能坚持说天主渗透并充满宇宙呢？正如圣经所说：\BibleQuote{我岂不充满天地么？主说}\BibleRef{耶 23:24}？因为一个形体渗透其他形体而不分割和被分割，不被包裹和对比，是不可能的，如同所有流体互相混合和交融一样。\par

但如果有人说形体是非物质的，如同希腊哲学家所说的第五形体（这个形体是不可能的），那么它将完全如同天一样服从运动。因为这是他们所说的第五形体的含义。那么是谁推动它呢？因为凡被动的都是被他物所动。又是谁推动那物呢？如此以至无穷，直到我们最终达到某个不动者。因为第一推动者是不动的，那就是天主。被推动者难道不是受空间限制的吗？所以，唯有天主是不动的，祂以不动推动宇宙。因此必须假定天主是非物质的。\par

但即使这样说——谓祂是非受生的、无始的、不变的、不朽坏的，并且具有我们通常归于天主及其环境的其他属性——也并未真正表达祂的本质。因为这些并不表明祂是什么，只表明祂不是什么。当我们想要解释某物的本质为何时，不能只作否定的陈述。然而，就天主而言，不可能解释祂在其本质上是什么，我们更适合讨论祂与万物的绝对分离。因为祂不属于存在物之类：并非祂不存在，而是祂超越一切存在物，甚至超越存在本身。因为如果一切知识都与存在物相关，那么那超越知识的必然也超越本质；反之，那超越本质的也必超越知识。\par

天主的无限和不可理解，而关于祂所能理解的一切就是祂的无限和不可理解。但我们能对天主作的一切肯定陈述，并不显示天主的本性，而只是显示祂本性的属性。因为当你称祂为善、公义、智慧等等时，你并没有说出天主的本性，只是说出祂本性的属性。此外，有些我们对天主的肯定陈述具有绝对否定的效力：例如，当我们指称天主为\Quote{黑暗}时，我们并非指黑暗本身，而是说祂不是光，且超越光；当我们称祂为\Quote{光}时，我们是在说祂不是黑暗。\par

% structure: p0020 source=h2 target=subsection reason=relative-heading-rank; base=h2

\subsection{第五章　证明天主是独一而非多位。}

我们已经充分证明了天主存在，且祂的本质是不可理解的。但天主是独一而非多位，这对于相信圣经的人来说，是毫无疑义的。因为主在律法的开端说：\BibleQuote{我是上主你的天主，曾将你从埃及地领出来。除我以外，你不可有别的神。}\BibleRef{出 20:2} 祂又说：\BibleQuote{以色列，你要听：上主我们的天主，是唯一的上主。}\BibleRef{申 6:4} 在依撒意亚先知书中我们读到：\BibleQuote{我是最初的天主，我也是最后的天主，除我以外没有别的天主。在我以前没有天主，在我以后也不会有天主，除我以外没有别的天主。}\EditorialNote{依撒意亚 43:10} 主也在神圣的福音中对祂的父说：\BibleQuote{永生就是：认识你，唯一的真天主。}\BibleRef{若 17:3} 但至于那些不信圣经的人，我们将以理性论证如下。\par

神性是完美的，在良善、智慧、能力上无有瑕疵，无始、无终、永恒、无限，简言之，在一切方面都是完美的。倘若我们说有很多神，那么就必须承认众多之间有差异。因为如果他们之间没有差异，那么他们是一而非多。但如果他们之间有差异，那么完美性何在？凡在良善、能力、智慧、时间或空间方面有所欠缺的，不能是天主。正是这各方面的一致表明神性是独一而非多。\par

再者，若有许多神，如何能维持天主是无限的？因为一处有了一个，另一处就不能有了。\par

此外，世界如何能被许多神治理，并在统治者之间争战可见时，免于瓦解和毁灭？因为差异引入争战。若有人说每个神统治一部分，那么是谁建立了这一秩序并赐予每个神其特定的领域？因为这或许更应是神。因此，天主是独一的、完美的、无限的，宇宙的创造者、保存者和治理者，超越并先于一切完美。\par

再者，由统一产生二元是自然之必然。\par

% structure: p0026 source=h2 target=subsection reason=relative-heading-rank; base=h2

\subsection{第六章　论圣言与天主子：一个论证性的证明。}

因此，这独一的天主并非无言（无圣言）。祂既拥有圣言，则这圣言并非无自立存在、非有起始、亦非将归于无有。因为从未有过天主无圣言之时：祂永远拥有从祂自己所生的圣言，不像我们的言语那样无自立存在而消散于空中，而是在祂内有自立存在、生命和完美，不离开祂自身而恒常居于祂内。因为若它要出去到祂之外，又能到哪里去？因我们的本性是可朽坏、易消散的，我们的言语也是无自立存在的。但既然天主是永恒而完美的，祂的圣言在祂内必然是有自立存在的、永恒的、生活的，且拥有生者的一切属性。正如我们的言语从心灵发出，既不完全与心灵相同，也不全然不同（就它从心灵发出而言，它不同于心灵；就它显示心灵而言，它不再绝对不同于心灵，而是与心灵同性，但对主体而言又不同于它），同样，天主的圣言在其独立存在中，也与那赋予它存在者有别；但就它在自身中显示与天主相同的属性而言，它与天主同性。正如在父内彻见绝对的完美，同样在从祂所生的圣言内也彻见完美。\par

% structure: p0028 source=h2 target=subsection reason=relative-heading-rank; base=h2

\subsection{第七章　论圣神：一个论证性的证明。}

再者，圣言亦必有圣神。因为就连我们自己的言语也并非没有气；但在我们身上，气是不同于我们本质的某物。因为有一种空气的吸力和运动，被吸入并呼出以使身体得以维持。正是这气在发声时刻成为清晰的言语，在自身中揭示言语的力量。但在神性本质中，它是单纯而不可复合的，我们必须以全部虔诚承认存在着天主的圣神，因为圣言并不比我们自己的言语更不完美。现在，我们不能虔诚地认为圣神是某种从外部进入天主的异物，如同我们这些复合本质的情形那样。相反，正如当我们听到天主的圣言时，我们将其视为并非无位格，也不是学习的产物，也不是单纯的声音表达，也不是消逝于空中而灭亡，而是本质性存在的，具有自由意志、行动和全能：同样，当我们得知天主的圣神时，我们将其视为圣言的同伴和祂行动的揭示者，而非无位格的纯粹气息。因为将居住于天主内的圣神设想为像我们自己的精神一样，就是将神性本质的伟大拖入最低的堕落深渊。但我们必须将其视为一种本质性的力量，存在于其自身固有的位格中，由圣父所发，安息于圣言内，并彰显圣言，既不能与祂所在的天主以及祂所伴随的圣言分离，也不会倾倒而消失于虚无，而是具有类似圣言的位格，赋有生命、自由意志、独立行动、能力，始终意愿那善的，并有能力在一切法令中与意志同步，无始无终。因为圣父从未在任何时候缺少圣言，圣言也从未缺少圣神。\par

因此，由于本质上的合一，希腊人认为天主是多的错误就被彻底摧毁了；而又因我们接受圣言和圣神，犹太人的教义就被推翻了；双方只留下有益的部分。一方面，犹太人的观念中我们有了天主本质的合一；另一方面，希腊人的观念中我们有了位格上的区别，而且仅此而已。\par

但倘若犹太人拒绝接受圣言和圣神，愿圣经驳斥他，勒住他的舌头。因为关于圣言，神圣的达味说：\BibleQuote{上主，祢的圣言永远安定在天}。又说：\BibleQuote{祂派遣自己的圣言，治好了他们}。关于圣神，同一位达味说：\BibleQuote{祢遣发祢的圣神，它们便被造成}。又说：\BibleQuote{因上主的话，诸天造成；因祂口中的气，万象生成}。约伯也说：\BibleQuote{天主的圣神造了我，全能者的气息赐给我生命}\BibleRef{约 33:4}。如今，那被派遣、创造、建立并保存的圣神，并非消散的气息，同样，天主的口也不是身体的器官。因为对两者的理解必须符合神性本质。\par

% structure: p0032 source=h2 target=subsection reason=relative-heading-rank; base=h2

\subsection{第八章 论天主圣三}

因此，我们相信独一的天主，一个元始，无始，非受造，非受生，不朽坏且不死，永恒，无限，无界限，无边际，无限能力，单纯，不可复合，无形体，无流变，无情感，不变，不改，不可见，是善与义的源泉，心智之光，不可接近；一种能力，不被任何尺度所认识，仅凭祂自己的意志即可度量（因为凡祂愿意的，祂都能），是一切可见或不可见受造物的创造者，是一切万物的维护者与保存者，是万物的供养者，是全能的统治者和君主和王，拥有无尽不朽的国度；没有敌对者，充满万有，不受任何包围，反而是万有的包围者、维护者和原初拥有者，占据一切本质而无损，超越万物，与一切本质分离，因为祂是超本质、超越一切，是绝对的天主、绝对的善、绝对的圆满；规定一切主权和等级，置于一切主权和等级之上，超越本质、生命、言语和思想；祂本身就是光、善、生命和本质，因为祂不是从另一来源获得自己的存在，即就存在的事物而言；但祂本身就是一切存在者的存在之源，生活者的生命之源，有理性者的理性之源；对万有一切善的起因；在一切尚未成为之前就已认识一切；一个本质，一个神性，一个能力，一个意志，一个行动，一个元始，一个权威，一个统治，一个主权，在三个完美的位格中被认知，并受一个崇拜的敬拜，被一切理性受造物所信仰和事奉，无混淆地合一，无分离地分开（这确实超越了思想）。（我们相信）圣父、圣子和圣神，我们也受洗归于祂们。因为如此，我们的主命令宗徒们施洗，说：\BibleQuote{把他们因父及子及圣神之名给他们授洗}\BibleRef{玛 18:19}。\par

（我们信）在一位圣父，万有的根源和原因：非由谁生：无原因或无\OriginalTerm{生育}{generation}，独自存在：万有的创造者：但\OriginalTerm{本性}{nature}上只是一位圣父，祂的\OriginalTerm{独生}{Only-begotten}子、我们的主、天主和救主耶稣基督，并发出至圣圣神。又在一位于天主之子，\OriginalTerm{独生}{Only-begotten}者，我们的主，耶稣基督：由圣父所生，在万世之前：光之光，真天主之真天主：受生而非受造，与圣父\OriginalTerm{同体}{consubstantial}，万物藉祂而受造：当我们说祂在万世之前，就表明祂的诞生是超时间、无开始的：因为天主之子并非从无中被造，祂是光荣的辉耀，圣父\OriginalTerm{位格}{subsistence}的印记，生活的智慧与能力\BibleRef{格前 1:24}，拥有内在\OriginalTerm{位格}{subsistence}的圣言，不可见天主的\OriginalTerm{本质}{essence}而完美而生活的肖像\BibleRef{希 1:3}。但祂始终与圣父同在、在圣父内，永远而无始地由祂所生。因为从来没有一个时候圣父存在而圣子不存在，而是圣父和由祂所生的圣子始终一同存在。因为圣父离开圣子就不能得到圣父的名号：因为祂若没有圣子，就不能是圣父：如果祂后来才有圣子，后来才成为圣父，先前并不是圣父，祂就由不是圣父改变成圣父。这是最恶劣的亵渎。因为我们不能说天主缺乏自然的\OriginalTerm{生育}{generation}能力：\OriginalTerm{生育}{generation}能力意味着，从自身、即从自己本有的\OriginalTerm{本质}{essence}产生与自己\OriginalTerm{本性}{nature}相似者。\par

在论及圣子的\OriginalTerm{生育}{generation}时，说时间介入、并且圣子的存在晚于圣父，这是不虔敬的行为。因为我们主张，圣子是出自祂，即出自圣父的\OriginalTerm{本性}{nature}而被\OriginalTerm{生育}{generation}。除非我们承认圣子从起初就与圣父同在，由祂所生，否则我们就引入圣父\OriginalTerm{位格}{subsistence}的变化，因为祂先前不是圣父，后来才成为圣父。因为受造之物，即便后来才出现，也并非从天主的\OriginalTerm{本质}{essence}而来，而是藉祂的旨意和能力从无中受造，变化不触及天主的\OriginalTerm{本性}{nature}。因为\OriginalTerm{生育}{generation}意味着生育者从自己的\OriginalTerm{本质}{essence}产出与\OriginalTerm{本质}{essence}相似的\OriginalTerm{本性}{nature}的后代。但\OriginalTerm{创造}{creation}和制造意味着创造者和制造者从外在之物、而非从自己\OriginalTerm{本质}{essence}产出具有绝对异质\OriginalTerm{本性}{nature}的\OriginalTerm{创造}{creation}物。\par

因此，在天主内，只有祂是无情感、不可改变、永不变易，始终如此，无论是\OriginalTerm{生育}{generation}还是\OriginalTerm{创造}{creation}都是无情感的。因为祂\OriginalTerm{本性}{nature}上无情感、无流变，因为祂是单纯、非复合的，所以不论在\OriginalTerm{生育}{generation}或\OriginalTerm{创造}{creation}中，都不受情感或流变的影响，也不需任何合作。但祂内的\OriginalTerm{生育}{generation}是无始而永恒的，是\OriginalTerm{本性}{nature}的工作，并出自自己的\OriginalTerm{本质}{essence}而产生，使生育者不经历变化，使祂不是起初是天主后来是天主，也不接受任何增益；而在天主造物的情形中，\OriginalTerm{创造}{creation}是意志的工作，不与天主同永恒。因为从无中受造之物，与无始而永恒者同永恒，是不自然的。实际上，人的制造与天主的制造有以下区别。人不能从无中产生任何东西，而他所制造的一切都需要预先存在的物质作为基础，并且他不只靠意志创造，而是先思考要成为什么，并在心中描绘，然后才用手塑造，经历劳苦和困难，并且常常失的，无法满意地产生他所追求的结果。但天主，仅凭意志的行使，就使万物从无中受造。现在，在\OriginalTerm{生育}{generation}和滋生方面，天主与人也有同样的区别。因为在无时间、无始、无情感、无流变、无形体、独一、无终的天主内，\OriginalTerm{生育}{generation}是无时间、无始、无情感、无流变、也不依赖两者的结合：祂自己不可理解的\OriginalTerm{生育}{generation}也没有开始或结束。它无始是因为祂是永不变易的；无流变是因为祂无情感和无形体；又不依赖两者的结合，是因为祂是无形体，也因为祂是独一唯一的天主，不需要任何合作；无终或无尽是因为祂无始、无时间或无终，并永恒保持同一。因为无始者无终；但藉\OriginalTerm{恩宠}{grace}而无尽者，确实不是无始的，例如天使。\par

因此，永恒天主产生祂自己的圣言，完美、无始无终，使那位\OriginalTerm{本性}{nature}及存在超时间的天主，不致在时间中\OriginalTerm{生育}{generation}。但人显然不同，因为人的\OriginalTerm{生育}{generation}涉及性别、毁灭、流变、增长，并且身体包裹着他，他拥有男性或女性的\OriginalTerm{本性}{nature}。因为男性需要女性的协助。但愿那超越万有、超越一切思想和理解的天主，垂怜我们。\par

圣而公教会、宗徒的教会教导：圣父与祂的独生子同时存在，圣子由圣父所生，无时间、无变化、无受苦，其方式不可理解，唯有宇宙的天主能领悟；正如我们承认火与从火发出的光同时存在：并非先有火而后有光，而是二者并存。又正如光常由火产生，常在于火，绝不与火分离，同样，圣子由圣父所生，绝不与祂分离，而常在于祂。然而，从火无分离地产生、常在于火的光，本身并无区别于火的本体（因为它乃是火的自然属性），但天主的独生子，由圣父无分离、无差别而生，常在于祂，却拥有区别于圣父的本有本体。\par

然后，\Quote{圣言}和\Quote{光辉}这些术语被使用，因为祂由圣父所生，没有两个的结合、没有受苦、没有时间、没有变化、没有分离；而\Quote{子}和\Quote{圣父本体的印}这些术语，因为祂是完美的，拥有本体，并且在各方面相似于圣父，除了圣父不是受生的；以及\Quote{独生子}这个术语，因为唯独祂是单独由单独的圣父所生。因为没有别的受生相似于天主子的受生，因为别的不曾是天子。因为虽然圣神由圣父所发，但这不是生发性的，而是发出性的。这是一种不同的存在方式，同样不可理解和未知，正如子的受生一样。因此，圣父所拥有的一切属性，圣子也有，除了圣父是未受生的，这一例外不涉及本质或尊严的差异，而只涉及存在方式的不同。我们有一个类比：\Person{亚当}不是受生的（因为天主亲自塑造了他），\Person{舍特}是受生的（因为他是亚当的儿子），\Person{厄娃}是从亚当的肋骨发出的（因为她不是受生的）。这些在本性上彼此没有区别，因为他们都是人；但他们在存在方式上不同。\par

因为必须认识到，只有一个\OriginalTerm{\GreekText{ν}}{\GreekText{ν}}的\OriginalTerm{非受造的}{\GreekText{ἀγένητον}}一词意指\Quote{非受造的}或\Quote{未被造的}，而写有双\OriginalTerm{\GreekText{ν}}{\GreekText{ν}}的\OriginalTerm{非受生的}{\GreekText{ἀγέννητον}}意指\Quote{非受生的}。按照第一种含义，本质与本质不同：因为一种本质是非受造的，即只有一个\OriginalTerm{\GreekText{ν}}{\GreekText{ν}}的\OriginalTerm{非受造的}{\GreekText{ἀγένητον}}，另一种是受造的，即\OriginalTerm{受造的}{\GreekText{γενητή}}。但在第二种含义中，本质与本质没有区别。因为所有生物的第一本体是\OriginalTerm{非受生的}{\GreekText{ἀγέννητος}}，但不是\OriginalTerm{非受造的}{\GreekText{ἀγένητος}}。因为它们是由造物主创造，藉祂的圣言而存在，但它们不是受生的，因为没有预先存在的像它们自身的形态，可从它们产生。\par

因此，在第一种含义中，圣三的三个绝对神圣的位格是一致的：因为它们同属一个本质且是非受造的。但在第二种含义中，则完全不同。因为只有圣父是\OriginalTerm{非受生的}{ingenerate}，没有别的位格赋予祂存在。圣子单独是\OriginalTerm{受生的}{generate}，因为祂由圣父的本质所生，无始无时。而只有圣神由圣父的本质\OriginalTerm{发出}{proceeds}，不是受生的，而只是发出。\BibleRef{若 15:26} 因为这是圣经的教导。但受生和发出的本性完全超出理解。\par

我们也应当知道，父职、子职、发出这些名号不是我们加给圣三的；相反，是由圣三传给我们，正如神圣的宗徒说：\BibleQuote{为此，我屈膝在圣父面前，天上地下的一切家族都由祂得名}。但如果我们说圣父是圣子的本源，并大于圣子，我们并非暗示圣父在时间上优先于圣子，或在本性上凌驾于圣子 \BibleRef{若 14:28}（因为圣父藉着圣子创造了世代），或任何其他方面的优越性，除了因果关系。我们以此意指圣子是由圣父所生，而非圣父由圣子所生，并且圣父按本性是圣子的原因：正如我们同样不说火从光发出，而是光从火发出。因此，每当我们听到说圣父是圣子的本源并大于圣子，让我们理解为就因果关系而言。正如我们不说火是一种本质而光是另一种，所以我们也不能说圣父是一种本质而圣子是另一种：二者是同一本质。正如我们说火因从它发出的光而明亮，并不认为火的光是服侍火的工具，而是它的自然力量：同样，我们说圣父藉祂的独生子创造一切，并非好像圣子只是服务于圣父旨意的工具，而是作为祂的自然而本体性的力量。正如我们说火发光，又可以说火的光发光，\BibleQuote{凡圣父所做的，圣子也同样做}。\BibleRef{若 5:19} 但是光没有区别于火的本有本体，而圣子是一个圆满的本体，与圣父的本体不可分离，如我们上文所显示。因为在受造物中绝不可能找到一个形象本身能完全精确地详细阐明圣三的本性。因为那受造的、复合的、可变的、可变化的、有限制的、有形式的、可朽坏的，怎能清楚地显示那超本质的神圣本质，而后者不受这些任何方面的影响呢？现在显然，所有受造物都容易受这些影响中的大部分，并且所有受造物由于自身的本性都趋向朽坏。\par

同样，我们也信唯一圣神，主及赋予生命者：由圣父所发，寓于圣子；与圣父圣子同受钦崇与荣耀，因祂与圣父圣子同体、同永恒；天主之神，直接、权威，智慧、生命与圣德的源泉；与圣父圣子一同存在并受呼求的天主；非受造、圆满、创造、统御万有、施行万有、全能、无限能力、一切受造之主，而非受制于任何主；使人成神，而非被成神；充满，而非被充满；被分享，而非分享；圣化，而非被圣化；代祷者，收纳众人的祈求；在一切事上与圣父圣子相似；由圣父所发，通过圣子传递，为一切受造物所分享，藉自身创造、赋与本质、圣化并维系宇宙；拥有自立体，存在于自己专有而独特的自立体中，与圣父圣子不可分不可离，并拥有圣父圣子所有的一切属性，唯独不具备非生或非受生。因为圣父是无因、非受生的：祂不从任何事物而来，而是从自身获得存在，也不从其他事物获得任何属性。相反，祂本身就是一切事物以确定而自然的方式存在的起始与原因。圣子则是从圣父以生的方式而来，同样圣神也是从圣父而来，但不是以生的方式，而是以发的方式。我们得知生与发之间有区别，但区别的本质我们全然不知。此外，圣子由圣父所生与圣神的发是同时的。\par

那么，圣子和圣神所有的一切皆来自圣父，甚至他们的存在本身也是如此：若非圣父存在，圣子与圣神亦不存在。若非圣父拥有某一属性，圣子与圣神亦不拥有该属性；并且藉着圣父，即因圣父的存在，圣子与圣神存在；藉着圣父，即因圣父拥有那些属性，圣子与圣神拥有他们所有的一切属性，惟独非受生、生与发的属性除外。因为只有在这\OriginalTerm{位格的}{hypostatic}或\OriginalTerm{位格的}{personal}属性上，三个圣自立体彼此不同，它们虽不可分地分开，却不是按本质，而是按它们自己专有而独特的自立体之区别标记。\par

此外，我们说三者各有完美的自立体，以便我们理解的不是一个由三个不完美元素组成的复合完美本性，而是一个单一、单纯、超越且先于完美的本质，存在于三个完美的自立体中。因为凡是由不完美元素组成的东西必然复合。但从完美的自立体中不可能产生复合物。因此，我们不谈论形式来自自立体，而是形式存在于自立体中。我们称那些不保持由其完成之物的形式的事物为不完美。例如，石头、木头和铁各自在其本性上完美，但相对于由它们完成的建筑，它们各自不完美：因为其中任何一个本身都不是房屋。\par

因此，我们说这些自立体是完美的，以免我们设想神性本质为复合。因为复合是分离的开端。同样，我们谈论三个自立体彼此寓居，以免引入一群或许多的神。藉着三个自立体，没有复合或混淆；而藉着它们拥有同一本质、彼此寓居、在意志、能力、力量、权柄和运动上相同（可以这么说），我们认识到神的不可分性和合一性。因为确实只有一位天主，及祂的圣言和圣神。\par

\EditorialNote{手稿边缘注释：论三个自立体的区别；以及论事物本身、我们的理性与思想与之的关系。}\par

此外，应当认识到，观察事物本身是一回事，以理性与思想观察是另一回事。在一切受造物中，自立体的区别是在现实中被观察到的。因为实际上，我们看到\Person{伯多禄}与\Person{保禄}是分开的。但共同性、联系与统一则是通过理性与思想被把握的。因为我们凭心智感知\Person{伯多禄}与\Person{保禄}具有同样的本性，共享一个共同本性。因为两者都是活物，有理性且可朽；两者都是肉躯，赋有理性和理解的精神。因此，这种本性的共同性是通过理性观察到的。因为在这里，自立体并不彼此寓居。但每个自立体私下地、个别地，即在其自身中，完全分离，有许多点将它与其他自立体分开。因为它们在空间上分离，在时间上不同，在思想、能力、形状或形式、习惯、气质、尊严、追求以及所有区分属性上被划分，尤其是它们并非彼此寓居，而是分离。由此，我们可以谈论两个、三个或多个人。\par

这种情形可以在整个受造界中看到，但是在神圣、超本质且不可理解的天主圣三身上（祂远远超越万物），情况恰好相反。因为在那里，透过自立的共永，以及它们拥有同一本质、同一行动、同一意志、同一心意的一致，并且在权威、能力和良善上相同——我不是说相似，而是相同——并且在同一推动下运动，我们实际上观察到共通与合一。因为只有一个本质、一个良善、一个能力、一个意志、一个行动、一个权威，同一个（我再重复），不是三个彼此相似。但三个自立具有同一个运动。因为其中的每一个都与另一个紧密相连，如同与自己相连一样：也就是说，圣父、圣子、圣神在各方面都是一体的，除了非受生、受生和发出。但差异是藉着思想而被认知的。因为我们承认一位天主：但只有在父性、子性和发出（即原因与结果，以及自立的圆满，即存在方式）的属性上，我们才察觉到差异。因为关于无限的天主性，我们不能像在我们自身中所做的那样谈论空间上的分离。因为自立的彼此寓居，毫不混淆，而是紧密相连，按照主的话：\BibleQuote{我在父内，父在我内}\BibleRef{若 14:11}；也不能承认在意志、判断、行动、能力或任何其他在我们中产生实际绝对分离的事情上有什么差异。因此，我们不谈论三位神——圣父、圣子、圣神——而只谈论一位神，即圣三，圣子和圣神被归因于一个原因，并且不是按照撒伯里乌的融合而复合或合并。因为，如我们所说，它们被合而为一，不是以便混合，而是以便彼此依附，并且它们彼此寓居，没有任何合并或混合。圣子和圣神也不分开，也不按照亚流的分裂在本质上被切割。因为天主性在分开的事物中是未分开的（简而言之）：就像三个太阳彼此依附而不分离，发出混合而联合的光。因此，当我们注视天主性、第一原因、主权、一体性和（可以这么说）运动与意志的同一性，以及本质、能力、行动和主宰的同一性时，我们看到的是合一。但是当我们观察那些天主性所在的事物（或更准确地说，就是天主性本身的事物），以及那些通过第一原因而无时间、无荣耀区别或无分离地存在于它之内的事物，即圣子和圣神的自立时，我们所敬拜的似乎是天主圣三。圣父是一位父，且是无始的，即没有原因：因为祂不来自任何事物。圣子是一位子，但不是无始的，即不是没有原因的：因为祂来自圣父。但是，如果你从时间中消除开始的概念，祂也是无始的：因为时间的创造者不能受时间支配。圣神是一位神，从圣父发出，不是以子性的方式，而是以发出的方式；这样，圣父没有因为祂生了而失去祂无始的属性，圣子也没有因为祂是由无始者所生而失去祂受生的属性（因为怎么可能呢？），圣神也没有因为祂发出了并且是神而变成圣父或圣子。因为一个属性是相当恒常的。因为如果一个属性是可变的、可移动的并能变成别的东西，它怎么能持久呢？因为如果圣父是圣子，祂严格来说就不是圣父：因为严格来说只有一位圣父。如果圣子是圣父，祂严格来说就不是圣子：因为严格来说只有一位圣子和一位圣神。\par

此外，应该理解，我们不称圣父来自任何人，而称祂为圣子之父。我们不称圣子为原因或父，而称祂既来自圣父，又是圣父之子。同样，我们称圣神来自圣父，并称祂为圣父之神。我们不称圣神来自圣子：但我们仍称祂为圣子之神。\BibleQuote{因为凡没有基督之神的，就不属于祂}\BibleRef{罗 8:9}，神圣的宗徒说。我们宣认祂是通过圣子显示并赐予我们的。\BibleQuote{祂向门徒嘘了一口气，说：你们领受圣神吧}\BibleRef{若 20:29}。这就像太阳的情况，从太阳发出光线和光辉（因为太阳本身是光线和光辉的源头），通过光线，光辉传到我们，而光辉本身使我们得光照，我们有分于它。此外，我们不称圣子来自圣神，也不称圣子出于圣神。\par

% structure: p0051 source=h2 target=subsection reason=relative-heading-rank; base=h2

\subsection{第九章。论关于天主所肯定的事。}

天主性单纯而不可复合。但由许多不同元素构成的事物是复合的。那么，如果我们谈论非受造、无始、无形体、不朽、永恒、善、创造性的等等特质作为天主本质上的差异，那么由如此多特质构成的事物就不会是单纯的，而必须是复合的。但这是极不虔敬的。因此，关于天主的每一肯定都应该被认为不是表示祂本质是什么，而是表示某种无法说明的东西，或是与某些相对之物或某些本性之果的关系，或是一个行动。\par

由此可见，在天主所有名称中，最恰当的是\Quote{\OriginalTerm{自有者}{He that is},}正如天主自己在山上答覆梅瑟时所说：\BibleQuote{你要对以色列子民说：\InnerQuote{自有者}派遣了我。}\BibleRef{出 3:14}因为天主怀抱一切存有，犹如一片无限而不可见的本质海洋。或者，正如\Person{圣狄奥尼修斯}所说：\Quote{\OriginalTerm{善者}{He that is good}.}因为不能说天主首先是存有，其次才是善。\par

天主第二个名称是\OriginalTerm{天主}{\GreekText{ὁ} \GreekText{Θεός}}，源自\OriginalTerm{奔跑}{\GreekText{θέειν}}，因为天主运行于万物之中；或源自\OriginalTerm{燃烧}{\GreekText{αἴθειν}}，因为\BibleQuote{天主是吞灭一切邪恶的烈火}\BibleRef{申 4:24}；或源自\OriginalTerm{观看}{\GreekText{θεᾶσθαι}}，因为天主无所不见\BibleRef{加下 10:5}。因为没有任何事物能逃避祂，祂洞察一切。祂在万物存在之前就已看见它们，将它们在永恒的思想中保存；每一事物依照祂自由而永恒的思想——这思想构成预定、肖象与模型——在预定的时候进入存在。\par

因此，第一个名称表达祂的存在及其存在的本性；第二个名称则包含行动的概念。此外，\Quote{无始}、\Quote{不朽}、\Quote{无生}，以及\Quote{非受造}、\Quote{无形体}、\Quote{不可见}等术语，说明祂不是什么：也就是说，它们告诉我们，祂的存在没有开端，祂不朽、非受造、无形体、不可见。再者，善、正义、虔诚等这类名称属于本性，但并不解释祂的实际本质。最后，主、王以及这类名称表示与对立面的关系：因为主的名称与主所统治的人相关，王的名称与受君王权威管辖的人相关，造物主的名称与受造物相关，牧者的名称与所牧放的羊相关。\par

% structure: p0056 source=h2 target=subsection reason=relative-heading-rank; base=h2

\subsection{第十章 论天主的合一与分离。}

因此，所有这些名称都应理解为整个神性的共通称谓，且包含同一、单纯、不可分与合一的概念；而父、子、圣神，以及无因、有因，无生、受生，与发源等名称则包含分离的概念：因为这些术语并不解释祂的本质，而是说明相互的关系与存在的方式。\par

当我们领悟了这些并从这些上升到神圣本质时，我们并非把握本质本身，而只是本质的属性：正如即使我们知道灵魂是无形体、无大小、无形状的，我们也并未把握灵魂的本质；同样，当我们知道身体是白的或黑的时，也并未把握身体的本质，而只是本质的属性。再者，真确的教义教导：神性是单纯的，具有一种单纯、善而运行于万物的行动，正如太阳的光线，温暖一切事物，并按照每样事物的自然禀赋与接受能力而在其中运行，这光线从它的创造者天主那里获得了这种行动的形式。\par

但关于圣言神圣而仁慈的降生成人的一切，则截然不同。因为在这事上，圣父与圣神都毫无参与，除非就赞同以及天主圣言——祂如同我们一样成为人，作为不变的天主与天主子——所行的不可思议的奇迹而言。\par

% structure: p0060 source=h2 target=subsection reason=relative-heading-rank; base=h2

\subsection{第十一章 论及好像天主有身体那样的肯定说法。}

由于我们在圣经中发现许多关于天主的用语是象征性地使用的，且更适用于有形体的东西，我们应当承认，我们这些被厚厚的血肉之躯所包裹的人，除非借由取自我们生活的形象、类型和象征，否则绝不可能理解或谈论天主性无形体的、崇高而超性的运行。因此，所有关于天主的有形体含义的陈述都是象征，但具有更高的意义：因为天主是单纯而无形式的。因此，由天主的眼睛、眼睑和视线，我们应理解祂监察万物的能力和知识，什么也逃不过；因为在我们身上，这个感官使我们的知识更完整、更确切。由天主的耳朵和听觉，意指祂乐意被安抚并接受我们的祈求；因为这个感官也使我们慈祥地对待恳求者，更和蔼地垂听他们。天主的嘴和言语是祂表明自己旨意的方式；因为我们正是通过嘴和言语表达内心的思想：天主的食物和饮料是我们对祂旨意的同意，因为我们也通过味觉满足自然欲望的需求。天主的嗅觉是祂对我们思想和祂善意的欣赏，因为我们通过这个感官欣赏甜美的香气。天主的面容是通过祂的作为显示和彰显祂自己，因为我们的彰显是通过面容。天主的双手意味着祂能力的有效本性，因为我们用自己的双手完成最有用和最有价值的工作。祂的右手是祂在顺境中的援助，因为我们也是用右手制造形状优美或价值极高的东西，或需要大力气的地方。祂的触摸是祂精确辨别和追究的能力，即使是最微小最隐秘的细节，因为我们所触摸过的人不能向我们隐瞒他们内在的任何东西。祂的脚和行走是祂的来临和临在，或是为了援助需要者，或是为了报复仇敌，或是为了执行任何其他行动，因为我们是用脚到达任何地方的。祂的起誓是祂旨意的不变性，因为我们是用誓言确认彼此之间的契约。祂的愤怒和暴怒是祂对一切邪恶的憎恨和厌恶，因为我们也憎恨违背我们心意的事物，并对此发怒。祂的遗忘、睡眠和打盹是对仇敌报复的延迟，以及对祂子民惯常帮助的推迟。简言之，所有关于天主的有形体含义的陈述都有某种隐藏的意义，通过我们熟悉的事物教导我们超出我们理解的事物，但有关天主圣言肉身居留的陈述除外。因为祂为了我们的得救，亲自取了完整的人性——理性的神魂、身体以及人性的一切属性，甚至包括自然而无罪的情感。\par

% structure: p0062 source=h2 target=subsection reason=relative-heading-rank; base=h2

\subsection{第十二章 论同一主题。}

以下就是我们由神圣的训示所学得的奥秘，正如神圣的亚略巴古人\Person{狄约尼削}所说的：\Quote{天主是万有的原因和起始：一切有本质者的本质：一切有生命者的生命：一切理性者的理性：一切有理智者的理智：远离祂者的召回和恢复：败坏本性者之更新和变化：在邪恶中飘摇者之神圣根基：坚定站立者之稳固：那些朝向祂前行者的道路，以及引领他们向上的伸出的手。}我还要补充说，祂也是祂所有受造物的父亲（因为从无中把我们带入存在的天主，在更严格的意义上是我们的父亲，胜于我们的父母，因为父母也从祂那里获得了存在和生育），是跟随祂并受祂牧养的人之牧人，是被光照者的光辉，是被启迪者的启迪，是被神化者的神化，是不和者的和平，是喜爱纯朴者的纯朴，是崇拜合一者的合一，是万始之始，\OriginalTerm{超本质}{super-essential}，因为超越一切起始，以及隐藏之事的良善启示，即对祂的认识，各人按其所合法和所能达到的限度。\par

\Emph{进一步更精确地论天主的名号}\par

天主既然是无可理解的，也必定是无名的。因此，既然我们不知道祂的本质，就不要为祂的本质寻求名号。因为名号是对实际事物的说明。但是，天主是良善的，祂从无中把我们带入存在，使我们能分享祂的良善，并赐予我们认识的能力，不仅没有将祂的本质启示给我们，甚至连祂本质的认识也未曾赐予。因为本性不可能完全理解超性。此外，如果知识是对存在之物的认识，那么对超本质之物又如何能有知识呢？因此，由于祂不可言说的良善，祂乐意被我们能够理解的名号称呼，使我们不致完全与对祂的认识隔绝，而是对祂有一些概念，尽管模糊。因此，既然祂是无可理解的，祂也是不可名状的。但既然祂是万有的原因，并在自身内包含万有的理性和原因，祂就从万有中获得名号，甚至从对立的事物：例如，祂被称为光明和黑暗、水和火，为使我们知道这些不是祂的本质，而是祂是\OriginalTerm{超本质}{super-essential}和不可名状的；但既然祂是万有的原因，祂从祂的一切效果中获得名号。\par

因此，在天主的名称中，有些具有否定的意义，表明祂是超本质的：例如\Quote{非本质的}、\Quote{无时间的}、\Quote{无始的}、\Quote{不可见的}——这并不是说天主逊于什么或缺乏什么（因为万物都属于祂，并从祂而来，借着祂而存在，并在祂内得以存留\BibleRef{哥 1:17}），而是说祂超越一切，与一切分离。因为祂不是存在物中的一种，而是超越万有。另一些名称则具有肯定的意义，表明祂是万有的原因。作为一切存在和一切本质的原因，祂被称为\OriginalTerm{存在}{Ens}和\OriginalTerm{本质}{Essence}。作为一切理性、智慧、有理性和有智慧者的原因，祂被称为理性和有理性的、智慧和有智慧的。同样，祂被称为理智和有理性的、生命和有生命的、能力和有能力的，以及其他类似的名称。更准确地说，那些源自最宝贵、与祂最相似之物的名称对祂最为恰当。非物质的比物质的更宝贵、更相似于祂，纯洁的比不纯洁的、神圣的比不神圣的，因为它们更多地分有祂。因此，太阳和光比黑暗、白天比黑夜、生命比死亡、火、灵和水（作为有生命者）比大地、尤其是良善比邪恶，更能作为祂的名称——这也就是说，存在比非存在更合适。因为良善是存在，并且是存在的原因，而邪恶则是良善的否定，即存在的否定。这些就是肯定和否定的名称，但最甜美的名称是两者的结合：例如，超本质的本质、超越天主的\OriginalTerm{神性}{Godhead}、超越开始的开始，等等。此外，还有些关于天主的肯定说法，在卓越的意义上具有否定的力量：例如，黑暗——这并不是说天主是黑暗，而是说祂不是光，而是超越光。\par

天主被称为理智、理性、灵、智慧和能力，因为祂是这些的原因，并且是非物质的、万有的创造者、全能者。这些名称，无论是肯定的还是否定的，都适用于整个天主性。它们也以完全相同的方式和完全的意义用于天主圣三的每一位位格。因为当我思考其中一个位格时，我认出祂是完全的天主和完全的本质；但当我把三者结合并加在一起时，我知道一个完全的天主。因为天主性不是复合的，而是在三个完美的位格中，是一个完全不可分割、非复合的天主。当我思考三个位格彼此之间的关系时，我领悟到圣父是超本质的太阳、良善的泉源、本质的无底深渊、理性、智慧、能力、光明、神性：是隐藏于自身内的良善的产生和生殖的源头。祂本身就是理智、理性的深渊、圣言的生出者，并通过圣言发出启示的圣神。简言之，圣父除了圣子——圣父的唯一能力、宇宙创造的直接原因——没有其他理性、智慧、能力、意志；圣子是完美的位格，由完美的位格以一种祂自己所知道的方式所生，祂是并被称为圣子。圣神是圣父的能力，启示祂神性的隐秘奥秘，由圣父藉圣子以祂自己所知道、但不同于生发的方式而发出。因此，圣神是宇宙创造的完成者。因此，所有适合圣父的术语，如原因、泉源、生发者，都应归于圣父；适合受生的圣子、圣言、直接能力、意志、智慧的术语，应归于圣子；适合受发的、启示的、完成的能力的术语，应归于圣神。圣父是圣子和圣神的泉源和原因：圣父单独是圣子的父，并是圣神的发出者。圣子是子、圣言、智慧、能力、肖像、光辉、圣父的印迹，并由圣父所出。但圣神不是圣父的子，而是圣父的圣神，因为由圣父发出。因为没有圣神就没有动力。我们也称圣神为圣子的圣神，不是由于祂从圣子发出，而是由于祂藉圣子从圣父发出。因为只有圣父是原因。\par

% structure: p0068 source=h2 target=subsection reason=relative-heading-rank; base=h2

\subsection{第十三章　论天主的位置：唯有神性是不受限制的。}

身体的位置是包容者的界限，被包容者由此被包容：例如，空气包容，而身体被包容。但并不是整个包容的空气都是被包容身体的位置，而是包容空气的界限，即它与被包容身体接触的地方；原因很清楚，因为包容者不在它所包容的事物之内。\par

但也有心智空间，理智在此活动，而心智及无形体的本性存在：那里是理智居住、运作并被容纳之处，不是以有形体的方式，而是以心智的方式。因为它是无形状的，因此不能像形体那样被容纳。那么，天主既非物质且\OriginalTerm{无界限的}{uncircumscribed}，便没有地方。因为祂是自己的地方，充满万有，超越万有，并亲自保全万有。但我们仍说天主有地方，天主的场所就是祂的\OriginalTerm{能力}{energy}彰显之处。因为祂渗透万物而不与万物混合，并按各物的相称与接受能力，将祂的能力分施给所有：我是指一种既属自然又属自愿的纯洁。因为无形之物比有形之物更纯洁，有德之物比与恶相连之物更纯洁。因此，所谓天主的场所，就是指在祂的能力和恩宠中更有份的地方。为此，天堂是祂的宝座。因为在其中，天使们奉行祂的旨意，时常光荣祂。这是祂的安息之所，而大地是祂的脚凳。\EditorialNote{依撒意亚 66:1}因为祂在肉身中居住在人们中间。\BibleRef{巴 3:38}而且祂的圣肉身被称为天主的脚。教会也被称为天主的场所：因为我们把它分别为圣，为光荣天主，作为一种祝圣的地方，我们在其中也与祂交谈。同样，那些祂的能力借着肉身或脱离肉身而向我们彰显的地方，也被称为天主的场所。\par

但必须明白，天主性是不可分的，处处以其整体完全临在，不像有形体那样被一部分一部分地分割，而是完全在万物之内，又完全在万物之上。\par

\Emph{页边注：关于天使和神灵的地方，以及关于无界限者。}\par

天使虽然不像形体那样以成形的形式被容纳在地方中，但仍被称为在地方中，因为他有心灵的临在，并按自己的本性\OriginalTerm{运作}{energise}，并非他处，而是在他运作之处有其心灵的界限。因为不可能同时在不同地方运作。因为唯有天主才有权能同时处处运作。天使凭借其本性的迅速以及他改变地方的敏捷和速度，在不同地方运作；但\OriginalTerm{天主性}{Deity}，祂处处临在且超越万有，以单一简单的\OriginalTerm{能力}{energy}同时以不同方式运作。\par

此外，灵魂与身体相系，整体与整体，而非部分与部分：它不被身体所容纳，而是如火热透铁般容含身体，且在身体中以其自身的固有\OriginalTerm{能力}{energies}运作。\par

凡被地方、时间或领悟所包含的，是\OriginalTerm{有界限的}{circumscribed}；而不被这些都包含的，是\OriginalTerm{无界限的}{uncircumscribed}。因此，唯有\OriginalTerm{天主性}{Deity}是无界限的，无始无终，包含万有，绝不被领悟。因为唯有祂是不可领悟、无界限的，不在任何人的认识之内，只被祂自己所默观。但天使却在时间方面（因其存在有起始）、地方方面（然而如前所述，是心智空间）和领悟方面都是有界限的。因为他们以某种方式彼此认识对方的本质，且其界限被造物主完美界定。简言之，形体在起始和终结、有形的地方和领悟方面都是有界限的。\par

\Emph{页边注：摘自各处，关于天主、圣父、圣子和圣神。以及关于圣言和圣神。}\par

那么，天主性是完全不变和不易的。因为一切不属我们掌控的事，祂皆以预知预先定了，各按其适当而特定的时间与地方。因此，\BibleQuote{父不审判任何人，但已将一切审判交给了子}。\BibleRef{若 5:22}因为显然，父和子以及圣神皆以天主身份审判。但子自己将降来，以人之形体坐在荣耀的宝座上（因为下降和就坐需要有界限的身体），并以正义审判普世。\par

万物远离天主，不在于地方，而在于本性。在我们身上，思虑、智慧、谋略作为存在状态而来去。在天主身上则不然：因祂没有生成或消逝：祂是不变和不易的：论及祂，不可说偶发性。因为善与\OriginalTerm{本质}{essence}相偕。凡常渴慕天主的，就看见祂：因为天主在万有之内。存有之物依赖那存有者，除非在那存有者内，无物能存在。天主因此与万物相融，维持它们的本性：而在祂的圣肉身中，\OriginalTerm{天主圣言}{God-Word}于\OriginalTerm{自立体}{subsistence}内成为一，并与我们的本性相合，却毫无混淆。\par

\BibleQuote{没有人看见过父，除了圣子和圣神}。\BibleRef{若 6:46}\par

子是父的谋略、智慧和能力。因为论及天主，不可谈论性质，以免暗示祂是\OriginalTerm{本质}{essence}与\OriginalTerm{性质}{quality}的复合体。\par

子来自父，并从父获得所有属性：\Emph{因此祂自己不能做什么}。因为祂没有专属于己且异于父的\OriginalTerm{能力}{energy}。\par

那位本性不可见的天主，借着祂的\OriginalTerm{能力}{energies}成为可见的，我们从世界的组织与管理可知。\BibleRef{智 12:5}\par

子是父的\OriginalTerm{肖像}{image}，圣神是子的肖像，基督借此住在人内，使人肖似祂自己的肖像。\par

圣神是天主，介于\OriginalTerm{非受生}{unbegotten}与\OriginalTerm{受生}{begotten}之间，并借着子与父联合。我们称祂为天主之神、基督之神、基督的心、主之神、主本身、义子之神、真理之神、自由之神、智慧之神（因为祂是这一切的创造者）：以\OriginalTerm{本质}{essence}充满万有，维持万有，以本质充满宇宙，而宇宙却不是祂能力的量度。\par

天主是永恒不变的\OriginalTerm{本质}{essence}，万有的创造者，以虔敬的思念受朝拜。\par

天主也是圣父，永不受生，因非由任何人所生，却生了与自己同永的圣子；天主同样是圣子，永远与圣父同在，由圣父无时间地、永远地、无流变、无受难、也无分离地生出。天主也是圣神，是圣化之能、\OriginalTerm{自立存在}{subsistential}，由圣父无分离地发出，并安息于圣子，与圣父、圣子本质同一。\par

圣言是始终本质性地存在于圣父内的那一位。再者，言也是心灵的天然运动，心灵借此运动、思想、默虑，犹如自身的光与光辉。再者，言是仅在心中讲述的思想。再者，言是作为思想信使的言语。因此，天主是本质性的、\OriginalTerm{实体性存在}{enhypostatic}的圣言；而另外三种言是灵魂的能力，不被视为拥有自身的自立存在。其中第一种是心灵的自然产物，永远自然地由心灵涌出；第二种是思想；第三种是言语。\par

神有多种含义。有圣神；但圣神的德能也被称为神；善天使也是神；魔鬼也是神；灵魂也是神；有时心灵也被称为神。最后，风是神，空气也是神。\par

% structure: p0089 source=h2 target=subsection reason=relative-heading-rank; base=h2

\subsection{第十四章 天主性体的属性}

非受造、无始、不死、无限、永恒、非物质、善、创造、公义、光照、不变、无受难、无界限、不可量、无限制、未定、不可见、不可想、无所欠缺、自为准则与权能、统御万有、赐予生命、全能、无限能力、包容并维持宇宙、为万有提供照顾：这些及类似属性，神性本有，并非从别处领受，而是祂将一切善按各受造物的能力分施给祂自己的造化。\par

各位格彼此寓居并牢固地立于对方之内。因为它们是不可分的，不能彼此分离，却在彼此之内各自循其轨道，不相融合、不相混杂，却互相依附。因为圣子在圣父与圣神内，圣神在圣父与圣子内，圣父在圣子与圣神内，却没有融合、混杂或混淆。并且有同一运动：因为三位格有同一推动、同一运动，这在任何受造性体中都无法观察到。\par

再者，神圣的光辉与\OriginalTerm{能量}{energy}，即是一、单纯、不可分，按其善良在可分的事物中呈现多种形态，并将自身本质的组成部分分配给各物，却仍保持单纯，在分者中无限地倍增，又将分者聚集并归回自身的单纯。因为万有渴慕它，且万有存在其中。它也按各物的本性给予存在，且它自身是存在物的存在，活物的生命，理性者的理性，思想者的思想。但它本身超越思想、理性、生命与本质。\par

再者，天主性体具有渗透万物而不与之混杂的属性，且自身不可被任何事物渗透。此外，还有以单纯的认识知晓万物，并以神圣、遍览、非物质的眼目观看万物，包括现在、过去以及未来的事物，在它们存在之前。\BibleRef{达 2:22}它也是无罪的，并能赦免罪过，带来救恩；凡它所愿意的，都能成就，但并非愿意成就它所能成就的一切。因为它能毁灭宇宙，但它不愿如此行。\par

\SourceRecord{Translated by E.W. Watson and L. Pullan.；Nicene and Post-Nicene Fathers, Second Series；Vol. 9.；Edited by Philip Schaff and Henry Wace.；Buffalo, NY: Christian Literature Publishing Co.,；1899.；<http://www.newadvent.org/fathers/33041.htm>.}

% structure: p0001 source=h1 target=section reason=page-title

\section{\WorkTitle{正统信仰阐释（卷二）}}

% structure: p0002 source=h2 target=subsection reason=relative-heading-rank; base=h2

\subsection{第一章 论永世}

祂创造了永世，祂自己在永世之前就存在，神圣的达味这样向祂说：\Quote{从永世到永世，祢是。}神圣的宗徒也说：\Quote{藉着祂也创造了永世。}\BibleRef{希 1:2}\par

必须理解，永世一词有各种含义，因为它表示许多事物。每个人的生命被称为一个永世。再者，一千年的时期被称为一个永世。再者，今生的全部历程被称为一个永世；未来的生命，即复活后不朽的生命，也被称为一个永世。再者，永世一词被用来表示，不是时间，也不是由太阳的运行和轨道所度量的一部分时间（即由日夜组成的时间），而是那种与永恒同延的时间性的运动和间隔。因为永世之于永恒的事物，正如时间之于暂时的事物。\par

论及这个世界的七个永世，即从天地受造直到人类的普遍终结和复活。因为有一个部分的终结，即每个人的死亡；但也有一个普遍和完全的终结，那时人类普遍复活将要发生。第八个永世是未来的永世。\par

在世界形成之前，当还没有太阳划分昼夜时，并没有一个可以度量的永世，而是存在与永恒同延的那种时间性的运动和间隔。在这种意义上只有一个永世，天主被称为\OriginalTerm{永世}{\GreekText{αἰών}}。\par

\SourceRecord{Translated by E.W. Watson and L. Pullan.；Nicene and Post-Nicene Fathers, Second Series；Vol. 9.；Edited by Philip Schaff and Henry Wace.；Buffalo, NY: Christian Literature Publishing Co.,；1899.；<http://www.newadvent.org/fathers/33042.htm>.}

% structure: p0001 source=h1 target=section reason=page-title

\section{\WorkTitle{正统信仰阐述（第三卷）}}

% structure: p0002 source=h2 target=subsection reason=relative-heading-rank; base=h2

\subsection{第一章 论天主的\OriginalTerm{救恩计划}{Œconomy}及天主对我们的眷顾，并论我们的救恩}

人就这样被恶魔之首的袭击所诱捕，违背了造物主的诫命，被剥夺了恩宠，失去了对天主的信赖，用辛劳生活的荆棘（因为这是无花果叶的含义）遮盖自己，并且被死亡所包裹，即必死性和肉体的粗重（因为这就是皮衣所象征的）；因天主公义的审判，他被逐出乐园，被判死亡，且成为腐朽的奴隶。尽管如此，然而，那位赐予他存在、并仁慈地赐予他幸福生活的天主，因祂的怜悯，并没有弃绝人。相反，祂首先用多种方式训练他，并藉着呻吟和战栗、藉着洪水淹没几乎全族\BibleRef{创 6:13}、藉着语言的混乱和分歧、藉着天使的统治、藉着城市的焚烧、藉着天主的形象显现、藉着战争和胜利与失败、藉着征兆和奇迹、藉着多种才能、藉着法律和先知，一再呼唤他回头。因为天主藉这一切手段，竭力将人从广泛束缚的罪恶奴役中释放出来——罪恶已使生命成为一团不义——并使人回归幸福的生命。因为正是罪恶将死亡如同一头野蛮的野兽引入世界\BibleRef{智 2:24}，以毁灭人的生命。但救赎者必须是无罪的，并非因罪而受死亡支配；此外，祂的本性应得到加强和更新，藉着劳动被训练，并教导那通往从腐朽到永生的美德之路；最终，祂那浩瀚的挚爱之海洋才得以彰显。因为造物主和主亲自为祂亲手所造的工程承担斗争，并藉着辛劳学会成为主人。既然敌人藉着神性的希望诱捕人，他自己反被肉体的帷幕所诱捕，这样便同时显明了天主的良善、智慧、公义和全能。天主的良善显现在祂没有忽视自己手的脆弱，反而动了怜悯之心，向堕落的人伸出援手；祂的公义显现在，当人被击败时，祂没有另立他人战胜暴君，也没有强行从死亡中夺走人，而是以祂的良善和公义，使那因自己罪恶成为死亡奴隶的人，自己再次成为征服者，并以相似之物解救相似之物，尽管这样做极为困难；祂的智慧则显现在祂为困难设计出最恰当的解决方案。因着我们天主圣父的美意，圣父怀中的独生子、圣言和天主\BibleRef{若 1:18}，与圣父和圣神同质，在万世以前就已存在，无始，在起初就已存在，在天主圣父面前，是天主，且具有天主的形状\BibleRef{斐 2:6}，祂屈尊降贵，从天降下，就是说，祂谦卑了祂那无法谦卑的崇高地位，以无法言喻、无法理解的屈尊俯就了祂的仆人（因为这就是\Quote{下降}的含义）。完美的天主成为完美的人，成就了万新之中最新的事\BibleRef{训 1:10}，即太阳底下唯一的新事，藉此显明了天主的无限大能。因为还有什么比天主成为人更伟大的呢？圣言成了血肉，没有改变，由圣神和玛利亚——那位圣洁、永远童贞的天主之母所生。祂成为天主与人之间的中保，唯一的热爱世人者，在童贞女纯洁的子宫中受孕，没有意志或欲望，没有与男子的结合或快感的生产，而是藉着圣神和亚当的首个后裔。祂成为与我们相似的、顺服圣父的那一位，为我们的不顺服找到了补救，这补救来自祂从我们身上所取的，并成为我们顺服的榜样——没有顺服，便不可能获得救恩。\par

% structure: p0004 source=h2 target=subsection reason=relative-heading-rank; base=h2

\subsection{第二章 论圣言成胎的方式及祂神圣的降生成人}

主的使者被派遣到这位出自达味后裔的圣童贞那里\BibleRef{路 1:27}。\BibleQuote{因为我们主明确是从犹大出来的，这支派中没有人曾在祭坛前服务}\BibleRef{希 7:14}，正如神圣的宗徒所说。关于这一点，我们后面会更准确地谈论。使者向她传报喜讯，说：\BibleQuote{万福！你充满恩宠，主与你同在}\BibleRef{路 1:28}。她因这话而惊慌，天使对她说：\BibleQuote{玛利亚，不要害怕，因为你在天主前获得了恩宠。你将怀孕生子，要给祂起名叫耶稣}，因为祂要把自己的民族从他们的罪恶中拯救出来\BibleRef{玛 1:21}。因此，\Quote{耶稣}一词的解释是\Quote{救主}。当她困惑地问：\BibleQuote{这事怎能成就？因为我不认识男人}\BibleRef{路 1:34}？天使又回答说：\BibleQuote{圣神要临到你，至高者的能力要庇荫你。因此，那要诞生的圣者，将称为天主的儿子}\BibleRef{路 1:35}。她对他说：\BibleQuote{我是主的婢女，愿照你的话成就于我吧}。\par

于是，在童贞玛利亚同意之后，圣神便按主藉天使所说的话降临在她身上，净化了她，并赐予她能力以接受圣言的神性，以及产生祂的能力。然后，至高天主的\OriginalTerm{成为位格的}{enhypostatic}智慧与能力，即与圣父\OriginalTerm{同质}{like essence}的天主圣子，作为神圣的种子，荫庇了她，并从她圣洁至极的血液中，形成了具有理性与思想的灵的肉身，即我们\OriginalTerm{复合本性}{compound nature}的初果：不是藉生育，而是藉创造，通过圣神；不是逐步发展身体的样式，而是立刻使之完美，祂自己，即天主圣言本身，与肉身处于\OriginalTerm{位格}{subsistence}的关系中。因为神圣的圣言不是与预先独立存在的肉身结合，而是居住在童贞玛利亚的胎中，祂毫无保留地在自己的位格内，通过永远贞女纯净的血液，为自己取了一个具有理性与思想的灵的肉身之体，如此，祂为自己取了我们人类复合本性的初果，圣言本身成为肉身中的位格。因此，祂同时是肉身，同时是天主圣言的肉身，也是具有理性与思想的灵的肉身。故此，我们不说是人成为了天主，而是天主成为了人。因为祂本性是完美的天主，自然也成为了完美的人；祂没有改变自己的本性，也没有使救恩计划成为空洞的表演，而是毫無混淆、改变或分离地，与那由童贞玛利亚受孕、具有理性与思想、并在祂内获得存在的肉身成为同一位格，同时祂没有将祂神性的本性改变为肉身的本质，也没有将肉身的本质改变为祂神性的本性，也没有从祂的神性本性和祂所采取的人性本性中造成一个复合的本性。\par

% structure: p0007 source=h2 target=subsection reason=relative-heading-rank; base=h2

\subsection{第三章 论基督的\OriginalTerm{二性}{two natures}，反对那些主张祂只有\OriginalTerm{一性}{one nature}的人。}

因为这两个本性彼此结合，没有改变或变化，神性本性没有离开其本来的单纯性，人性也没有被改变为神性本性或归于非存在，也没有从两者中产生一个复合的本性（\OriginalTerm{复合性}{compound nature}）。因为复合的本性不能与组成它的任何本性本质相同，因为它们是从其他东西合成的一物：例如，身体由四种元素组成，但与火、气、水、土本质不同，也不保持这些名称。因此，如果结合后，如同异端者所主张，基督的本性是一个\OriginalTerm{复合一体}{compound unity}，那么祂就从\OriginalTerm{单纯性}{simple nature}变成了复合的本性（\OriginalTerm{复合性}{compound nature}），并且与本性单纯的圣父本质不同，也与母亲（她不是神性与人性的复合体）本质不同。那么祂将既不是神性也不是人性，也不是神也不是人，而只是基督；并且\Quote{基督}一词将不是位格的名称，而是他们所谓的单一本性的名称。\par

然而，我们不认为基督的本性是复合的，也不认为祂是由其他东西组成而与之不同，如同人由灵魂和身体组成，或身体由四种元素组成那样，而是认为，尽管祂由这些不同的部分组成，祂仍是同一个。因为我们承认，祂无论在神性还是人性中，都是完全的天主，同一个存有，并且祂由两个本性组成，存在于两个本性中。此外，我们理解\Quote{基督}一词是\OriginalTerm{位格}{subsistence}的名称，不是作为种类的意义，而是表示两个本性的存在。因为祂在自己的位格中\OriginalTerm{傅油}{anointing}自己：作为天主，用祂自己的神性傅油祂的身体；作为人，被傅油。因为祂自己是神也是人。而傅油就是祂人性的神性。因为如果基督是一个复合的本性（\OriginalTerm{复合性}{compound nature}），与圣父本质相同，那么圣父也必然是复合的，与肉体本质相同，这既荒谬又极其亵渎。\par

同一个本性怎能容纳对立且本质的差异呢？同一个本性怎能同时是\OriginalTerm{受造}{created}与\OriginalTerm{非受造}{uncreated}、\OriginalTerm{可朽}{mortal}与\OriginalTerm{不朽}{immortal}、\OriginalTerm{有限}{circumscribed}与\OriginalTerm{无限}{uncircumscribed}？\par

但如果那些宣称基督只有一个本性的人也说那本性是单纯的，他们必须承认祂要么是纯粹单纯的天主，从而将降生成人视作仅属虚表，要么根据聂斯多略，祂只是一个人。那么祂被说成是\Quote{在完美神性和完美人性中}（\OriginalTerm{在完美神性和完美人性中}{perfect in divinity and perfect in humanity}）又当如何？如果他们主张结合后祂有一个复合的本性，那么什么时候才能说基督有两个本性？因为显然，在结合之前，基督的本性是一个（\OriginalTerm{一个本性}{one nature}）。\par

这就是异端者误入歧途的原因，即他们把本性和位格（\OriginalTerm{本性}{nature}, \OriginalTerm{位格}{subsistence}）视为同一。因为当我们说人类的本性是一个时，请注意，这样说时我们并不考虑灵魂和身体的问题。因为当我们比较灵魂和身体时，不能说它们本性相同。但由于有许多人的位格，而所有位格都具有同一种本性：因为所有位格都由灵魂和身体组成，都分享灵魂的本性，拥有身体的本质，并具有共同的形式；所以我们称这许多不同位格的本性为一个；而每一位格则有两个本性，并在两个本性中实现自己，即灵魂和身体。\par

但在主耶稣基督的情形，不能承认一个共同的形式。因为从未有过、现在没有、将来也绝不会有另一位基督，由神性和人性构成，同时存在于神性和人性中，是完美的天主和完美的人。因此，在主耶稣基督的情形，我们不能说有一个由神性和人性合成的本性，如同个别个体由灵魂和身体合成那样。因为在后一种情形，我们面对的是一个个体，但基督并非一个个体。因为并没有一个可称为基督性的可表谓形式，是祂所拥有的。因此，我们坚持认为，两个完美的本性——一个神性、一个人性——已经结合；不是如受天主咒骂的狄奥斯柯若、欧提基、塞维鲁以及所有那邪恶团伙所说的那样，带着混乱、混淆、混合、掺杂；也不是如天主憎恶的聂斯多略、狄奥多若、摩普绥提亚的狄奥多若及其魔鬼般的族类所说的那样，以位格或关系的方式，或基于尊严、意志一致、同等尊荣、同名、善意；而是经由综合，即按位格，没有变化、混淆、改变、差异或分离；我们宣认在完美的两个本性中，只有一个降生成人的天主圣子的位格；坚持祂的神性和人性同属一个且同一个位格，并承认两个本性在结合后仍得以保存，但我们不认为每个本性是分离且独立的，而是它们在同一个复合的位格中彼此结合。因为我们视这结合为本质性的，即真实而非想象的。我们称之为本质性的，又不是指两个本性产生一个复合的本性，而是指它们真实地结合在天主圣子的一个复合位格中，并且我们坚持它们本质的差异得以保存。因为受造的仍是受造的，非受造的非受造；可朽的仍是可朽的，不朽的不朽；有限制的有限制，无限制的无限制；可见的可见，不可见的不可见。\Quote{一方因奇迹而完全荣耀，另一方则受侮辱。}\par

此外，圣言将人性的属性归为己有：因为凡属于祂圣肉身的一切，都是祂的；祂也藉各部分相互渗透、按位格合一的通信方式，将祂自己的属性赋予肉身，因为那位既是天主又是人、生活并行动、取此形又与此形交往的，是同一位。因此，光荣之主被说是被钉在十字架上\BibleRef{格前 2:8}，尽管祂的神性从未经历十字架；而人子也被承认在受难前已在天上，如主自己所说\BibleRef{若 3:13}。因为光荣之主与那位在本性上和真理上为人子者——即成为人的那位——是同一位；祂的奇迹和苦难我们都知晓，尽管奇迹是按祂天主的能力而行，苦难是按人而受。因为我们知道，正如祂的位格是独一的，本性的本质差异也得以保存。因为如果那彼此相异的事物本身不被保存，差异如何得以保存？差异就是相异之物间的差别。就基督的本性彼此相异而言，即在本质方面，我们承认基督在自己内结合了两个极端：就祂的神性而言，祂与父及圣神相连；就祂的人性而言，祂与母亲及全人类相连。就祂的本性联合而言，我们承认祂一方面与父及圣神不同，另一方面与母亲及其余人类不同。因为本性在祂的位格内联合，具有一个复合的位格，由此祂既与父及圣神不同，也与母亲及我们不同。\par

% structure: p0015 source=h2 target=subsection reason=relative-heading-rank; base=h2

\subsection{第四章 论互相通传的方式}

我们已多次说过，本质是一回事，位格是另一回事；本质表示同类位格的共同和普遍形式，如天主、人；而位格标示个体，即父、子、圣神，或伯多禄、保禄。那么，请注意：神性和人性之名指代本质或本性；而天主和人之名既按本性使用，如我们说天主是不可理解的本质，天主是独一的；也按位格使用，即较具体者被冠以较普遍者的名称，如圣经说：\Emph{因此，天主，你的天主，用喜油傅抹了你}，又如\Emph{在乌兹地方有一个人}\BibleRef{约 1:1}，因为所指的仅是约伯。\par

因此，在我们主耶稣基督的情形，既然我们承认祂有两个本性，却只有一个由两者合成的位格，那么当我们默观祂的本性时，我们讲论祂的神性和人性；但当我们默观由本性合成的位格时，我们有时使用指向祂双重本性的术语，如\Quote{基督}、\Quote{同时是天主和人}、\Quote{降生成人的天主}；有时使用只暗示祂一个本性的术语，如单独\Quote{天主}或\Quote{天主子}，以及单独\Quote{人}或\Quote{人子}；有时使用暗示祂崇高性的名字，有时使用暗示祂卑微性的名字。因为同一位既是天主又是人，祂永恒地由父而来，没有原因，但后来因祂对人的爱而成为人。\par

因此，当我们论及祂的天主性时，我们并不将人性的属性归于天主性。因为我们不说祂的天主性是可受苦难的或被造的。同样，我们也不将天主性的属性归于祂的肉体或人性：因为我们不说祂的肉体或人性是非受造的。但当我们论及祂的位格时，无论我们使用一个同时指称两种性体的名称，还是只指其中一种性体的名称，我们仍将两种性体的属性归于它。因为基督这一名称同时指称两种性体，被称为天主和人、受造和非受造、可受苦难和不可受苦难；当祂仅就其一种性体被称为天主子及天主时，祂仍保留共存性体的属性，即肉体的属性，被称为受苦的天主和被钉的光荣的主\BibleRef{格前 2:8}，这并非就祂是天主而言，而是就祂同时是人而言。同样，当祂被称为人和人子时，祂仍保留天主性体的属性和荣耀，即万古之先的婴孩和不知起始的人；然而，祂之万古之先并非作为婴孩或人，而是作为天主，并在末世成为婴孩。这就是属性互通的方式：由于位格的同一性和各部分之间的相互渗透，每一性体将其属性交换给另一性体。因此，我们可以论基督说：\Quote{这位我们的天主曾在地上显现，与人同住}，以及\Quote{这个人是非受造的、不可受苦难的、且不受限制的。}\par

% structure: p0019 source=h2 target=subsection reason=relative-heading-rank; base=h2

\subsection{第五章　论性体的数目}

因此，就天主性而言，我们宣认只有一个性体，但坚持有三个实际存在的位格，并坚持一切属于本性及本质的事物都是单纯的，并仅在三个属性上认识位格的区别：无因而自有的父性、有因而子的子性、有因而圣神的发引。我们还知道这些位格彼此不可分割、不可分离，合而为一，无混杂地互相渗透。是的，我再说，无混杂地合一，因为它们虽合而为一，仍是三个；虽不可分离，仍是分别的。因为虽然每一个都有独立存在，即是说，是一个完美的位格，并有其个别性，即有特殊的存有方式，但它们在本质、本性属性以及从父的位格不可分离、不可分割方面是一，并且是、也被称为独一天主。同样地，在神圣而不可言喻的、超过一切思想和理解的安排（我指的是圣三独一天主圣言，我们的主耶稣基督的降生成人）中，我们宣认有两个性体，一为天主性，一为人性，彼此结合，并在位格上合一，以致两个性体形成一个复合位格：但我们坚持即使在结合之后，在两个性体在一个复合位格（即在独一基督内）中仍得以保存，并且它们实际存在并拥有其本性属性；因为它们是无混杂地结合，并且被区别和计数而不分离。正如圣三的三个位格无混杂地结合，被区别和计数而不分离（计数不引起它们之间的分裂、分离、疏远或割裂——因为我们承认独一的天主父、子及圣神），同样，基督的各个性体虽然结合，却是无混杂地结合；虽然彼此渗透，却不允许一个变成另一个或互相转化。因为每一个都严格保持其本性的个别性不变。因此，它们可以被计数而不引入分裂。因为数目的本性不是造成分离或合一，而是表明被计数之物的数量，无论它们是合一的还是分离的：例如，当五十块石头砌成一堵墙时，我们有合一；但当五十块石头散落在地时，我们有分离；又如，当我们说煤炭有两种本性，即火和木时，我们有合一，但分离在于火的本性是一回事，木的本性是另一回事；因为这些事物之合一与分离并非由于数目，而是由于另一种方式。因此，正如即使天主性的三个位格彼此合一，我们也不能称它们为一个位格，因为那样会混淆并消除位格之间的区别；同样，即使基督的两个性体在位格上合一，我们也不能称它们为一个性体，因为那样会混淆并消除、甚至毁灭两个性体之间的区别。\par

% structure: p0021 source=h2 target=subsection reason=relative-heading-rank; base=h2

\subsection{第六章　论天主性在其中一个位格中整体地与人性的整体结合，而非部分与部分}

共通而普遍者被用以称谓其所包含的个别者。本质作为形式是共通的，而位格则是个别的。它不是因其拥有本质的一部分而缺少其余部分成为个别，而是在数目上成为个别，如同个体。位格间的差异被认为在于数目，而不在于本质。因此，本质被用来称谓位格，因为在同一形式的每个位格中，本质都是圆满的。因此，位格之间在本质上并无差异，而在于那些确为其特征属性的偶性，但这些特征属于位格而非本性。实际上，他们定义位格为本质连同偶性。所以，位格既包含普遍又包含个别，并且有独立存在；而本质没有独立存在，而是在位格中被认识。因此，当其中一个位格受苦时，整个本质因能受苦而被认为在其一个位格中受苦，正如该位格所受苦的程度，但这并不必然意味着同一类别中所有位格都与受苦的位格一同受苦。\par

因此，我们宣认天主性的本性完全而圆满地在它的每一个位格中，完全在圣父内，完全在圣子内，完全在圣神内。为此，圣父是圆满的天主，圣子是圆满的天主，圣神是圆满的天主。同样，在圣三之中唯一圣言的降生成人中，我们主张在本性的一个位格内，天主性的本性完全而圆满地与全部的人性结合，而非部分与部分结合。诚然，神圣宗徒说：\BibleQuote{在祂内住着整个天主性的圆满，有形体地} \BibleRef{哥 2:9}，即是在祂的肉体内。而祂那深受神圣启迪的门徒狄奥尼修斯，对神圣事物有如此深刻的认识，曾说天主性在自己的一位格中与我们全体共融。但我们不会被驱使而主张神圣天主性的所有位格——即三者——与人性的所有位格结合为一个位格。因为圣父和圣神参与天主圣言的降生成人，无非是基于善意和喜悦。但我们主张，整个天主性的本质与整个人性结合。因为天主圣言在起初塑造我们时，并未省略任何祂赋予我们本性之物，而是全部接纳了它们，包括身体和灵魂，无论是有智力的还是有理性的，以及它们的一切属性。因为缺少其中一样，受造物就不是人。但祂以其圆满，接纳了圆满的我，与整个我结合，以便祂能以恩宠将救恩赐予整个人。因为没有被接纳的，就不能被治愈。\par

天主圣言遂通过理性——它介于天主的纯全与肉身的粗钝之间——与肉身结合。因为理性主宰灵魂和身体，但理性是灵魂中最纯全的部分，而天主则是理性的最纯全者。当被那更卓越者允许时，基督的理性显示其自身的权能，但它受制于且顺从于那更卓越者，并实行天主圣意所定之事。\par

再者，理性成了与它结合于位格的天主性的居所，正如身体显然也是如此，但并非如同一位寄居者——这是异端者所犯的亵渎错误，他们声称一个斗不能装两个斗，因为他们以物质的标准判断非物质的。基督如何能被称作圆满的天主和圆满的人，并被说成与圣父和我们在本质上相似，如果只有部分的天主性在祂内与部分的人性结合？\par

此外，我们主张，我们的本性已从死者中复活，升了天，并坐在圣父的右边：并非所有人性个体都已复活并坐在圣父的右边，而是在基督的位格中，这已发生在我们的整个本性上。诚然，神圣宗徒说：\BibleQuote{天主使我们一同复活，并在基督内使我们一同坐在天上} \BibleRef{弗 2:6}。\par

此外，我们主张结合是通过共通的本质而发生的。因为每一个本质对于它所包含的位格都是共通的，并且不能找到一种部分而特殊的本性，即本质：否则我们就必须主张同样的位格同时在其本质上相同又不同，而且圣三在神性方面同时在其本质上相同又不同。所以，同样的本性在每个位格中都被观察到，当我们说圣言的本性成了血肉——如真福亚大纳削和济利禄所说的——我们意指天主性与肉身结合。因此我们不能说：\Quote{圣言的本性受苦了}，因为在祂内的天主性并未受苦，而是我们说人性——决非意指人的所有位格——在基督内受苦，并且我们进一步承认基督在祂的人性内受苦。因此，当我们谈论圣言的本性时，我们所指的正是圣言本身。而圣言既有本质的普遍元素，又有位格的个别元素。\par

% structure: p0028 source=h2 target=subsection reason=relative-heading-rank; base=h2

\subsection{第七章 论天主圣言的单一复合位格}

因此我们持守，\OriginalTerm{神圣位格}{divine subsistence}——即天主圣言的位格——先于万有而存在，无时间、永恒、\OriginalTerm{单纯而不复合}{simple and uncompound}，非受造、无形体、不可见、不可触、无限定，拥有圣父所拥有的一切，因为祂与圣父本质相同，仅在受生方式及与圣父位格的关系上与圣父的位格有别；祂亦是成全的，从未与圣父的位格分离；在末世，祂未离开圣父的怀抱，以无限定的方式，不经种子，以只有祂自己知道的、不可理解的方式，居于童贞玛利亚的胎中，并使那从童贞玛利亚所取的肉身在万世之前就已存在的同一（位格）中得以存在。\par

所以，祂既在万有之内又在万有之上，也居于天主之母的胎中，但在其中是藉着降生成人的功能。因此，祂成了肉身，并藉此取了我们\OriginalTerm{复合性体}{compound nature}的初果，即由理智与理性的灵魂所赋予生命的肉身，以致天主圣言的位格变成了肉身的位格，而圣言的位格，原先是单纯的，成为了复合的，即由两个完美的性体——天主性与人性——所复合，且带有天主圣言天主子嗣的特有属性，藉此祂有别于圣父和圣神；也带有肉身的特有属性，藉此祂有别于母亲和其余的人类；此外，还带有天主性体的属性，藉此祂与圣父和圣神合一，以及人性的标记，藉此祂与母亲和祂与我们相合。再者，祂同时是天主与人，因而有别于圣父、圣神、母亲和我们。因为我们知道这是基督位格的最特别属性。\par

因此，即使在降生成人之后，我们仍宣认祂是天主独一之子，同样亦是人之子，独一基督，独一主，天主独生圣言，独一主耶稣。我们敬拜祂的两次受生：一次是在时间以前、超越原因、理性、时间和本性，由圣父所生；另一次是在末世，为了我们，与我们相似又超越我们。为了我们，因为是为我们的救恩；与我们相似，因为祂是由女人在足月时所生的人；超越我们，因为不是藉着种子，而是藉着圣神和童贞玛利亚，超越了分娩的规律。我们宣认祂并非只是天主而没有我们的人性，也非只是人而剥夺了祂的天主性，亦非两个不同的位格，而是同一个，同时是天主与人，\Quote{完全的天主}和\Quote{完全的人}，\Quote{整个天主}和\Quote{整个人}，同一个既是整个天主，即使祂也是肉身，也是整个人，即使祂也是至高的天主。藉着\Quote{完全的天主}和\Quote{完全的人}，我们意指强调性体的圆满与不匮；藉着\Quote{整个天主}和\Quote{整个人}，我们意指着重位格的单一性与个体性。\par

我们也宣认天主圣言的一个成肉身的性体，照真福济利禄所说，用\Quote{成肉身}一词表达肉身的本质。因此，圣言成了肉身，却并未放弃祂自己的非物质性：祂成了整个肉身，却仍保持完全无限定。就祂是身体而言，祂被缩小而局限于狭窄的界限内，但就祂是天主而言，祂是无限定的，祂的肉身与祂无限定的天主性并无共同广延。\par

因此，祂是整个完全的天主，但不是单纯的天主：因为祂不但是天主，也是人。祂也是整个完全的人，但不是单纯的人，因为祂不但是人，也是天主。因为此处\Quote{单纯地}是指祂的性体，而\Quote{完全地}是指祂的位格；正如\Quote{另一物}是指性体，而\Quote{另一位}是指位格。\par

但是要注意，虽然我们认为主的性体互相渗透，但我们知道这渗透出自天主性体。因为是它渗透并贯穿万物，按祂的意愿，而没有任何东西渗透它；也是它将自身的特有荣耀赋予肉身，同时自身保持不可受情，不参与肉身的情。因为如果太阳将它的能力赋予我们，却不参与我们的能力，那么太阳的创造主和主更当如何呢？\par

% structure: p0035 source=h2 target=subsection reason=relative-heading-rank; base=h2

\subsection{第八章。答覆那些问主的性体是属于连续量还是非连续量的人。}

如果有人问主的性体是否属于连续量或非连续量，我们要说，主的性体既不是一个身体，也不是一个表面，也不是一条线，也不是时间，也不是地点，因此不能归结为连续量。因为这些是连续计量的东西。\par

进一步注意，数目涉及有差异的事物，而毫无差异的事物完全不能被计数；它们正是按照差异而被计数的。例如，伯多禄和保禄就他们是一而言，不是分别被计数的。因为就本质而言，他们是一，不能被说成两个性体；但就位格而言，他们有差异，因而被说成两个位格。因此数目涉及差异，而有差异的对象正是按照彼此差异的程度而被计数的。\par

主的两性，按位格而言，是毫无混淆地结合，并按差异的方法或方式而毫无分离地分开。它们并非按结合的方式被计数，因为我们并不认为基督有两性是指位格方面：而是按毫无分离地分开的方式被计数，因为按差异的方法或方式，基督有两性。因为它们在位格上结合并互相渗透，毫无混淆地结合，各性始终保持其独特而自然的差异。因此，既然它们按差异的方式——且仅按此——被计数，它们必须被归入不连续量之下。\par

因此，基督是唯一、完美的天主和完美的人。我们与圣父及圣神一起以同一敬拜朝拜祂，甚至钦崇祂无玷的肉体，并不认为肉体不配受敬拜：因为它实际上是在圣言的唯一位格中被敬拜，这圣言确实成了肉体的位格。但在这一点上，我们并非敬拜受造物。因为我们朝拜祂，不是作为单纯的肉体，而是作为与神性结合的肉体，并且因为祂的两性被归入天主圣言的同一位和同一位格。我惧怕触及煤炭，因为火与木材相连。我敬拜基督的双重性，因为神性在祂内与肉体相连。因为我并未在天主圣三中引入第四位。绝不允许！但我明认天主圣言及其肉体的一个位，并且天主圣三即使在圣言降生成人后仍是天主圣三。\par

% structure: p0040 source=h2 target=subsection reason=relative-heading-rank; base=h2

\subsection{答问：两性是归入连续量还是不连续量？}

主的两性既不是一个物体，也不是一个面，也不是一条线，也不是地点，也不是时间，以致被归入连续量：因为这些都是连续计量的东西。但主的两性按位格是毫无混淆地结合，并按差异的方法或方式而毫无分离地分开。它们按结合的方式不被计数。因为我们不说基督的两性是两位格或位格上的两个。而是按毫无分离地分开的方式，它们被计数。因为按差异的方法或方式，有两性。因为它们在位格上结合并互相渗透，毫无混淆地结合，既没有彼此改变，而是在结合后各保持其自然的差异。因为受造的仍然是受造的，非受造的非受造的。因此，按差异的方式——且仅按此——它们被计数，因而被归入不连续量。因为凡在各方面互无差异的事物不能被计数，而仅在其相异的程度上被计数；例如，\Person{伯多禄}和\Person{保禄}在他们是同一的方面不被计数：因为在本性上是同一的，他们不是两性，也不如此称呼。但就他们在位格上有差异而言，他们被称为两位格。所以差异是数的原因。\par

% structure: p0042 source=h2 target=subsection reason=relative-heading-rank; base=h2

\subsection{第九章 答问：是否有无位格之性？}

因为虽然没有无位格之性，也没有离位之体（因为实际上，本体和性是在位和位格中被思量的），但未必能推论说，那些在一位格中彼此结合的各性，应各自有其自己的位格。因为当它们结合为一个位格后，可能既不使它们缺乏位格，也不使各有其独特的位格，而是两者共有一个相同的位格。因为既然圣言的同一个位格变成了各性的位格，两者中哪一个都不允许缺乏位格，也不允许有彼此不同的位格，或有时有此性的位格，有时有彼性的位格，而是总是毫无分开或分离地共有同一个位格——一个不分裂为部分或分开的位格，使得一部分属此，一部分属彼，而是完全地、绝对地整体归属于此也整体归属于彼。因为天主圣言的肉体并非作为独立的位格而存在，也没有产生另一个不同于天主圣言的位格，而是它存在于那个位格中，反而成为一个存在于另一位中的位格，而非独立的位格。因此，它既完全不缺乏位格，也并未因此在天主圣三中引入另一个位格。\par

% structure: p0044 source=h2 target=subsection reason=relative-heading-rank; base=h2

\subsection{第十章 论\OriginalTerm{三圣颂}{Trisagium}（\Quote{三圣}）}

既然如此，我们说那狂妄的\Person{彼得·富勒}对\OriginalTerm{三圣颂}{Trisagium}或\Quote{三圣}颂所作的添加是亵渎的；因为它给天主圣三引入了一个第四位，将天主子（祂是圣父的真实自立体能）置于一个单独的位置，并将被钉十字架的那位（好像祂不同于\Quote{大能者}）置于另一个单独的位置，或者好像天主圣三被认为是可以受苦难，圣父和圣神与圣子一同在十字架上受了苦。你们当远离这亵渎而无稽的插入！因为我们持守\Quote{圣天主}一词是指圣父，并不将天主性的称号局限于祂自己，也承认圣子和圣神是天主；\Quote{圣而大能}一词我们归于圣子，并不剥夺圣父和圣神的大能；\Quote{圣而不朽}一词我们归于圣神，并不剥夺圣父和圣子的不朽。因为，实际上，我们仿效神圣宗徒的话，将一切神圣名号单纯而无条件地应用于每一位格。\Emph{\BibleQuote{但是为我们只有一个天主，就是圣父，万物出于他，而我们也归于他；也只有一个主，就是耶稣基督，万物借他而有，我们也借他而有。}}\BibleRef{格前 8:5}。然而，我们跟随神学家\Person{额我略}，他说：\Quote{但是为我们只有一个天主，圣父，万物出于他；一个主耶稣基督，万物借他而有；一个圣神，万物在他内。}因为\Quote{出于他}、\Quote{借他}和\Quote{在他内}这些词并不分割诸本性（因为介词和名称的顺序都不能改变），而是刻画了一个不混淆的本性的属性。这从那事实清楚可见：即它们再次被聚集为一，只要人仔细阅读同一位宗徒的话：\Emph{\BibleQuote{万物出于他、借他、在他内；愿光荣归于他，至于无穷之世。阿们。}}\BibleRef{罗 11:36}。\par

因为\Quote{三圣颂}并非仅指圣子，而是指天主圣三，神圣而圣善的\Person{亚大纳削}、\Person{巴西略}和\Person{额我略}，以及全体受神感召的教父们都作证：因为，事实上，藉着三重圣洁，神圣的色辣芬向我们暗示了超越本质的天主性之三位格。但藉着唯一的宰制，他们表明至极神圣的天主圣三的一个本质和权能。神学家\Person{额我略}确实说：\Quote{这样，至圣者，完全被色辣芬遮盖，并以三个祝圣而受光荣，在一个宰制和一个天主性中结合为一。}这是我们的另一位前辈的最美最高深的哲学。\par

教会史家们说，有一次，当\Place{君士坦丁堡}的百姓向天主献上祈祷以化解一场迫在眉睫的灾祸时，正值\Person{普罗克洛斯}担任总主教之职，有个小孩从人群中被提去，由天使教师们教授了\Quote{三圣}颂：\Quote{圣天主，圣而大能者，圣而不朽者，求你垂怜我们。}当他再次回到地上时，他讲述了自己所学到的，全体百姓都唱这首颂歌，于是迫近的灾祸便化解了。在第四次神圣大公会议，即\Place{加采东}大公会议上，我们被告知这首颂歌正是以这种形式被咏唱的；因为这次神圣会议的记录如此记载。因此，这\Quote{三圣}颂，由天使教导我们，由化解灾祸所证实，由如此众多圣教父的会议所批准和确立，并由色辣芬首先歌唱作为天主性三位格的宣告，竟被篡改，而且俨然按照愚昧的富勒的意见加以修正，仿佛他比色辣芬还高，这是引人嘲笑和戏弄的。但是，哦！何等的傲慢！还不止于愚蠢！然而我们仍然这样宣称，即使魔鬼要把我们撕成碎片：\Quote{圣天主，圣而大能者，圣而不朽者，求你垂怜我们。}\par

% structure: p0048 source=h2 target=subsection reason=relative-heading-rank; base=h2

\subsection{第十一章 论从种相和个体角度看待的本性，以及结合与降生成人的区别：并如何理解\Quote{降生成人的天主圣言的一个本性}。}

本性或被视为纯思维中的抽象（因为它没有独立存在），或同等地出现在同一类属的所有位格中作为它们的结合纽带，然后被称为从种相看待的本性；或普遍地被视为相同，但加上个别属性，在一个位格中，并被称为从个体看待的本性，这与从种相看待的本性是同一的。因此，降生成人的天主圣言并非取了那被视为纯思维中抽象的本性（因为这不是降生成人，只是冒充和虚构的降生成人），也不是从种相看待的本性（因为祂没有取所有位格），而是取了从个体看待的本性，这与从种相看待的本性是同一的。因为祂将我们复合本性的要素取在自己身上，这些要素并非当作独立存在或原来是一个个体、如此被祂所取，而是存在于祂自己的位格中。因为天主圣言的位格本身成为了血肉的位格，因此\BibleQuote{圣言成了血肉}\BibleRef{若 1:14}显然毫无变化，同样，血肉成了圣言而无改变，天主成了人。因为圣言是天主，而人也是天主，由于拥有同一且相同的位格。因此，可以谈论同一事物既是圣言的本性，又是个体中的本性。因为它严格且唯独既不指个体（即位格），也不指位格的共同本性，而是指在其中一个位格中被看待和呈现的共同本性。\par

结合是一回事，降生成人则是另一回事。因为结合只表示联结，而完全不表明与什么结合。但降生成人（这正如有人所说的\Quote{穿上人性}）则表明这联结是与血肉，即与人结合，如同铁的加热表示它与火的结合一样。事实上，真福济利禄本人在解释\Quote{天主圣言降生成人的一个本性}这一短语时，在\WorkTitle{致苏琛苏的第二封书信}中说：\Quote{因为如果我们只说\InnerQuote{圣言的一个本性}，然后沉默，而不加上\InnerQuote{降生成人}这个词，而是可以说完全排除了救恩计划，那么他们假装提出的问题就有一定道理：\InnerQuote{如果一个本性就是全部，那么人性的圆满何在？或者说，与我们相似的本性如何存在？}但是，既然人性的圆满和与我们相似的本质的显示都包含在\InnerQuote{降生成人}这个词中，他们就必须停止依赖一根稻草。}因此，他将圣言的本性置于本性本身之上。因为如果他用本性替代位格，那么省略\Quote{降生成人}就不会荒谬。因为当我们只说天主圣言的一个位格时，我们并没有错。同样，拜占庭的良提乌也认为这一短语指的是本性，而非位格。但在真福济利禄为回应德奥多肋对第二谴责的攻击而写的\WorkTitle{辩护}中，他说：\Quote{圣言的本性，即位格，就是圣言本身。}所以，\Quote{圣言的本性}既不是指位格本身，也不是指\Quote{位格的共同本性}，而是指\Quote{在圣言的位格中整体观之的共同本性}。\par

那么，已经说过，圣言的本性成了血肉，即与血肉结合了；但圣言的本性在血肉中受苦，我们至今从未听说过，尽管我们被教导基督在血肉中受苦。所以，\Quote{圣言的本性}并不意味着\Quote{位格}。因此，剩下的说法是：成为血肉就是与血肉结合，而圣言成了血肉则意味着圣言的位格本身无改变地成了血肉的位格。也说过天主成了人，人成了天主。因为作为天主的圣言无改变地成了人。但天主性成了人，或成了血肉，或穿上人性，我们从未听说过。我们确实学到的是：天主性在其一位格中与人结合，并说天主取了另一种形式或本质，即我们自身的。因为\Quote{天主}这个名适用于每一个位格，但我们不能将\Quote{天主性}一词用于位格。因为我们从未被告知天主性单独是父，或单独是子，或单独是圣神。因为\Quote{天主性}意味着\Quote{本性}，而\Quote{父}意味着位格，正如\Quote{人性}意味着本性，\Quote{伯多禄}意味着位格。但\Quote{天主}指示本性的共同元素，并可派生地用于每一个位格，正如\Quote{人}是那样。因为拥有天主性者是天主，拥有人性者是人。\par

除此之外，请注意：圣父和圣神在圣言降生成人中除了在奇迹方面，以及在善意和目的方面外，一点也不参与。\par

% structure: p0053 source=h2 target=subsection reason=relative-heading-rank; base=h2

\subsection{第十二章。论童贞圣母是真正天主之母：反驳奈斯多略异端者的论证。}

此外，我们宣认童贞圣母是严格真正的天主之母。因为既然由她所生的是真天主，那么这位生了真天主降生成人的她就是真正的天主之母。因为我们坚信天主由她所生，这并非暗示圣言的天主性由她开始存在，而是说天主圣言本身（他由父在万世之前无时间地所生，并与父和圣神无始无终地同在）为了我们的救恩在末世居住在童贞女的胎中，无改变地成了血肉，并由她所生。因为童贞圣母所生的不是单纯的人，而是真天主；并且不是单纯的天主，而是降生成人的天主，他并没有从天上带下他的身体，也没有简单地像通过渠道一样经过童贞女，而是从她那里接受了与我们同性同体并自立存在的血肉。因为如果身体是从天降下，而没有分享我们的本性，那么他成为人有什么用呢？因为天主圣言成为人的目的，正是为了那同一本性（它犯了罪、堕落了、败坏了）能战胜欺骗的暴君，从而从败坏中获得自由，正如神圣的宗徒所说：\BibleQuote{因为死亡既由一人而来，死者的复活也由一人而来}\BibleRef{格前 15:21}。如果第一个是真的，第二个也必然是真的。\par

然而，尽管他说：\BibleQuote{第一个亚当出于地，属于土；第二个亚当是主，来自天上}\BibleRef{格前 15:47}，他并不说他的身体是从天上来的，而是强调他不是单纯的人。因为请注意，他称他既是亚当，又是主，从而表明他的双重本性。因为亚当的解释是\Quote{出于地的}，并且显然，人的本性是出于地的，因为人是由土形成的；但\Quote{主}的称号指明他的天主性体。\par

再者，宗徒说：\BibleQuote{天主派遣了自己的独生子，生于女人}\BibleRef{迦 4:4}。他没有说\Quote{由女人所造}。因此，神圣的宗徒的意思是：天主的独生子和天主，就是那位由童贞玛利亚所生的人，而由童贞玛利亚所生的，就是天主之子与天主。\par

但祂是按照身体的方式出生的，因为祂成了人，并不是预先形成一个人（如同在先知中）然后居住在其中，而是祂自己在本性上和真理上成了人，也就是说，祂使有理智和理性的灵魂所赋予生命的肉身，在自己（的）位格中成存，并且祂自己成为这肉身的位格。因为这就是\BibleQuote{生于女人}的意义。因为天主的圣言自己怎能成为法律的属下，除非祂成了与我们有相同本性的真人呢？\par

因此，我们公正而真实地称圣玛利亚为天主之母。因为这个名字包含了整个救恩计划的奥秘。因为如果生祂的那位是天主之母，那么由祂所生的那位确实是天主，同样也是人。因为那在万世之前的天主，除非祂成了人，怎能由女人所生呢？因为人子显然必须是人自己。但如果由女人所生的那位本身就是天主，那么显然，按神圣且无始的本性的法则由天主圣父所生的那位，与在末世按有始且受时间限制的本性（即人的本性）的法则由童贞女所生的那位，必是同一的。这个名字确实指明我们的主耶稣基督的独一位格、两个本性和两次诞生。\par

但我们从未说圣童贞女是基督之母，因为这是为了废除天主之母的称号，并侮辱那位真实独尊、超乎一切受造的天主之母，才由不洁可憎的犹化者\Person{聂斯多略}（那羞辱的器皿）发明了这个名字作为侮辱。因为\Person{达味}王和司祭长\Person{亚郎}也被称为\Emph{基督}，因为按惯例，君王和司祭是由傅油而立；此外，每位受天主启示的人也可称为\Emph{基督}，但他并非本性是天主；是的，可咒骂的\Person{聂斯多略}称由童贞女所生的那位为\OriginalTerm{携带天主者}{God-bearer}，以此侮辱祂。愿我们绝口不提或想祂只是携带天主者，祂实际上是成了肉身的天主。因为圣言亲自成了血肉，确是由童贞女受孕，但作为神带着所取的本性而出，这本性一经产生，就立即被祂神化了，以致这三件事同时发生：取了我们的人性、进入存在、以及所取的人性被圣言神化。因此，圣童贞女被思考并称为天主之母，不仅是因为圣言的本性，还因为人性的神化，受孕和存在的奇迹同时被成就，即圣言的受孕和肉身在圣言自己内的存在。因为天主之母本身以某种奇妙的方式成了塑造万有之创造者的工具，并将人性赋予万有的天主和创造者，祂神化了祂所取的本性，而结合保存了那些被结合的事物，正如它们被结合一样——也就是说，不仅保存了基督的神性，也保存了祂的人性，不仅是那超越我们的，也是出于我们的。因为祂不是先成为像我们一样，然后才超越我们，而是从祂最初存在开始，就具有双重本性，因为从受孕之初，祂就存在于圣言自己内。因此，按祂自己的本性，祂是人，但以某种奇妙的方式，祂又属于天主，是神圣的。此外，祂具有活肉身的属性：因为由于这救恩计划，圣言领受了这些按照自然运动的秩序真正自然的属性。\par

% structure: p0060 source=h2 target=subsection reason=relative-heading-rank; base=h2

\subsection{第十三章 论两个本性的特性。}

因此，我们宣认同一的耶稣基督、我们的主，是完美的天主和完美的人，我们相信祂具有圣父的一切属性，除了\OriginalTerm{非受生}{ingenerate}；并具有第一亚当的一切属性，除了罪，这些属性是身体和理智理性的灵魂；此外，对应于两个本性，祂具有属于两个本性的两套自然品质：两个自然意志，一个神圣一个属人；两个自然能力，一个神圣一个属人；两个自然自由意志，一个神圣一个属人；以及两种智慧和知识，一个神圣一个属人。因为祂与天主圣父本质相同，祂作为天主自由地意愿和行动；祂也与我们本质相同，祂同样作为人自由地意愿和行动。因为神迹属于祂，受苦的状态也属于祂。\par

% structure: p0062 source=h2 target=subsection reason=relative-heading-rank; base=h2

\subsection{第十四章 论我们的主耶稣基督的意志和自由意志。}

既然基督有两个本性，我们相信祂也有两个自然意志和两个自然能力。但由于祂的两个本性有一个位格，我们相信是同一个位格在两个本性中自然地意愿和行动，祂就是由这两个本性所成、在其中、且就是那一位基督我们的主；而且祂意愿和行动是不分离的，而是作为统一的整体。因为祂在每一形式中都与另一形式紧密共融地意愿和行动。因为本质相同的事物，其意志和能力也相同；而本质不同的事物，其意志和能力也不同；\Emph{反之亦然}，意志和能力相同的事物，本质也相同；意志和能力不同的事物，本质也不同。\par

是故，对于圣父、圣子与圣神，我们因其在意志与能力上的同一，而承认其本性的同一。但在天主救恩计划中，我们因其在意志与能力上的差异，而承认两个本性的差异；正如我们察觉到两个本性的差异，我们也承认意志与能力各不相同。因为正如同一基督的本性数目，当以敬虔之心加以思虑和言说时，并不导致对唯一基督的划分，而只是揭示出即便在结合中，本性间的差异仍得以保持；同样，就其本性而言的意志与能力的数目也是如此。（因为祂为了我们的救恩，在两种本性中都赋有愿意与行动的能力。）这并不引入分裂：决不允许！而只是揭示出，即便在结合中，它们之间的差异也被确保并保存。因为我们坚持意志与能力属于本性，而非属于\OriginalTerm{位格}{subsistence}；我指的是由愿意者与行动者所运用的那些意志与能力。因为如果我们允许它们属于位格，我们便不得不声称天主圣三的三个\OriginalTerm{位格}{subsistences}具有不同的意志与不同的能力。\par

应当注意，意志行动与意志行动的方式并非同一回事。因为愿意是本性的一种官能，正如观看一样，所有人皆拥有之；但意志行动的方式却不取决于本性，而取决于我们的判断，正如观看的方式（无论是好是坏）也是如此。因为并非所有人都以相同方式愿意，也并非所有人都以相同方式观看。这一点我们也要在能力方面予以承认。因为愿意、观看或行动的方式，乃是运用意志、观看与能力官能的模式，仅属于运用它们的人，并藉着普遍认可的差异将其与他人区分开来。\par

因此，单纯的愿意被称为\OriginalTerm{意志能}{volition}或\OriginalTerm{意志能力}{faculty of will}，即一种\OriginalTerm{理性倾向}{rational propension}与\OriginalTerm{自然意志}{natural will}；而特定方式的愿意，或作为意志能基础者，则是\OriginalTerm{意志对象}{object of will}，以及\OriginalTerm{判断意志}{will dependent on judgment}。此外，那内在具有意志能官能的，被称为\OriginalTerm{有能力意愿}{capable of willing}：例如，神性有能力意愿，人性也同样。但行使意志能者，即\OriginalTerm{位格}{subsistence}（例如伯多禄），被称为愿意者。\par

既然基督是唯一的，祂的\OriginalTerm{位格}{subsistence}也是唯一的，那么，无论作为天主还是作为人而愿意的那一位，也是同一且唯一的。并且，既然祂有两个赋有意志的本性（因为凡理性的皆赋有意志与\OriginalTerm{自由意志}{free-will}），我们便应在祂内设定两个\OriginalTerm{意志}{volition}或\OriginalTerm{自然意志}{natural will}。因为祂亲自按照两种本性能够愿意。因为祂承受了那本性上属于我们的意志官能。并且，既然那亲自按照两种本性而愿意的基督是唯一的，我们便在祂的情况设定同一意志对象，并非好像祂只愿意那些祂作为天主自然愿意的事（因为愿意吃喝等并不属于神性），而是祂也愿意那些人性为维持自身所需的事，且不包含任何判断上的对立，而仅仅出自本性的个别性。因为那时祂正是如此自然地愿意，当祂的神性意志如此愿意并容许肉身遭受和行出那属于它自身之事时。\par

至于意志能本性上植根于人，从以下可见。排除神性生活，有三种生命形式：\OriginalTerm{植物性生命}{vegetative life}、\OriginalTerm{感觉性生命}{sentient life}与\OriginalTerm{理智性生命}{intellectual life}。植物性生命的属性是营养、生长与繁衍；感觉性生命的属性是\OriginalTerm{冲动}{impulse}；理性与理智性生命的属性是\OriginalTerm{自由意志}{freedom of will}。因此，如果营养本性上属于植物性生命，冲动属于感觉性生命，那么自由意志本性上属于理性与理智性生命。但自由意志无非就是意志能。因此，\OriginalTerm{圣言}{Word}成了血肉，赋有生命、理智与自由意志，也就赋有\OriginalTerm{意志能}{volition}。\par

再者，\OriginalTerm{自然的}{natural}并非出于训练：因为无人学习如何思想、生活、饥饿、口渴或睡眠。我们也并非学习如何\OriginalTerm{愿意}{willing}：所以愿意是自然的。\par

再者：如果在无理性受造物中，\OriginalTerm{本性}{nature}支配，而在人（他凭自己的\OriginalTerm{自由意志}{free-will}与\OriginalTerm{意志能}{volition}而行动）中，本性被支配，那么，人本性上便是赋有意志能的。\par

再者：如果人是按照有福且\OriginalTerm{超越本质的神性}{super-essential Godhead}的\OriginalTerm{肖像}{image}而受造，并且如果\OriginalTerm{神性本性}{divine nature}本性上赋有\OriginalTerm{自由意志}{free-will}与\OriginalTerm{意志能}{volition}，那么，人作为其肖像，本性上是自由且赋有意志能的。因为\OriginalTerm{教父}{fathers}将自由定义为意志能。\par

再者：如果愿意是每一个人的本性的一部分，而非有些有、有些无，并且如果那被视为众人共有的乃是属于该类个体之本性的特征，那么，人确是本性地赋有意志能。\par

再者：如果本性既不增多也不减少，所有人都同等赋有意志能，而非有人更多、有人更少，那么人本性上是赋有意志能的。既然人本性上是赋有意志能的，主也必本性上是赋有\OriginalTerm{意志能}{volition}的，不仅因为祂是天主，也因为祂成了人。因为正如祂承受了我们的本性，祂也本性上承受了我们的意志。正是如此，\OriginalTerm{教父}{Fathers}说祂在自己内形成了我们的意志。\par

如果意志并非\OriginalTerm{自然的}{natural}，它必是\OriginalTerm{位格的}{hypostatic}或\OriginalTerm{非自然的}{unnatural}。但如果是位格的，那么圣子就必然与圣父有不同的意志：因为属于位格的东西仅是本\OriginalTerm{位格}{subsistence}的特性。而如果是非自然的，那么意志必是本性的缺陷：因为非自然的东西破坏自然的东西。\par

万物的天主圣父要么作为圣父而愿意，要么作为天主而愿意。如果作为圣父，祂的意志将不同于圣子的意志，因为圣子不是圣父。但如果作为天主，圣子是天主，同样圣神是天主，所以\OriginalTerm{意志能}{volition}是祂本性的一部分，即是\OriginalTerm{自然的}{natural}。\par

此外，按照教父们的观点，如果具有同一意志者亦具有同一本质，而基督的天主性与人性具有同一意志，则二者必然也具有同一本质。\par

再者：按照教父们的观点，如果在单一意志中看不出本性的区别，那么，当我们论及一个意志时，就必须停止谈论基督内不同的本性；或者，当我们谈论基督不同的本性时，就必须停止谈论一个意志。\par

此外，神圣福音记载：\BibleQuote{主来到提洛和漆冬境内，进了一家，不愿任何人知道，但祂不能隐藏。}\BibleRef{谷 7:24} 那么，如果祂的天主性旨意是无所不能的，但祂虽愿意却不能隐藏，那么必然是作为人祂愿意而不能，因此祂作为人必须具有意志。\par

再者，福音告诉我们：\BibleQuote{祂来到那地方，说\InnerQuote{我渴}。他们给了祂一些搀了苦胆的醋，祂尝了，却不愿意喝。}那么，如果一方面作为天主祂忍受了口渴并且尝了以后不愿意喝，那么祂作为天主也必然受苦难，因为口渴和品尝都是苦难。但如果祂口渴不是作为天主，而是完全作为人，同样地祂作为人必然具有意志。\par

此外，蒙福的宗徒保禄说：\BibleQuote{祂服从致死，且死在十字架上。}\BibleRef{斐 2:8} 但服从是真实意志的服从，而非虚假意志的服从。因为无理性者不能被称为服从或不服从。但主成为服从圣父的，不是作为天主而是作为人。因为作为天主，祂不能被称为服从或不服从。因为这些是受制于人的事，正如受默感的\Person{国瑞}所说。因此，基督作为人具有意志。\par

然而，当我们断言意志是自然的，我们并不认为它受必然性支配，而是自由的。因为如果它是理性的，它必然绝对自由。因为不仅天主性和非受造的本性不受必然性束缚，而且理智的受造本性也是如此。这是显而易见的：因为天主，按其本性是善、是造物主且是天主，这一切并非出于必然。因为谁能引入这种必然性呢？\par

还需要留意，自由意志有多种含义：一种用于天主，一种用于天使，一种用于人。用于天主时，须以超本质的方式理解；用于天使时，须理解为选择与状态同时发生，且不容任何时间间隔；因为天使虽本性拥有自由意志，却能毫无阻碍地使用它，既无身体的违抗需要克服，也无任何攻击者。再者，用于人时，须理解为在时间上状态先于选择。因为人本性是自由的且拥有自由意志，但他也有魔鬼的攻击来阻碍他，以及身体的冲动；因此，由于攻击和身体的重量，选择在状态之后才出现。\par

那么，如果亚当是凭自己的意志服从，凭自己的意志吃，那么在我们内，意志确实是首先受苦的部分。如果意志首先受苦，而降生成人的圣言没有连同我们其余的本性一同承受这意志，那么就得出我们未从罪中得释放。\par

此外，如果本性中的自由意志能力是祂的作为，而祂却没有承受它，那么祂要么谴责自己的作品为不善，要么吝惜这能力所带来的安慰，从而剥夺了我们的圆满益处，并表明祂自己受制于苦难，因为祂不愿意或不能成就我们完美的救恩。\par

此外，我们不能像\OriginalTerm{位格}{hypostasis}是由两个本性合成那样，说有一个由两个意志合成的复合体。首先，因为合成是位格中的事物的合成，而不是在不同范畴中的、不在其特有范畴中的事物的合成；其次，因为如果我们谈论意志和能力的合成，我们将不得不谈论其他自然属性的合成，例如非受造的和受造的，不可见的和可见的，等等。那么，由两个意志合成的意志将称作什么呢？因为复合体不能以组成它的元素的名称来称呼。否则，我们应把由本性合成的称为本性而非位格。此外，如果我们说在基督内有一个复合意志，我们就使祂在意志上与圣父分离，因为圣父的意志不是复合的。因此，只能说基督的位格是复合且共有的，如同在本性方面一样，在自然属性方面也如此。\par

如果我们想要准确无误，就不能说基督拥有\OriginalTerm{判断}{\GreekText{γνώμη}}和偏好。因为判断是对未知之事经过探究和深思熟虑后，即经过商议和决定后，所做出的决定的一种倾向。判断之后是偏好，它选择并选取其一而非其他。但主不仅是人，也是天主，祂知晓一切，无需探究、调查、商议和决定，祂本性上将所有善的视为自己的，将一切恶的视为外在于自己的。因为依撒意亚先知这样说：\Quote{在孩子知道偏好邪恶之前，他会选择善；因为孩子尚未知道善恶，就因选择善而拒绝邪恶。}因为\Quote{之前}一词表明，祂并非像我们一样通过调查和深思熟虑，而是作为天主，并以神圣的方式存在于肉身中，也就是说，与肉身在自立体上结合，并且因祂的存在本身和无所不包的知识，祂本性上就拥有善。因为诸德是自然的品质，且本性上平等地植入所有人内，即使我们并非都平等地运用我们的自然能力。因着过犯，我们从自然被驱赶到不自然。但主将我们从这不自然带领回自然。这正是\BibleQuote{按我们的肖像，按我们的模样}\BibleRef{创 1:26}的含义。今生的纪律和劳苦，并非设计来使我们获得那本性上陌生的德性，而是使我们能够丢弃那外在于并违反我们本性的邪恶：正如费力地从钢铁上除去并非其本性、而是因疏忽而获得的锈蚀时，我们揭示了钢铁的自然光泽。\par

此外，请注意\Quote{判断}（\OriginalTerm{判断}{\GreekText{γνώμη}}）一词有多种用法和多重含义。有时它表示劝诫：如神圣的宗徒说：\BibleQuote{关于童女，我没有主的命令，但我给出我的判断}\BibleRef{格前 7:25}；有时它表示建议，如先知达味说：\Quote{他们狡猾地商议谋害你的百姓}；有时它表示法令，如我们在达尼尔书中读到：\Quote{这无耻的法令是关于谁（或什么）发出的？}；其他时候，它用在信仰、意见、或目的的意义上，简而言之，\Quote{判断}一词有二十八种不同含义。\par

% structure: p0088 source=h2 target=subsection reason=relative-heading-rank; base=h2

\subsection{第十五章 论我们主耶稣基督内的能量}

此外，我们坚持在我们主耶稣基督内有两种能量。因为祂作为天主，与圣父同性同体，拥有天主的能量；同样，因祂成为人，与我们同性同体，祂也拥有适合人性的能量。\par

但要注意，能量、能量的能力、能量的产物、以及能量的主体都是不同的。能量是本性有效（\OriginalTerm{能动的}{\GreekText{δραστική}}）的本质活动；能量的能力是能量所由出的本性；能量的产物是能量所造成的结果；能量的主体是使用能量的人或自立体。此外，有时\Quote{能量}用在\Quote{能量的产物}的意义上，而\Quote{能量的产物}用在\Quote{能量}的意义上，正如\Quote{创造}和\Quote{受造物}这两个术语有时互换使用。因为我们说\Quote{一切受造}，指的是受造物。\par

也要注意，能量是一种活动，是\Quote{被运作}而不是\Quote{运作}；正如神学家额我略在其关于圣神的论著中所说：\Quote{如果能量存在，它显然必须被运作，而不是运作；一旦被运作，它就会停止。}\par

应当留意，生命本身就是能量，是的，是生物的首要能量，同样，生物的整个体制——营养和生长的功能，即其本性的植物性方面；由冲动激发的运动，即其感觉性方面；以及其理智和自由意志的活动——也是能量。此外，能量是能力的完全实现。如果我们默观在基督内所有这些，那么无疑我们必须也承认祂拥有人的能量。\par

我们最初生起的思维称为活动：这是一种简单的活动，不涉及任何关系，是理智以独立而无形的方式发出其特有的思想，因为若非如此，它就不能恰当地被称为理智。再次，通过清晰言语来揭示和展开思想，也被称为活动。但这不再是那种不涉及关系的简单活动，而是在关系中加以考虑，因为它由思想和言语构成。此外，行事者与所成就之事之间的那种关系本身是活动；而所成就之事本身也被称为活动。第一种只属于灵魂，第二种属于利用身体的灵魂，第三种属于由理智所驱动的身体，最后一种是效果。因为理智预先看见将要发生之事，然后通过身体如此实行。因此，主导权属于灵魂，因为它使用身体作为工具，引导并约束它。但身体的活动大不相同，因为身体是被灵魂引导和推动的。关于效果，对所成就之事的触摸、把握，以及可谓拥抱，属于身体，而赋形与形成则属于灵魂。因此，关于我们的主耶稣基督，行奇迹的能力是祂天主性的活动，而祂双手的工作、意愿以及说：\BibleQuote{我愿意，你洁净了吧} \BibleRef{玛 8:3}，是祂人性的活动。至于效果，擘饼 \BibleRef{若 6:11} 以及癞病人听到\BibleQuote{我愿意}，属于祂的人性，而饼的增多和癞病人的洁净属于祂的天主性。因为借着两者，即通过身体的活动和灵魂的活动，祂显示了同一个、同源的、同等的天主性活动。正如我们看见，祂的性体结合并互相渗透，但我们并不否认它们不同，反而列举它们，虽然我们知道它们不可分离；同样，关于意志和活动，我们知道它们的结合，也承认它们的区别，并加以列举，而不引入分离。因为正如肉身被神化而未改变其本性，同样，意志和活动也被神化而未逾越其固有界限。因为无论祂是此或彼，祂是同一个；无论祂以这种或那种方式愿意和活动，即作为天主或作为人，祂是同一个。\par

那么，我们必须坚持，基督因祂的双重性体而有两种活动。因为本性相异的事物，其活动也不同；活动相异的事物，其本性也不同。反之亦然，本性相同的事物，其活动也相同；活动相同的事物，其本性也相同——这是教父们的观点，他们宣示了神圣的含义。那么，其中之一必为真：要么，如果我们认为基督有一种活动，我们也必须认为祂只有一个性体；要么，如果我们关心真理，并按照福音和教父的教导宣认祂有两个性体，我们也必须宣认祂有与之相应并伴随的两种活动。因为祂在天主性上与天主父性体相同，所以祂在活动上也与父相等。同样，祂在人性上与我们性体相同，所以祂在活动上也与我们相等。因为真福的尼撒主教额我略说：\Quote{凡活动相同的事物，其能力也绝对相同。}因为一切活动都是能力的效果。但受造的本性和非受造的本性不能有同一个本性、能力或活动。然而，如果我们认为基督只有一种活动，我们就会把理智灵魂的激情——即流泪、忧伤和痛苦——归于圣言的天主性。\par

如果他们确实说，圣教父们在论及天主圣三的辩论中曾说：\Quote{凡性体相同的事物，其活动也相同；性体不同的事物，其活动也不同。}并且不宜将关乎神学的事转移到\OriginalTerm{经世}{dispensation}上，我们将回答说：如果这是教父们仅仅针对神学而言的，并且圣子在降生成人后也没有与父相同的活动，那么祂肯定也不能有相同的性体。但我们将这些话归于谁呢？\BibleQuote{我父作工直到现在，我也作工} \BibleRef{若 5:17}；以及\BibleQuote{凡父所作的事，子也照样作}；以及\BibleQuote{你们纵然不信我，也应信这些事}；以及\BibleQuote{我所做的事为我作证}；以及\BibleQuote{父怎样复活死人并赐予生命，子也照样随自己的意愿赐予生命}。因为这一切不仅显示祂即使在降生成人后也与父性体相同，而且也有相同的活动。\par

再者：如果涵盖所有受造物的天主眷顾不仅属于父和圣神，也属于圣子，即使在降生成人后，既然那是活动，那么祂即使在降生成人后也必定与父有相同的活动。\par

但如果我们从奇迹中得知基督与父有相同的性体，而奇迹恰好是天主的活动，那么祂即使在降生成人后也必定与父有相同的活动。\par

但如果祂的天主性和人性有一种活动，那将是复合的活动，并且要么与父的活动不同，要么父也有一种复合的活动。但如果父有复合的活动，那么显然祂也必须有复合的性体。\par

但是，若他们声称，随着能量也引入了位格性，我们要答复说：若位格性随同能量被引入，则真实的逆命题也必成立，即能量也随同位格性被引入；那么，正如有三个位格或自立体，也将有三种能量，或者，正如只有一个能量，也将只有一个位格和一个自立体。的确，圣教父们异口同声地主张：具有相同本质的事物也具有相同的能量。\par

此外，若位格性随同能量被引入，那么那些认为基督的能量既不可说一也不可说二的人，并不主张基督的位格可说一或二。\par

试以火焰剑为例：正如其中火与铁的本性保持区别，它们的两种能量及其效果也是如此。因为铁的能量是切割力，火的能量是燃烧力，切割是铁之能量的效果，燃烧是火之能量的效果；这些在烧焦的切痕和切开的灼伤中保持区别，尽管二者结合后，燃烧离不切割，切割也离不燃烧；我们并不因双重的本性能量而主张有两把火焰剑，也不因火焰剑的统一而混淆能量的本质区别。同样，在基督身上，祂的神性拥有神圣全能的一种能量，而祂的人性拥有如同我们的一种能量。祂人性的效果是祂拉着女孩的手把她带到自己身边，而祂神性的效果是使她复活。这两种效果彼此截然不同，尽管在神人能量中它们彼此不可分离。但是，若因为基督有一个位格，祂就必须有一个能量，那么，因为祂有一个位格，祂也必须有一个本质。\par

再者：若我们认为基督只有一个能量，这能量要么是神性的，要么是人性的，要么都不是。若我们认为是神性的，我们就必须主张祂是独一的天主，剥除了我们的人性。若我们认为是人性的，我们便犯了不敬之罪，说祂不过是人。若我们认为既非神性也非人性，我们也就必须主张祂既非天主也非人，在本质上既不像圣父也不像我们。因为合一的结果是产生了位格的同一性，但本性的差异并未被消除。既然本性的差异得以保存，显然本性的能量也将得以保存。因为不存在缺乏能量的本性。\par

若我们的主基督有一个能量，它必须是受造的或非受造的；因为在这两者之间没有能量，如同没有本性一样。如果它是受造的，它便只指向受造的本性；但如果它是非受造的，它便只表明非受造的本质。因为那属于本性的必须完全与其本性相符：不能存在有缺陷的本性。但与本性和谐的能量并不属于外在之物；这很明显，因为除了与本性和谐的能量之外，任何本性既不能存在也不能被认识。因为借着各样事物彰显其能量的方式，不变性确认了它自己的本性。\par

若基督有一个能量，它必须是同一个能量既施行神性的行动又施行人性的行动。但没有任何东西在保持其自然状态时能以相反的方式行动：因为火不会冻结和煮沸，水也不会干燥和弄湿。那么，那位性本为天主、又成为人的，怎能以同一个能量既施行奇迹又承受苦难呢？\par

那么，如果基督接纳了人的心灵，即理智而理性的灵魂，祂无疑拥有思想，并将永远思想。但思想是心灵的能量：因此，基督作为人，被赋予能量，并将如此直到永远。\par

的确，至智、至大、至圣的\Person{若望·基所口}在其\WorkTitle{宗徒大事录注释}的第二篇讲道中说：\Quote{若有人甚至称祂的苦难为行动，也不会错：因为在祂承受一切的事上，祂成就了那伟大而奇妙的工作，即推翻死亡，以及祂所有的其他工作。}\par

如果一切能量都被定义为某种本性的本质运动，正如精通这些事物的人所说，那么我们在哪里能察觉到任何没有运动、完全缺乏能量的本性呢？或者在哪里能找到不是自然能力的运动的能量呢？但是，正如真福\Person{济利禄}所说，任何有理智的人都不能承认天主和祂的受造物只有一个自然能量。使\Person{拉匝禄}从死者中复活的，并非祂的人性；流泪的，也并非祂的神能：因为流泪是人性特有的，而生命是性位格的生命特有的。然而，它们彼此相通，因为位格的同一性。因为基督是一个，祂的位格或自立体也是一个，但祂有两个本性，一个属于祂的人性，另一个属于祂的神性。的确，从祂的神性自然发出的荣耀，借着位格的同一性成为二者共有的；同样，因祂的肉身，卑微的也成为二者共有的。因为祂是二者之一，即天主或人，是同一的，神性和人性都属于祂自己。因为当祂的神性施行奇迹时，并非没有肉身；当祂的肉身履行卑微的职分时，也并非没有神性。因为祂的神性与受苦的肉身结合，却仍不受苦，承受了救赎的苦难；圣心灵与圣言的能动神性结合，感知并知道所成就的事。\par

因此，祂的神性将其自身的荣耀通传给身体，而神性本身却不参与肉身的苦难。因为祂的肉身并非因神性而受苦，正如神性藉肉身而行动一般。肉身乃作为神性的工具。因此，虽然从最初受孕起，两种形态之间绝无分离，且两种形态的一切行动在全部时间内都成为同一位的行动，但我们绝不混淆那些虽无分离却发生的事，而是从其工作的性质认出某形态具有何种样式。\par

基督，于是按照祂的两性而行动，每一性在祂内都与其他性相通而行动：圣言藉其神性的权威和能力，履行所有属于圣言的行动，即一切至高和主权之行动；而身体则履行所有属于身体的行动，顺服于与它结合的圣言的旨意，圣言已成为它分明的部分。因为祂并非自行趋向于自然的情感，亦非以那种方式回避痛苦之事，祈求脱离它们，或忍受从外而来的遭遇，而是按祂的本性被推动，圣言愿意并容许祂\OriginalTerm{以救恩计划的方式}{œconomically}受那些苦，并做属于祂的事，好使真理得藉本性之工得以证实。\par

此外，正如祂在童贞女的诞生中接受了超本质的本质，同样，祂也以超人的方式显示祂的人性活动：用尘世的脚行走在不安定的水上，并非将水变为土，而是藉其神性的超越能力使水既不流散，也不在物质之脚的重量下屈服。因为祂行人事，并非仅按人的方式：祂不仅是人，也是天主，故连祂的苦难也带来生命与救恩；祂行神事，亦非严格按神的方式，因为祂不仅是天主，也是人，故藉触摸、言语及类似方式施行奇迹。\par

但若有人说：\Quote{我们说基督只有一性，并非为了取消祂的人性活动，而是因为人性活动与神性活动相对，被称为\OriginalTerm{被动}{\GreekText{πάθος}}。}我们将回答：按此推理，那些主张祂只有一性的人，也并非为了取消祂的人性，而是因为人性与神性相对，被称为\OriginalTerm{可被动的}{\GreekText{παθητική}}。但天主禁止我们在区分人性活动与神性活动时，称人性活动为被动。一般而言，没有一样事物的存在是通过比较或对照而被认识或定义的。若果真如此，现存事物就会彼此互为原因。因为若人性活动因神性活动是主动而被称为被动，那么人性也必因神性善而成为恶；并且，若神性活动因人性活动被称为被动而被称为主动，那么神性也必因人性恶而为善。于是，一切受造物都必为恶，那说\BibleQuote{天主看了祂所造的一切，看，都很好}的人就是说谎了\BibleRef{创 1:31}。\par

因此，我们主张，圣教父们依据潜在的概念给予人性活动各种名称。他们称之为能力、活动、区别、运作、属性、特质、被动，并非为了与神性活动区分，而是：能力，因为它是保守且不变的力量；活动，因为它是区别的标志，揭示同一类事物之间的绝对相似；区别，因为它加以区分；运作，因为它彰显；属性，因为它是构成性的且唯独属于它，不属于任何其他；特质，因为它赋予形式；被动，因为它被推动。因为所有出于天主且在天主之后的事物，在被推动方面有所承受，因为它们自身不具有运动或能力。因此，如前所述，它被命名并非为了区分彼此，而是按照宇宙的创造者以创造方式植入其中的计划。因此，当它与神性一同被提及，教父们称它为活动。因为那说\Quote{两种形态在密切相通中各按其活动}的人，与那说\BibleQuote{祂禁食四十天后，就饿了}\BibleRef{玛 4:2}的人所为截然不同（因为祂在愿意时，按自己固有的方式让本性活动），也与那些主张祂内有一种不同的活动、或有两种活动、或时而一种时而另一种的人不同。因为这些说法，随着术语的改变，表明两种活动。的确，数目常藉术语的改变以及称它们为神圣的和人性的而得以指明。因为区别在于不同的事物中，但不存的事物如何区别？\par

% structure: p0113 source=h2 target=subsection reason=relative-heading-rank; base=h2

\subsection{第十六章　答覆那些说\Quote{若人有两性及两种活动，则基督应被认为有三性及同样多的活动}的人}

每一个体的人，既然由两性——灵魂与身体——组成，且这两性在他内是不变的，便可适当地被称为两性：因为他在结合后仍保存每一方的自然属性。因为身体并非不朽，而是可朽的；灵魂亦非可朽，而是不朽的；身体非无形，灵魂非肉眼可见：灵魂是理性的、理智的、无形的，而身体是稠密的、可见的、无理性的。但本质彼此对立的事物没有同一性，因此灵魂与身体不能有同一本质。\par

再者：若人是有理性的可朽动物，且每个定义都说明潜在的本性，而理性与可朽按本性计划并不相同，那么人按其自身定义之规则，当然不能有单一的本性。\par

但若人有时被称为具有一个\OriginalTerm{本性}{nature}，则\Quote{本性}一词在此用作\Quote{\OriginalTerm{物类}{species}}之意，如同我们说人与人之间在本性上毫无差别时那样。但由于所有的人都以同样方式被造，由灵魂和身体组成，且每人都有两个不同的本性，因此他们都被归入同一个定义。而这也并非不合理，因为圣亚大纳削在论及一切受造物时，曾称它们具有同一个本性，因为都是受造的；在其\WorkTitle{驳斥亵渎圣神者演讲}中如此表达说：\Quote{圣神超越一切受造物，与受造事物的本性不同，且属于天主性所特有，这一点我们又能看出。因为凡看见为多物所共有、且不在这一个中多而在另一个中少的，就被称为\OriginalTerm{本质}{essence}。既然每一个人都由灵魂和身体组成，因此我们称人具有一个本性。但我们不能称主的\OriginalTerm{位格}{subsistence}为一个本性，因为每一本性即使在结合后仍保持其本性特质，且我们无法找到一类基督。因为没有别的基督既由天主性又由人性而生，同时既是天主又是人。}\par

再者：人在\OriginalTerm{物类}{species}上的统一，与灵魂和身体在\OriginalTerm{本质}{essence}上的统一并不相同。因为人在物类上的统一显明了所有人之间的绝对相似，而灵魂和身体在本质上的统一则是对其本身存在的侮辱，并将它们归于虚无：因为要么一方必须转变为另一方的本质，要么从不同的事物中产生出某种不同之物，于是两者都将改变；或者，如果它们保持各自的界限，则必然有两个本性。因为就本质的本性而言，有形体的与无形体的并不相同。因此，尽管我们认为人有一个本性，不是因为灵魂和身体的本质性质相同，而是因为包括在该物类下的个体完全相同，但我们不必认为基督也有一个本性，因为在这种情况下，并没有一个包含许多位格的物类。\par

再者，每一个复合体都被说是由其直接组成部分所构成。因为我们不说房屋是由土和水构成的，而是由砖和木料构成的。否则，就必须说人至少由五样东西组成，即四大元素和灵魂。同样，对于我们的主耶稣基督，我们不看组成部分的组成部分，而看祂直接由之构成的划分，即\OriginalTerm{天主性}{divinity}和\OriginalTerm{人性}{humanity}。\par

再者，如果因为人说有两个\OriginalTerm{本性}{natures}，我们就必须认为基督有三个，那么，你因为说人由两个本性组成，也必须认为基督由三个本性组成；而对于\OriginalTerm{能力}{energies}也一样，因为能力必须与本性相应。额我略神学家作证说，人被称为具有且确实具有两个本性，他说道：\Quote{天主和人乃是两个本性，因为灵魂和身体确实也是两个本性。}而在他的\WorkTitle{论洗礼}演讲中，他说：\Quote{既然我们由两个部分，即可见的和不可见的本性，即灵魂和身体组成，那么净化也同样是双重的，即借着水和圣神。}\par

% structure: p0120 source=h2 target=subsection reason=relative-heading-rank; base=h2

\subsection{第十七章 论吾主肉身的本性及其意志的神化}

值得注意的是，主的肉身并非因任何本性上的改变、变化、转变或混杂而被神化、成为与天主相等且是天主；正如额我略神学家所说：\Quote{其中一位神化了，另一位被神化了，并且，大胆地说，成为与天主相等；那傅油的成了人，那被傅油的成了天主。}这些话并不表示本性上的任何改变，而是指经世上的结合（我指的是位格上的结合，藉此祂与天主圣言不可分离地结合），以及本性之间的互相渗透，正如我们看见燃烧渗透了钢铁一样。因为正如我们宣认天主毫无变化或改变地成了人，同样我们也认为肉身毫无改变地成了天主。因为圣言成了血肉，并没有超出祂自己天主性的界限，也没有放弃属于祂的天主性荣耀；另一方面，肉身在被神化时，也没有在其本性或其本性特质上发生改变。因为即使在结合之后，两个本性仍保持不混杂，其特质也不受损。但主的肉身通过与圣言最纯洁的结合，即位格上的结合，领受了神圣能力的丰富，而不致损及其任何本性属性。因为它并非凭自己的任何能力，而是借着与祂结合的圣言，显出了神圣的能力；因为燃烧的钢铁之所以燃烧，不是因为它本身被赋予了物理性的燃烧能力，而是因为它通过结合于火而获得了这种能力。\par

因此，同一肉身按本性是可死的，而通过位格上与圣言的结合，又是赋予生命的。我们认为意志的神化也是如此；因为其本性的活动并未改变，而是与祂神圣而全能的意志结合，并成了成为人的天主的意志。因此，祂虽愿意，却不能自行逃脱\BibleRef{谷 7:24}，因为天主圣言乐意使那确实在祂内的人性意志的软弱得以显示。但祂能照自己的意愿洁净癞病人\BibleRef{玛 8:3}，因为与神圣意志的结合。\par

此外还应注意，本性与意志的神化最明确、最直接地指向两个本性和两个意志。因为正如燃烧并不把被烧之物的本性变成火，而是使被烧的和燃烧的各自清晰可辨，并表明不是只有一个，而是有两个本性；同样，神化也不产生一个复合本性，而是两个，以及它们的位格结合。额我略神学家确实说：\Quote{其中一位神化了，另一位被神化了}，而通过\Quote{其中}、\Quote{一位}、\Quote{另一位}这些词，他确实表明了有两个本性。\par

% structure: p0124 source=h2 target=subsection reason=relative-heading-rank; base=h2

\subsection{第十八章 再论意志与自由意志：亦论理智、知识与智慧。}

当我们说基督是完全的天主和完全的人时，我们确实将属于圣父和母亲的所有本性属性都归于祂。因为祂成了人，是为了使那被克服的能够克服。因为那位全能者，在祂全能的权威和能力中，并非没有能力把人从暴君手中拯救出来。但是，如果天主在战胜人之后却用武力对抗他，暴君就有理由抱怨。因此，天主出于怜悯和对人的爱，愿意显露出堕落的人自己作为征服者，并成了人，以相似者恢复相似者。\par

但没有人会否认人是有理性有智力的动物。那么，如果祂取了一个没有灵魂的肉身，或一个没有理智的灵魂，祂怎能成为人呢？因为那并不是人。再者，如果那首先受苦的没有被拯救，也没有被天主性结合所更新和强化，祂成为人对我们又有什么益处呢？因为没有被取纳的，就没有被医治。因此，祂取纳了整个人，甚至是他身上最美好的部分，这部分已经患病，以便祂能将救恩赐予整个人。事实上，绝不可能存在一个没有智慧、缺乏知识的理智。因为如果它没有能力或运动，它就完全化为乌有了。\par

因此，天主圣言，愿意恢复那按自己肖像所造的，便成了人。但按自己肖像所造的是什么呢？难道不是理智吗？所以祂放弃了那更好的，取纳了那较差的。因为理智处于天主与肉身之间的边界，它确实与肉身共融，同时又是天主的肖像。因此，理智与理智相融合，理智居于天主的纯净与肉身的稠密之间。因为如果主取纳了一个没有理智的灵魂，祂就取纳了一个非理性动物的灵魂。\par

但如果福音作者说\BibleQuote{圣言成了血肉}\BibleRef{若 1:14}，请注意，在圣经中，有时人被称为灵魂，例如\Quote{雅各伯带着七十五个灵魂来到埃及}；有时人被称为血肉，例如\Quote{所有血肉都要看见天主的救恩}。因此，主不是成了没有灵魂或没有理智的血肉，而是成了人。事实上，祂自己说：\Quote{你们为什么寻求杀害我，一个把真理告诉你们的人？}\BibleRef{若 8:40}所以，祂取纳了具有理性灵魂和理智的血肉，这理智管辖肉身，但自身隶属于天主圣言的天主性。\par

因此，祂作为天主也作为人，本性上拥有意志的能力。但祂的人性意志顺从并隶属于祂的天主性意志，不是按照自己的倾向引导，而是意志那属于天主性意志所意志的事物。因为祂因着天主性意志的允许，才本性上承受那属于祂自身的事物。因为当祂祈求逃脱死亡时，是祂的天主性意志自然愿意并允许祂如此祈祷、痛苦和恐惧；而当祂的天主性意志定意祂的人性意志选择死亡时，苦难就对祂成为自愿的。因为祂不只是作为天主，也是作为人，自愿交付自己于死亡。因此，祂也赐予我们面对死亡的勇气。所以，在祂救世的苦难之前，祂说：\Quote{父啊！如果可能，请让这杯离开我}，显然如同祂作为人而非作为天主来喝这杯。那么，作为人，祂愿意这杯离开祂：但这些是自然怯懦的话语。\Quote{然而}，祂说，\Quote{不要照我的意志}，这就是说，不是就我与你有不同本质而言，\Quote{但要照你的意志成就}，这就是说，我的意志和你的意志，就我与你同体而言。这些话是勇敢心灵的话语。因为主的神，既然祂真正地按自己的旨意成了人，在初次体验其自然软弱时，就感受到与身体分离的自然共苦，但被天主性的意志所强化，便又勇敢地面对死亡。因为祂自己虽然是全然的天主，但同时也是全然的人；虽然是全然的人，但同时也是全然的天主。祂自己作为人，在祂内并藉着祂自己，将祂的人性服从于天主父，并成为顺服于父的，从而为我们树立了最卓越的典范和榜样。\par

此外，祂自由地运用了祂的天主性意志和人性意志。因为自由意志确实铭刻于每一个理性本性之中。因为如果它不能自由地推理，它拥有理性又有何用呢？因为创造者甚至在无理性的野兽中也植入了自然欲望，迫使它们维持自身的本性。因为缺乏理性，它们无法引导自己的自然欲望，反而被它所引导。因此，一旦对某物的欲望产生，行动的动力也随之而来。因此，它们追求美德并不赢得赞美或幸福，作恶也不受惩罚。但理性本性，虽然拥有自然欲望，却能在遵守自然法则的地方以理性引导和训练它。因为理性的优势在于此，即自由意志，我们以此指理性主体中的自然活动。因此，在追求美德时，它赢得赞美和幸福；在追求邪恶时，它赢得惩罚。\par

因此，主的灵魂凭借其自由意志而意志，但它自由地意志那属于祂的天主性意志所意志它意志的事物。因为肉身并非如梅瑟和众圣者受天象指示那样，由圣言示意而运动。相反，祂自己，虽是一位，却既是天主又是人，按照祂的天主性意志和人性意志来意志。因此，主的两个意志之间的区别不在于倾向，而在于自然能力。因为祂的天主性意志是无始的、全能的，具有与之相随的能力，且不受激情影响；而祂的人性意志在时间上有开始，自身承受了自然的、无罪的激情，并且本性上不是全能的。但由于它真实而自然地源自天主圣言，它又是全能的。\par

% structure: p0132 source=h2 target=subsection reason=relative-heading-rank; base=h2

\subsection{第十九章 论神人合一之能动。}

当荣福\Person{狄约尼削}说基督向我们显示了一种新的\OriginalTerm{神人而神性之能}{theandric energy}时，他并非通过说一种能力来自神性与人性能力的结合而取消了自然能力：因为同样我们可以说一种新的本性来自神性与人性本性的结合。因为，按照圣教父们，凡具有一种能力的也具有一个本质。但他希望指出基督自然能力显示自身的那种新奇而不可言喻的方式，一种适合基督诸本性彼此渗透的不可言喻的方式，并且进一步指出祂的为人生活是多么奇特、奇妙和在本性上未知，最后指出源于不可言喻结合的彼此互换的方式。因为我们主张能力并非分割，诸本性并非分别活动，而是每一方与另一方完全共融地以其自身适当的能力活动。因为人的部分并不仅仅以人的方式活动，因为祂不是纯人；神的部分也并不仅仅以神的方式活动，因为祂不是纯粹的天主，而是同时是天主和人。因为正如在本性方面我们既承认它们的结合又承认它们本性的区别，同样对自然的意志和能力也是如此。\par

因此，请注意，在我们的主\Person{耶稣基督}身上，我们有时谈论祂的两个本性，有时谈论祂的一个位格：二者都指向同一个概念。因为两个本性是一个基督，而一个基督是两个本性。因此，我们说\Quote{基督按照祂的任何一个本性活动}或\Quote{任何一个本性在基督内与另一个共融中活动}是完全相同的。那么，神性在其活动中与肉体共融，因为肉体的受苦和行它所特有的事是由神性意志的善意允许的，并且因为肉体的活动完全是救恩性的，这并非人的能力的属性，而是神性能力的属性。另一方面，肉体在其活动中与圣言的神性共融，因为神性活动是通过身体的器官进行的，并且因为同时作为天主和人活动的那一位是同一个。\par

此外，请注意祂的神圣心智也施行其自然的活动，思考并认知它是天主的心智，被一切受造物朝拜，并记起祂在世上度过的时光和祂所受的一切苦，但祂在其活动中与圣言的神性共融，安排并治理宇宙，思考、认知并安排，并非像单纯的人的心智，而是作为与天主在实体上结合、作为天主的心智而行动。\par

因此，\OriginalTerm{神人而神性之能}{theandric energy}表明，当天主成为人，即当祂降生成人时，祂的人的活动是神圣的，即被神化的，并且并不缺乏祂的神性活动，而祂的神性活动也不缺乏祂的人的活动，而是每一方都与另一方共同被观察。这种说话方式被称为\OriginalTerm{迂回说法}{periphrasis}，即当一句话中包含两件事时。正如在火焰剑的情况下，我们称切割燃烧为一体，燃烧切割为一体，但仍认为切割和燃烧有不同的能力和不同的本性，燃烧具有火的本性，切割具有钢的本性；同样，当我们谈论基督的一个\OriginalTerm{神人而神性之能}{theandric energy}时，我们理解为其两个本性的两种不同能力：属于祂神性的神性能力和属于祂人性的人性能力。\par

% structure: p0137 source=h2 target=subsection reason=relative-heading-rank; base=h2

\subsection{第20章 论自然而无辜的情感。}

那么，我们承认祂承受了人所特有的所有自然而无辜的情感。因为祂承受了完整的人以及人所有的一切属性，除了罪。因为罪不是自然的，也不是造物主植入我们内的，而是由于魔鬼的进一步植入在我们的生活方式中自愿产生的，虽然它不能强制胜过我们。因为自然而无辜的情感是不在我们的能力范围内的，而是由于因过犯而受的惩罚进入了人的生活；例如饥饿、口渴、疲倦、劳苦、眼泪、朽坏、对死亡的退缩、恐惧、带着血汗的痛苦、由于本性的软弱而得到天使的援助，以及其他类似属于每一个人本性的情感。\par

祂承受了一切，为圣化一切。祂受了试探并得胜，为给我们预备胜利，给本性力量去战胜它的对手，为使那从前被战胜的本性能以它自己被战胜的武器去战胜它从前的征服者。\par

因此，那恶者从外面发起攻击，不是通过内在的思想，正如对亚当那样。因为亚当不是被内在思想攻击，而是被蛇攻击。但主击退了攻击，像蒸气一样驱散了它，为使那些攻击祂并被克服的情感能被我们轻易制服，并使新亚当拯救旧亚当。\par

的确，我们的自然情感在基督内既是合乎自然的，又是超乎自然的。因为当祂允许肉体承受它所特有的事时，这些情感按自然方式被激起；但它们超乎自然，因为那自然的事在主内并没有控制意志。因为在祂内没有强制，一切都是自愿的。因为祂是自愿饥饿、口渴、畏惧和死亡的。\par

% structure: p0142 source=h2 target=subsection reason=relative-heading-rank; base=h2

\subsection{第21章 论无知与奴役。}

值得注意的是，祂所取的是无知和奴仆的本性。因为人的本性就是成为天主、其造物主的奴仆，并且人没有对未来的知识。因此，按照神学家\Person{格列高利}的观点，如果你将可见领域与思想领域分开，那么肉体可以说既是奴仆的又无知的；但由于\OriginalTerm{位格同一}{identity of subsistence}和不可分的结合，主的灵魂被赋予了关于未来的知识，并拥有其他奇迹般的能力。因为正如人的肉体就其本性而言并非生命的给予者，而我们的主的肉体——与天主圣言本身\OriginalTerm{位格合一}{union in subsistence}的——尽管没有脱离其本性中的可死性，却因与圣言位格合一而成为生命的给予者，我们不能说它过去不是、现在不是永远生命之给予者；同样，祂的人性本质上并不拥有对未来的知识，但主的灵魂因与天主圣言本身的结合以及位格同一，如我所说，被赋予了关于未来的知识以及其他奇迹般的能力。\par

此外，请注意，我们不能称祂为奴仆。因为奴仆与主人的称谓并非本性的标记，而是表明与某物的关系，如父职与子职的关系。因为这些并不表示本质，而是表示关系。\par

正如我们关于无知所说的，如果你以精细的思考——即精细的想象——将受造者与不受造者分开，那么肉体是奴仆，除非它与天主圣言结合。但是，一旦它在位格中结合，它又怎能是奴仆呢？因为基督是唯一的，祂不能是祂自己的奴仆和主人。因为这些不是简单的断言，而是相对的。那么，祂会是谁的奴仆呢？是祂父亲的吗？那么，圣子就不会有圣父的一切属性，如果祂是圣父的奴仆，而在任何方面都不是祂自己的。此外，宗徒论到我们这些因祂而成为义子的人，又怎能说：\BibleQuote{使你们不再是奴仆，而是儿子}\BibleRef{迦 4:7}，如果祂自己是个奴仆呢？因此，奴仆一词只是作为一个称号使用，并非严格的意义；但为了我们的缘故，祂取了\OriginalTerm{奴仆的形像}{form of a servant}，在我们中间被称为奴仆。因为虽然祂不受苦难，却为了我们成了苦难的奴仆，并成为我们救恩的臣仆。因此，那些说祂是奴仆的人，是将唯一的基督分为两个，正如纳斯托利（\Person{纳斯托利}）所做的那样。但我们宣认祂是万有的主宰和主，唯一的基督，同时是天主和人，且是全知的。\BibleQuote{因为在祂内蕴藏着智慧和知识的一切宝藏，隐藏的财富}\BibleRef{哥 2:3}。\par

% structure: p0146 source=h2 target=subsection reason=relative-heading-rank; base=h2

\subsection{第二十二章 论祂的成长。}

此外，祂被说成在智慧、年龄和恩宠上增长\BibleRef{路 2:52}，在年龄上增长，并借着年龄的增长显出祂内在的智慧；不仅如此，祂也将人类在智慧和恩宠上的进步，以及圣父旨意的完成——即人类对天主的认识和人类的救恩——视为祂自己的增长，处处把属于我们的当作祂自己的。但是，那些认为祂在智慧和恩宠上有所进步、有所增益的人，并没有说结合发生在肉体的最初起源，也没有重视位格合一，而是追随愚昧的纳斯托利，想象某种奇怪的相对合一和单纯内居（\OriginalTerm{相对性的合一}{relative union}，\OriginalTerm{内居}{indwelling}），\BibleQuote{他们不明白自己所说的是什么，也不知道自己所主张的是什么}\BibleRef{弟前 1:1}。因为如果事实上肉体从它的最初起源就与天主圣言结合，或者如果说它存在在祂内并与祂位格同一，那么它怎能不完全赋有一切智慧和恩宠呢？不是它本身领受恩宠，也不是因恩宠分享属于圣言的东西，而是因着位格合一，既然属于人类和属于天主的都属于同一基督，并且祂自己同时是天主和人，如泉水般将祂的恩宠、智慧和一切福祉的圆满倾注于宇宙。\par

% structure: p0148 source=h2 target=subsection reason=relative-heading-rank; base=h2

\subsection{第二十三章 论祂的恐惧。}

恐惧一词有双重含义。因为有一种本性的恐惧，当灵魂不愿意与身体分离时，由于造物主最初植入灵魂中的本性的同情和紧密联系，使得它恐惧并抗拒死亡，祈求逃脱死亡。可以这样定义：\OriginalTerm{本性的恐惧}{natural fear}是使我们紧附存有而畏缩的力量。因为如果万物都是造物主从无中造出的，那么它们都本性上渴望存有而非非存有。此外，对支持存在之物的倾向是它们的一种本有属性。因此，当天主圣言成为人时，祂也有这种渴望；一方面在支持存在之物上显出本性的倾向，渴望饮食和睡眠，并自然地经验了这些；另一方面在导致朽坏之物上显出本性的厌恶，在受难的时刻自愿地面对死亡而畏缩。因为虽然所发生的事是依照自然律，但不像我们那样是出于必然。因为祂自愿地、自发地接受了那属于本性的东西。所以，恐惧本身、战栗和痛苦都属于本性的和\OriginalTerm{无辜的情愫}{innocent passions}，并不受罪恶的统治。\par

还有一种恐惧源于理性和信德的缺失以及对死亡时刻的无知，例如我们在夜间被偶然的响声惊吓而感到恐惧。这是\OriginalTerm{非本性的恐惧}{unnatural fear}，可以这样定义：非本性的恐惧是一种出乎意料的畏缩。我们的主没有取这种恐惧。因此，祂从未感到恐惧，除了在祂受难的时刻，尽管祂按照救恩计划常常体验到畏缩感。因为祂并非不知道指定的时间。\par

但\Person{圣亚大纳削}在反对\Person{亚玻里纳留}的言论中说，祂确实感到了恐惧。\Quote{因此主说：\Emph{\InnerQuote{现在我心神烦乱}}。\BibleRef{若 12:27}\InnerQuote{现在}一词确是指\InnerQuote{当祂愿意的时候}，但同时也指向了实际发生的事。因为祂并非谈论不存在的事，好像它存在似的，仿佛所说的事只是表面发生。因为所有事都是自然且真实地发生的。}随后，在论及其他事之后，他又说：\Quote{祂的神性绝不受无受苦身体的苦难，也不显露出无痛苦和烦忧之灵魂的痛苦与烦忧，也不经受无受苦和祈祷之心灵的痛苦与祈祷。因为，虽然这些事件并非由于本性的颠覆，但它们发生是为了显示祂的真实存有。}\Quote{这些事件并非由于祂本性的颠覆}这句话证明祂并非不自愿地忍受这些事。\par

% structure: p0152 source=h2 target=subsection reason=relative-heading-rank; base=h2

\subsection{第二十四章　论主的祈祷。}

祈祷是心灵向\Person{天主}的提升，或向\Person{天主}祈求合乎义理的事物。那么，我们的主在\Person{拉匝禄}的事上以及在受难时刻是如何祈祷的呢？因为祂的圣心既不需要任何对天主的提升——它早已与天主圣言在自立体上合一——也不需要任何对天主的祈求。因为基督是独一的。但祂之所以如此，是因为祂将我们的人格据为己有，在我们身上印上祂的印记，并成为我们的榜样，教导我们祈求天主并向往祂，并通过祂自己的圣心引导我们走上通往天主的道路。因为正如祂受难，为我们成就了对苦难的胜利，同样祂也祈祷，引导我们（如我所说）走上通往天主的道路，并为我们的缘故\Quote{满全一切义德\BibleRef{玛 3:15}}，正如祂对\Person{若翰}所说，并使祂的圣父与我们和好，尊敬圣父为元始和根源，并证明祂不是天主的仇敌。因为当祂在\Person{拉匝禄}的事上说：\Quote{\Emph{圣父！我感谢祢，因为祢俯听了我。我本来知道祢常常俯听我，但我说这话，是为了四周站立的群众，叫他们相信是祢派遣了我}\BibleRef{若 11:42}}，难道不是向所有人最清楚地表明祂说这话是为了尊敬圣父，视圣父为甚至圣父自己的根源，并显示祂不是天主的仇敌吗？\par

再者，当祂说：\Quote{\Emph{父啊！如果可能，请把这杯从我这里移开；但不要照我的意愿，而要照祢的意愿}\BibleRef{玛 26:39}}，难道不是对所有人显而易见的：祂说这话是为了教导我们在试炼中只向\Person{天主}求助，并以\Person{天主}的意愿为优先，同时也是证明祂确实将我们本性的属性据为己有，并且祂确实拥有两个意愿——自然的，且与祂的本性相配——但彼此绝不反对？\Quote{父}暗示祂与圣父同体，但\Quote{如果可能}并非指祂不知（因为对\Person{天主}来说什么是不可可能的？），而是教导我们将天主的意愿置于我们自己的意愿之上。因为唯独违反天主意愿和许可的事是不可能的。\Quote{但不要照我的意愿，而要照祢的意愿}，因为就祂是天主而言，祂与圣父相同，就祂是人而言，祂彰显了人类的自然意愿。因为正是这自然意愿本能地寻求逃避死亡。\par

此外，祂说\Quote{\Emph{我的天主，我的天主，祢为什么舍弃了我}\BibleRef{玛 27:46}}，这是将我们的人格据为己有。因为除非有人用精细的心智分辨所见与所思，否则天主不会被视为我们的父；祂也从未被祂的神性舍弃；相反，是我们被舍弃和忽视。因此，祂以将我们的人格据为己有的方式献上了这些祈祷。\par

% structure: p0156 source=h2 target=subsection reason=relative-heading-rank; base=h2

\subsection{第二十五章　论据为己有。}

应当注意，有两种\OriginalTerm{据为己有}{appropriation}：一种是自然的和本质的，一种是位格的和相对的。自然的和本质的据为己有，是指我们的主出于对人类的热爱而接受了我们的本性及我们一切自然的属性，成为本性上和真理上的人，并体验了那自然的；而位格的和相对的据为己有，是指当一个人出于怜悯或爱而相对地担当另一个人的位格，并且代替他说出与他自身无关的话。我们的主正是以这种方式将我们的咒骂和我们的被舍弃，以及其他非自然之事据为己有：并非祂本身是或成为这样，而是祂担当了我们的人格，并将自己列为我们中的一员。\Quote{\Emph{为我们成了咒骂}\BibleRef{迦 3:15}}这句话的意义也应如此理解。\par

% structure: p0158 source=h2 target=subsection reason=relative-heading-rank; base=h2

\subsection{第二十六章　论主身体的受难与祂神性的不能受苦。}

因此，天主圣言亲身在肉身上经历了这一切，而祂那唯一不受苦难的神性却保持不受苦难。因为既然独一的基督——由神性和人性复合而成，存在于神性和人性之中——确曾受苦，那么那能受苦的部分自然地受苦了，但那不能受苦的部分并未分担苦难。因为灵魂虽能受苦，却分担身体割伤的痛苦和折磨，尽管它本身并未被割伤，而只是身体被割伤；但神性部分不能受苦，并不分担身体的苦难。\par

观察，此外，我们说是天主在肉身内受苦，但绝不说祂的神性在肉身内受苦，或天主藉肉身受苦。因为，如果太阳照耀在一棵树上，斧头砍伐树木，太阳依然未被劈开且无受情，那么，结合于肉身之位格的圣言的无受情神性，当肉身受苦时，更将保持无受情。若有人将水浇在燃烧的钢铁上，那自然被水损害的（我说的是火）便会熄灭，但钢铁却不受影响（因为钢铁的本质不会被水毁灭）：那么，当肉身受苦时，祂的唯一无受情的神性，尽管与肉身不可分离，也远离一切受情。因为人不应该过分绝对和严格地理解这些比喻：事实上，在比喻中，必须考虑相似与不相似之处，否则它就不是比喻。因为，如果它们在各方面都相似，它们便是同一事物，而不是比喻，尤其在处理神圣事物时更是如此。因为无论是在神学还是救恩计划中，都不能找到一个在各方面都相似的比喻。\par

% structure: p0161 source=h2 target=subsection reason=relative-heading-rank; base=h2

\subsection{第二十七章 论及圣言的神性在我们主死亡时与灵魂和肉身不可分离，且祂的位格仍保持为一。}

由于我们的主耶稣基督是无罪的（因为\Quote{祂没有犯罪，祂除去了世界的罪，口中也没有诡诈}），祂不受死亡支配，因为死亡是因罪而进入世界的。\BibleRef{罗 5:12} 因此，祂死，是因为祂为我们承担了死亡，并且祂为我们将自己作为祭品献给圣父。因为我们得罪了祂，祂必须为我们接受赎价，这样我们才能从定罪中被释放。愿主血的赎价绝不能献给暴君。因此，死亡临近，吞噬肉身如同诱饵，被神性的钩子刺穿，在尝了无罪且赋予生命的肉身后，死亡灭亡，并将它从前所吞噬的一切都吐出来。因为正如黑暗在光出现时消失，同样，死亡在生命的攻击下退却，并为一切带来生命，而为毁灭者带来死亡。\par

因此，虽然祂作为人而死，且祂的圣灵魂与祂无玷的肉身分离，但祂的神性仍与两者（即祂的灵魂和肉身）不可分离，并且即使如此，祂的一个位格并没有分裂为两个位格。因为肉身和灵魂同时在其原初就在圣言的位格中获得了存在，虽然它们因死亡而彼此分离，但各自仍拥有圣言的那个位格。因此，圣言的单一一位格同样是圣言、灵魂和肉身的位格。因为灵魂或肉身从未有过一个与圣言不同的、自身的独立位格，而圣言的位格永远是单一的，从未是两个。所以，基督的位格始终是一个。因为，虽然灵魂在空间上与肉身分离，但在位格上，它们藉着圣言而结合。\par

% structure: p0164 source=h2 target=subsection reason=relative-heading-rank; base=h2

\subsection{第二十八章 论腐败与毁灭}

\Quote{腐败}一词有两种含义。因为它指称所有人类的痛苦，如饥饿、口渴、疲乏、被钉穿、死亡（即灵魂与肉身的分离）等等。在这个意义上，我们说主的肉身曾受腐败。因为祂自愿承受了这一切。但腐败也指身体完全分解为构成元素并彻底消失，许多人也称之为毁灭。主的肉身没有经历这种腐败，正如先知达味所说：\Quote{因为你不会将我的灵魂遗弃在阴府，也不允许你的圣者见到腐败}。\par

因此，像愚昧的朱利安和加伊安那样，说我们主的肉身在复活前（按第一种含义）是不朽的，这是不敬的。因为如果它是不可朽坏的，那么它就不是真实地，而只是表面上与我们同质，而福音告诉我们的那些事，即饥饿、口渴、钉子、肋旁的伤口、死亡，就没有实际发生。但如果它们只是表面上发生，那么救恩计划就是一个骗局和虚假，祂只是表面上成为人，而非实际，我们也只是表面上得救，而非实际。但愿天主禁止，愿那些这样说的人不得分享救恩。但我们已经获得并将获得真正的救恩。但在\Quote{腐败}一词的第二种含义上，我们宣认主的肉身是不朽的，即不可毁灭的，因为这是受启发的教父的传统。事实上，自从我们的救主从死者中复活后，我们说主的肉身即使在第一种含义上也是不朽的。因为主藉着自己的肉身将复活和随后的不朽的恩赐也赐给了我们自己的肉身，祂自己成为我们复活、不朽和无受情的初果。\BibleRef{格前 15:20} 因为正如神圣的宗徒所说：\Quote{这可朽坏的必须穿上不可朽坏的}。\par

% structure: p0167 source=h2 target=subsection reason=relative-heading-rank; base=h2

\subsection{第二十九章 论下降阴府}

当灵魂被神化后，它下降到阴府，以便如同正义的太阳\BibleRef{拉 4:2}为地上的人升起一样，祂也照亮那些坐在黑暗中死荫之下\BibleRef{依 9:2}的人；以便如同祂向地上的人传报和平、向囚犯传报释放、向盲者传报复明，并对于那些相信的人成为永远救恩的创始者，对于那些不相信的人成为他们不信的责备\BibleRef{伯前 3:19}一样，祂也向阴府中的那些人成为同样的：\BibleQuote{凡在天上、地上和地下的一切，都要向祂屈膝}。\BibleRef{斐 2:10} 这样，在释放了那些被捆绑了多年的人之后，祂立即从死者中复活，为我们显示了复活的道路。\par

\SourceRecord{Translated by E.W. Watson and L. Pullan.；Nicene and Post-Nicene Fathers, Second Series；Vol. 9.；Edited by Philip Schaff and Henry Wace.；Buffalo, NY: Christian Literature Publishing Co.,；1899.；<http://www.newadvent.org/fathers/33043.htm>.}

% structure: p0001 source=h1 target=section reason=page-title

\section{正信阐义（第四卷）}

% structure: p0002 source=h2 target=subsection reason=relative-heading-rank; base=h2

\subsection{第一章 关于复活后的事}

基督从死者中复活后，便放下祂的一切苦难，即腐朽、饥饿、口渴、睡眠、疲乏等。因为虽然祂在复活后也尝了食物\BibleRef{路 24:43}，但祂这样做并非由于本性的规律（祂并不感到饥饿），而是基于\OriginalTerm{经世}{economy}，为使我们确信复活是真实的，且是同一肉躯受苦并复活。但祂并未放下本性中的任何部分，无论是身体还是灵魂，祂仍同时拥有身体和灵魂，这灵魂是理智而理性的、有意愿且有\OriginalTerm{机能}{energies}的。如此，祂坐在圣父的右边，以天主和人的双重意志为我们的得救而行动：以天主性行使眷顾、保存和管理万物的\OriginalTerm{机能}{energies}，以人性记住祂在世上的时光，同时祂看见并知晓自己受一切受造理性物崇拜。因为祂的圣神知道自己与天主圣言是\OriginalTerm{同体}{substance}的，并作为天主之神而非单纯之神，分享应归于祂的崇拜。此外，祂从地上升至天上，又从天上降下，是祂有界限身体的\OriginalTerm{机能}{energies}的表现。\BibleQuote{因为祂要这样再来，就如你们看见祂升天一样。}\BibleRef{宗 1:11}\par

% structure: p0004 source=h2 target=subsection reason=relative-heading-rank; base=h2

\subsection{第二章 关于坐在圣父右边}

此外，我们相信基督以身体坐在天主圣父的右边，但我们并不认为圣父的右边是一个实际的地方。因为那不受限制者怎会有被地方限制的右边呢？右边和左边属于受限制者。我们理解圣父的右边是指天主性的荣耀和尊荣，天主子在万世之前即作为天主存在，与圣父\OriginalTerm{本质相似}{of like essence}，最终成了血肉，祂的身体也分享这荣耀，坐在那里。因为祂带着自己的肉躯，由一切受造物同受一个崇拜。\par

% structure: p0006 source=h2 target=subsection reason=relative-heading-rank; base=h2

\subsection{第三章 对那些人说\Quote{如果基督有两个性体，要么你们敬拜受造性体就是服事受造物，要么你们说有一个性体堪受敬拜，另一个不堪受敬拜}的回应}

我们与圣父和圣神共同崇拜天主子，祂在降生成人之前是无形的，如今亲自取了肉身并成为人，同时仍是天主。祂的肉躯，就其本性而言——若是我们对所见和所思作精细的理智区分——并不堪受崇拜，因为它是受造的。但因它与天主圣言结合，便为祂的缘故并在祂内受崇拜。正如国王无论是未穿袍还是穿袍都同样应受尊崇；又如紫袍本身只是一件紫袍，可被践踏和丢弃，但一旦成为王服便获得一切尊荣与荣耀，凡侮辱它的人通常被判死刑；又如木头本身并非不可触摸，但一旦与火接触便成为木炭，变得不可触摸——这不是出于其本性，而是因为与之结合的火，木头的本性并非不可触摸，而是木炭或燃烧的木头那样——同样，肉躯就其本性而言是不堪受崇拜的，但在降生成人的天主圣言中受崇拜，并非因自身，而是因与天主圣言在\OriginalTerm{位格}{subsistence}上的结合。我们不说我们崇拜单纯的肉躯，而是崇拜天主的肉躯，即天主降生成人。\par

% structure: p0008 source=h2 target=subsection reason=relative-heading-rank; base=h2

\subsection{第四章 为何是天主子，而非圣父或圣神，降生成人：以及祂降生成人所成就的}

圣父是父而非子；圣子是子而非父；圣神是神而非父或子。因为\OriginalTerm{个体性}{individuality}是不变的。如果个体性不断变化和改变，它怎能持续存在呢？因此，天主子成了人子，为使其\OriginalTerm{个体性}{individuality}得以保持。因为祂本是天主子，成为人子，由童贞玛利亚取了肉身，并未失去子的\OriginalTerm{子性}{Sonship}。\par

此外，天主子成了人，为能再次赐予人那创造他时所赐的恩惠。因为祂按自己的\OriginalTerm{肖像}{image}创造了人，赋予人理智和自由意志，并按自己的\OriginalTerm{模样}{likeness}，即使人尽可能达到完美的德行。因为以下特性可说是天主本性的印记：无挂虑、无分心、无诡诈、良善、智慧、正义、远离一切恶。因此，在使人与自己\OriginalTerm{共融}{communion}之后（因为祂使人归于不朽，\BibleRef{智 2:23}，并通过与自己的\OriginalTerm{共融}{communion}引导人不朽），当我们因违反诫命而混淆并抹去了天主\OriginalTerm{肖像}{image}的印记，变得邪恶时，我们便被剥夺了与天主的\OriginalTerm{共融}{communion}（\BibleQuote{因为光明与黑暗怎能相通？}\BibleRef{格后 6:14}），并且被逐出生命之外，屈服于死亡的腐朽。是的，既然祂赐予我们分享更好的部分，而我们未能保持它，祂便分享较次的部分，即我们的本性，为能通过祂自己并在祂内，更新那按祂\OriginalTerm{肖像}{image}和\OriginalTerm{模样}{likeness}所造的，也教导我们过德行的生活，通过祂自己为我们铺平道路，并通过生命的交流使我们脱离腐朽，成为我们复活的\OriginalTerm{初果}{firstfruits}，重造那无用和破旧的器皿，召叫我们认识天主，为救我们脱离魔鬼的暴政，并坚固和教导我们如何通过忍耐和谦卑推翻暴君。\par

于是，对魔鬼的崇拜终止了；受造物已被天主圣血圣化；偶像的祭台与庙宇被推倒；天主的认识已植入人心；同体天主圣三、非受造的神性、独一真神、万物之创造主和主宰，接受人的侍奉；德行得到培育；藉着基督的复活，复活的希望已赐下；魔鬼在从前受其控制的人们面前战栗。确实，奇妙的是，这一切都通过祂的十字架、苦难与死亡成功达成。遍及全地，认识天主的福音已被宣讲；没有动用战争、武器或军队来击溃敌人，只有少数赤身露体、贫穷、无学识、受迫害、受折磨的人，他们手持性命，宣讲那位肉身被钉死的主，并战胜了智者与强者。因为十字架的全能之力伴随着他们。死亡本身，曾是人类最大的恐惧，已被推翻；如今，那曾是仇恨与厌恶的对象，反被置于生命之上。这些是基督临在的成就；这些是祂权能的标记。因为祂所拯救的，不是单一民族，就像当年藉梅瑟分开大海，将以色列从埃及和法郎的奴役中解救出来\BibleRef{出 14:16}；相反，祂将全人类从死亡的腐败和罪恶的悲惨暴政中解救出来：不是用强力带领他们走向德行，不是用土压垮他们、用火焚烧他们，或命令用石头砸死罪人，而是以温和与坚忍劝导人选择德行，彼此争胜，并在为获取德行的奋斗中找到喜乐。因为从前，罪人受迫害，却反而更紧贴罪恶，罪恶被他们视为神；而如今，为了虔诚与德行，人们选择了迫害、钉十字架与死亡。\par

万福！基督，天主的圣言、智慧与德能，全能的天主！我们这些无助之人，能回报这一切恩赐什么呢？因为万物都属于祢，祢除了我们的救恩外一无所求，祢亲自赐予这救恩，却又因祢不可言说的美善而感激领受者。感谢祢，赐予我们生命，赐予我们幸福生活的恩宠，并在我们迷失时，藉祢不可言说的屈尊俯就，将我们恢复原状。\par

% structure: p0013 source=h2 target=subsection reason=relative-heading-rank; base=h2

\subsection{第五章　答覆那些问基督的位格是受造或非受造的人}

天主圣言的位格在降生成人之前是单纯非复合的、无形体且非受造的；但在成为肉身后，它也成为肉身的位格，并由祂始终拥有的神性和所取的人性复合而成；祂带有两个性体的特性，在两个性体中得以被认识：以致同一个位格在神性上是非受造的，在人性上是受造的，可见又不可见。否则，我们被迫要么分割独一的基督，讲两个位格；要么否认性体之间的区别，从而引入变化与混淆。\par

% structure: p0015 source=h2 target=subsection reason=relative-heading-rank; base=h2

\subsection{第六章　论基督何时被称为基督}

有人错误断言，在由童贞女降生成人之前，理智已与天主圣言结合，并从那时起称为基督。这是奥力振的荒谬无稽之谈，他主张灵魂先存的学说。但我们认为，天主子与圣言在入住祂圣洁的、终身童贞的母亲胎中之后，成为基督，未改变地成为肉身，且这肉身被神性所傅油。因为这是人性的傅油，正如额我略神学家所说。以下是最神圣的亚历山大济利禄写给德奥多西皇帝的话：\Quote{因为我确实认为，如果天主的圣言无有人性，就不应称之为耶稣基督；如果由女人所出的圣殿不与圣言结合，也不应称之为耶稣基督。因为当天主的圣言以不可言喻的方式在救恩计划的结合中与人结合时，就被理解为基督。}他又这样写信给皇后们：\Quote{有些人认为，\InnerQuote{基督}这名号正确地只赐给由天主圣父所生的圣言，单独看祂自身。但我们并未受教如此思想与言说。因为当圣言成为肉身时，我们才称祂为基督耶稣。因为祂被喜乐之油，即圣神，由那位是天主和父的一位所傅油，因此祂被称为基督。但没有任何习于正确思想的人会怀疑，这傅油是涉及祂作为人的行动。}此外，著名的亚大纳削在他题为\WorkTitle{论救恩显现}的讲词中这样说：\Quote{在肉身寄居之前的天主不是人，而是在天主内的天主，不可见、无感情；但当祂成为人时，祂因着肉身而额外接受了基督的名号，因为确实，苦难与死亡随同这名号而来。}\par

尽管圣经说：\BibleQuote{因此，天主，你的天主，用喜乐之油傅了你}，但应注意，圣经常常用过去时代替将来时，例如这里：\BibleQuote{此后祂在地上显现，并在人间居住}。\BibleRef{巴 3:38}因为当此话说出时，天主尚未被看见，也未住在人间。又如：\BibleQuote{在巴比伦河边，我们坐下，痛哭}。因为那时这些事尚未发生。\par

% structure: p0018 source=h2 target=subsection reason=relative-heading-rank; base=h2

\subsection{第七章　答覆那些问天主之母是否生了两个性体，以及两个性体是否被钉在十字架上的人}

\OriginalTerm{\GreekText{ἀγένητον}}{\GreekText{ἀγένητον}}和\OriginalTerm{\GreekText{γενητόν}}{\GreekText{γενητόν}}，拼写为一个\Quote{\OriginalTerm{\GreekText{ν}}{\GreekText{ν}}}，意为\Quote{非受造}和\Quote{受造}，关涉本性；但\OriginalTerm{\GreekText{ἀγέννητον}}{\GreekText{ἀγέννητον}}和\OriginalTerm{\GreekText{γεννητόν}}{\GreekText{γεννητόν}}，即\Quote{非受生}和\Quote{受生}，正如双写\Quote{\OriginalTerm{\GreekText{ν}}{\GreekText{ν}}}所指示的，关涉的不是本性，而是位格。神圣本性是\OriginalTerm{\GreekText{ἀγένητος}}{\GreekText{ἀγένητος}}，即非受造，但所有在神圣本性之后的事物都是\OriginalTerm{\GreekText{γένητα}}{\GreekText{γένητα}}，即受造的。因此，在神圣而非受造的本性中，\OriginalTerm{\GreekText{ἀγέννητον}}{\GreekText{ἀγέννητον}}或非受生的属性在圣父身上被思及（因为祂不是受生的），\OriginalTerm{\GreekText{γέννητον}}{\GreekText{γέννητον}}或受生的属性在圣子身上（因为祂永远由圣父受生），而\Quote{出发}（procession）的属性则在圣神身上。此外，在每一类生物中，最初的成员是\OriginalTerm{\GreekText{ἀγέννητα}}{\GreekText{ἀγέννητα}}（非受生），但不是\OriginalTerm{\GreekText{ἀγένητα}}{\GreekText{ἀγένητα}}（非受造）；因为它们是由造物主使之存在，而非像它们那样的生物的产物。因为\OriginalTerm{\GreekText{γένεσις}}{\GreekText{γένεσις}}是创造，而\OriginalTerm{\GreekText{γέννησις}}{\GreekText{γέννησις}}或受生，就天主而言，是出自圣父唯一来源的同质圣子的源头；就形体而言，是出自男女接触的同有位格的源头。因此我们看出，受生所指的不是本性，而是位格。因为如果它真指向本性，\OriginalTerm{\GreekText{τὸ} \GreekText{γέννητον}}{\GreekText{τὸ} \GreekText{γέννητον}}和\OriginalTerm{\GreekText{το} \GreekText{ἀγέννητον}}{\GreekText{το} \GreekText{ἀγέννητον}}，即受生和非受生的属性，就不能在同一本性中被思及。因此，天主之母怀了一个在两性中显示出来的位格；一方面，按其神性，出自圣父无时间地受生；另一方面，在最后，由她因时间而成为肉身，且按肉体诞生。\par

但如果我们的追问者暗示由天主之母所生者具有两性，我们回答说：\Quote{诚然！祂具有两性：因为祂在自己的位格内是天主和人。关于十字架受难、复活和升天也应如此说。因为这些所指的不是本性，而是位格。基督既然在两性内，就只在可受苦难的本性内受苦并被钉十字架。因为在十字架上悬挂的是祂的肉体，而非祂的神性。不然，让他们回答我们，当我们问两性是否死了。不，我们将会说。所以两性并未被钉十字架，但基督被生了，也就是说，圣言成为人，在肉身内被生，在肉身内被钉十字架，在肉身内受苦，而祂的神性继续不受苦难。}\par

% structure: p0021 source=h2 target=subsection reason=relative-heading-rank; base=h2

\subsection{第八章 天主的独生子为何被称为首生者。}

首先受生者被称为首生者，无论他是独生子还是众兄弟中的第一个。如果天主子被称为首生者，但并未被称为独生子，我们可能想象祂是受造物中的首生者，作为受造物。但既然祂同时被称为首生者和独生子，这两种含义都必须在祂身上保留。我们说祂是\Quote{一切受造物的首生者}\BibleRef{哥 1:15}，因为祂自己出自天主，受造物也出自天主，但既然祂自己单独且无时间地出自天主圣父的本质，祂有理由被称为独生子、首生者，而非首造者。因为受造物并非出自圣父的本质，而是凭祂的意志出于无。祂也被称为\Quote{在众兄弟中为首生者}\BibleRef{罗 8:29}，因为虽然是独生子，祂也由一位母亲所生。因为，祂如同我们一样，分享了血和肉，成为人，而我们同样藉着祂成了天主的子女，藉着圣洗被收养，那位本性的天主子成了我们中间的首生者，我们则是藉着义子和恩宠成为天主的子女，并与祂处于兄弟的关系。因此祂说：\Emph{我升到我的父和你们的父那里去}。\BibleRef{若 20:17}祂没有说\Quote{我们的父}，而是说\Quote{我的父}，显然指本性的父；又说\Quote{你们的父}，指恩宠的父。以及\Quote{我的天主和你们的天主}。祂没有说\Quote{我们的天主}，而是说\Quote{我的天主}；如果你以精微的思想区分所见之物和所思之物，也还有\Quote{你们的天主}，作为造物主和主。\par

% structure: p0023 source=h2 target=subsection reason=relative-heading-rank; base=h2

\subsection{第九章 论信德与圣洗。}

我们宣认一次圣洗，为得罪之赦与永生。因为圣洗宣示主的死亡。我们实在是\BibleQuote{与主一同埋葬了}\BibleRef{哥 2:12}，正如神圣的宗徒所说。因此，正如我们的主一次而永远地死了，我们也必须一次而永远地受洗，并按照主的话受洗：\Emph{\BibleQuote{因父、及子、及圣神之名}}\BibleRef{玛 28:19}，从而受教宣认父、子及圣神。那么，那些在受洗归入父、子及圣神之后，又受教知道在一个天主性内有三个\OriginalTerm{位格}{subsistences}，却重行受洗的人，正如神圣的宗徒所说，是把基督再次钉在十字架上。\Emph{\BibleQuote{因为，他说，对于那些一度被光照……等等，再使他们更新而悔改是不可能的：因为他们把基督再次钉在十字架上，公然羞辱祂。}}\BibleRef{希 6:4}但那些未受洗归入天主圣三的人，则必须重新受洗。因为虽然神圣的宗徒说：\Emph{\BibleQuote{我们受洗归入基督，归入祂的死亡}}\BibleRef{罗 6:3}，他并非意指圣洗的呼请必须用这些话，而是说圣洗是基督死亡的形象。因为三次浸水，圣洗象征主被埋葬的三天。因此，受洗归入基督，就是信徒归于祂。如果我们未受教宣认父、子及圣神，我们就不能信基督。因为基督是永生天主之子\BibleRef{玛 16:16}，圣父以圣神傅油了祂\BibleRef{宗 10:38}：如神圣的\Person{达味}所言：\Emph{\BibleQuote{因此天主，你的天主，以喜乐之油傅了你，胜过你的同伴}}。\Person{依撒意亚}也以上主的口吻说：\Emph{\BibleQuote{上主的神临于我，因为祂给我傅了油}}\BibleRef{依 61:1}。然而，基督教导祂自己的门徒那呼请，并说：\Emph{\BibleQuote{你们要因父、及子、及圣神之名给他们授洗}}\BibleRef{玛 28:19}。因为既然基督造我们是为了\OriginalTerm{不朽}{incorruption}，而我们违犯了祂救命的诫命，祂就判我们受死亡的腐朽，为使邪恶不致成为不朽；当祂因怜悯俯就祂的仆人，成为相似我们时，祂藉自己的苦难救赎我们脱离腐朽。祂从自己神圣无瑕的肋旁为我们涌出赦免之泉\BibleRef{若 19:34}：水为使我们\OriginalTerm{重生}{regeneration}，洗去罪恶与腐朽；血为使我们饮用，作为永生的保证。祂给我们颁布命令：要借着水和圣神而重生，藉祈祷和呼请，圣神亲近于水。因为人的本性是双重的，由灵魂与肉身组成，祂赐给我们双重的净化：水与圣神。圣神更新我们内那肖似祂肖像的部分，水藉圣神的恩宠洗净肉身脱离罪恶，解救它脱离腐朽。水确实表达死亡的形象，但圣神赐予生命的\OriginalTerm{凭据}{earnest}。\par

因为起初\Emph{\BibleQuote{天主的圣神运行在水面上}}\BibleRef{创 1:2}，并且圣经再次见证水有洁净的能力\BibleRef{肋 15:10}。在诺厄时代，天主用水洗去了世界的罪恶\BibleRef{创 6:17}。按照法律，任何不洁的人都藉水得以洁净，甚至衣服也要用水洗。\Person{厄里亚}在焚烧全燔祭时，将水倒在其上，显明了圣神与水的恩宠。并且按照法律，几乎一切都是用水洁净的：因为看得见的事物是看不见之事的象征。然而，重生是在圣神内发生的：因为信德有能力藉圣神使我们，一如受造物，成为天主的子女，并引导我们进入原初的福乐。\par

因此，罪之赦是藉着圣洗同样赐予众人；但圣神的恩宠则与信德和先前的洁净相称。如今，我们确实藉着圣洗领受了圣神的初果\OriginalTerm{初果}{firstfruits}，而重生为我们是另一生命的开端、印记\OriginalTerm{印记}{seal}、保证和光照\OriginalTerm{光照}{illumination}。\par

因此，我们理当全力以赴，恒常保持自己远离污秽的行为，免得我们\BibleQuote{像狗转回自己的呕吐物}\BibleRef{伯后 2:22}那样，再次成为罪的奴隶。因为\BibleQuote{没有行为的信德是死的，同样，没有信德的行为也是死的}\BibleRef{雅 2:26}。\par

如今，我们受洗归入天主圣三，因为受洗者需要天主圣三来支持并持续，而三个位格\OriginalTerm{位格}{subsistences}不能不彼此同在。因为天主圣三不可分割。\par

第一次圣洗是洪水，为除灭罪。第二次\BibleRef{创 7:17}是通过海洋和云彩：因为云是圣神的象征，海是水的象征。\BibleRef{格前 10:1}第三次是法律的圣洗：因为凡不洁的人用水洗自己，甚至洗自己的衣服，然后进入营中。\BibleRef{肋 14:8}第四次是若翰的圣洗，是预备性的，引导受洗者悔改，使他们信从基督：\Quote{\Emph{我，确实}}，他说，\Quote{\Emph{用水洗你们，但那在我以后来的，要用圣神及火洗你们}}。\BibleRef{玛 3:11}这样，若翰用水洗净是领受圣神的预备。第五次是我们主的圣洗，祂亲自受了洗。祂受洗不是因祂需要洁净，而是将我的洁净当作自己的，为在水上击碎龙的头，为洗去罪并在水中埋葬全部旧亚当，为使施洗者成圣，为满全法律，为启示天主圣三的奥迹，为成为我们圣洗的预像和榜样。而我们也在我们主完美的圣洗中受洗，即水和圣神的圣洗。此外，基督也被说成以火施洗：因为以火舌的形像将圣神的恩宠倾注给祂的圣门徒，正如主亲自说：\Quote{\Emph{若翰确实以水施洗，但你们不多几天后，将要以圣神及火受洗}}。\BibleRef{宗 1:5}或者是因为那将来的火的圣洗，我们将在其中受惩罚。第六次是藉补赎和泪水，这圣洗的确痛苦。第七次是藉血和殉道的圣洗，基督自己为了我们受了这圣洗\BibleRef{路 12:50}，祂是那么崇高和可颂，不可能被任何后来的污点玷污。第八次是最后的圣洗，不是救赎性的，而是毁灭邪恶：因为邪恶和罪不再有势，但它施行无尽的惩罚。\par

此外，圣神以鸽子的形体降下，指示我们圣洗的初果并尊崇身体：因为连这身体也因\OriginalTerm{神化}{deification}而成为天主；而且鸽子从前也常宣布洪水的止息。但圣神以火的形体降临到圣宗徒们身上：因为祂是天主，而\BibleQuote{天主是吞噬之火}。\BibleRef{申 4:24}\par

圣洗圣事中使用橄榄油，作为我们领受傅油的标记，并使我们成为受傅者，且通过圣神向我们宣布天主的怜悯：因为鸽子衔给从洪水中得救者的正是橄榄枝。\BibleRef{创 8:11}\par

若翰受洗时，将手放在他师傅的神圣头顶上，并用自己的血（受洗）。\par

当来者的信德由他们的行为所证实后，我们不应延迟圣洗圣事。因为那欺诈地前来领受圣洗圣事的人，将受到谴责而非益处。\par

% structure: p0034 source=h2 target=subsection reason=relative-heading-rank; base=h2

\subsection{第十章 论信德}

此外，信德是双重的。因为\BibleQuote{信德来自听闻}。\BibleRef{罗 10:17}因为藉着听闻圣经，我们相信圣神的教导。这信德藉着基督所吩咐的一切事而完美：在行为中信从，修养虔诚，并遵行那恢复我们者的命令。因为那不按天主教会传统而信，或藉怪异行为与魔鬼交往的人，是不信者。\par

但再次，\BibleQuote{信德是所希望之事的实质，是未见之事的确证}\BibleRef{希 11:1}，或者是对天主所应许我们的一切以及我们祈祷的美好结果坚定不移、毫不含混的希望。因此，前者属于我们的意志，而后者属于圣神的恩赐。\par

此外，请注意：藉着圣洗圣事，我们脱去我们从出生以来所穿的一切外衣，即罪，并成为属神的以色列人和天主的子民。\par

% structure: p0038 source=h2 target=subsection reason=relative-heading-rank; base=h2

\subsection{第十一章 论十字架，并在此进一步论信德}

\BibleQuote{\Emph{十字架} \Emph{为丧亡的人是愚妄，但为我们得救的人，却是天主的德能}}。\BibleRef{格前 1:23}因为\BibleQuote{属神的人判断一切，但属血气的人不接受圣神的事。}因为这为那些不信赖信德、不思考天主的良善和全能，而以人性及自然推理探究神圣之事的人，乃是愚妄。因为凡属天主的事，都超越本性、理性及想象。因为若有人考虑天主如何以及为何从无中生有，并藉自然推理来达到这一点，他就不领悟。因为这种知识属于邪灵和魔鬼。但若有人在信德指引下，思考天主的良善、全能、真理、智慧和公义，他会发现一切顺遂平坦，道路笔直。\BibleQuote{但若没有信德，不可能得救。}\BibleRef{希 11:6}因为藉着信德，一切事物，无论是人性还是属神的，得以维持。因为若没有信德，农夫不犁地，商人不将生命托付给惊涛骇浪中的一小片木头，婚姻也不缔结，生命中任何其他步骤也不进行。藉着信德，我们确信万物由天主的德能从无变有。我们藉着信德引导一切，无论是神圣还是人间的事。此外，信德是不带任何好事探究的认可。\par

因此，基督的每一项行动和所行的奇迹都是至大、神圣而奇妙的：但其中最奇妙的是祂珍贵的十字架。因为除了我们主耶稣基督的十字架，没有任何事物能征服死亡，消除原祖的罪，劫掠阴府，赐予复活，赋予我们轻视现世甚至死亡本身的能力，预备我们回归昔日的福乐，打开乐园的门，使我们的本性坐在天主的右边，并使我们成为天主的子女和继承人。因为借着十字架，一切都得以矫正。宗徒说：\BibleQuote{我们这许多受了洗归于基督耶稣的人，就是受洗归于祂的死}\BibleRef{罗 6:3}，以及\BibleQuote{凡受洗归于基督的，就是穿上了基督}\BibleRef{迦 3:27}。此外，\BibleQuote{基督是天主的德能，天主的智慧}\BibleRef{格前 1:24}。看！基督的死，即十字架，给我们穿上了天主位格性的智慧与德能。天主的德能就是十字架的道理，或因天主的威力即战胜死亡已借此启示给我们，或因正如十字架的四端被中间的榫钉牢牢固定联结在一起，同样，借着天主的德能，高和深、长和阔，即一切有形的和无形的受造物，得以维持。\par

这被赐予我们作为我们额上的标记，正如割损赐予以色列人一样：因为借此我们信徒与不信者分开并区别出来。这是抵御魔鬼的盾牌、武器和战胜他的凯旋。\BibleQuote{这是那灭命者不可碰你们的记号}\BibleRef{出 12:23}，如经上所说。这是那躺在死亡中者的复活，站立者的支持，软弱者的手杖，羊群的牧棍，热切者的安全通行，前进者的成全，灵魂和身体的救恩，一切邪恶的避免，一切美善的守护，罪恶的除去，复活的植物，永生的树。\par

所以，这同一真正珍贵而可敬的树木，基督曾在其上为我们的缘故将自己献为祭，因接触祂的圣身和圣血而被圣化，理应受敬拜；同样，钉子、长矛、衣服、祂神圣的居所——马槽、山洞、加耳瓦略（带来救恩之地）、赋予生命的坟墓、熙雍——诸教会的主要堡垒等，也应受敬拜。按照天主之父达味的话：\Quote{我们要进入祂的居所，在祂脚凳前敬拜}。而接下来的话清楚表明所指的就是十字架：\Quote{上主，请起来，进入祢的安息之所}。因为复活在十字架之后。如果连我们所爱的房屋、床铺和衣服都令人渴望，那么属于天主、我们的救主、并借以我们真正得救之物，岂不更当渴望吗？\par

此外，我们甚至敬拜那珍贵而赋予生命的十字架的像，即使它是由另一种树木制成的，并非尊崇树木（愿天主禁止），而是尊崇那作为基督象征的像。因为祂告诫祂的门徒说：\BibleQuote{那时人子的记号要出现在天上}\BibleRef{玛 24:30}，意指十字架。同样，复活的天使对妇女说：\BibleQuote{你们寻找那被钉在十字架上的纳匝肋耶稣}\BibleRef{谷 16:6}。宗徒也说：\BibleQuote{我们宣讲被钉在十字架上的基督}\BibleRef{格前 1:23}。因为有许多基督和许多耶稣，但只有一位被钉十字架者。祂没有说被刺，而是说被钉。因此，我们应当敬拜基督的记号。因为无论记号在哪里，祂也在那里。但我们不应敬拜制成十字架像的材料，即使它是黄金或宝石，在那像被毁之后（如果发生的话）。因此，一切奉献给天主的，我们都敬拜，将朝拜归于祂。\par

天主在乐园中所种植的生命树，预表了这珍贵的十字架。因为既然死亡来自一棵树，那么生命和复活也应由一棵树赐予，这是合适的。雅各伯敬拜若瑟的杖头时，首先描绘了十字架，当他交叉双手祝福他的儿子们时见：\BibleRef{希 11:21}，他极为清楚地划了十字记号。同样，梅瑟的杖以十字架的形状击打大海，拯救了以色列，同时使法郎淹没在深渊中；同样，双手交叉伸出击溃阿玛肋克；苦水被一棵树变甜，磐石被击开涌出水流\EditorialNote{《户籍纪》二十}，以及为亚郎表明大司祭尊荣的杖\EditorialNote{《出谷纪》四}；以及在树上被举起如同死了的铜蛇，那树为那些以信心看见仇敌已死的人带来救恩，正如基督被钉在树上，祂取了有罪的肉体，却不知罪。伟大的梅瑟曾呼喊：\Quote{你将要看见你的生命悬挂在你眼前的树上}，依撒意亚同样说：\BibleQuote{我整天向那悖逆和反抗的子民伸出我的手}\EditorialNote{《依撒意亚》65:2}。但愿我们这些敬拜十字架的人，得以分享被钉十字架的基督。阿们。\par

% structure: p0045 source=h2 target=subsection reason=relative-heading-rank; base=h2

\subsection{第十二章 论朝向东方敬拜}

我们朝东方敬拜，并非没有理由或出于偶然。但既然我们由可见和不可见的本性组成，即部分精神部分感官的本性，我们也向造物主呈上双重的敬拜；就如我们既以精神又以肉体的嘴唇歌唱，既以水又以圣神受洗，并以双重的样式与主结合，共享奥迹和圣神的恩宠。\par

既然天主是灵性之光\BibleRef{若一 1:5}，基督在圣经中被称为正义的太阳\BibleRef{拉 4:2}和晨曦，那么东方就必须被指定为朝拜祂的方向。因为一切美善都当归于祂，一切美善皆由祂而来。的确，神圣的达味也说：\Quote{世上的列国，请歌颂天主，请赞美上主，祂乘驾诸天之上的天，向东方。}此外，圣经又说：\Quote{上主天主在伊甸东部栽植了一个乐园，就把他所形成的人安置在那里。}\BibleRef{创 2:8}当他违犯祂的命令时，祂就驱逐了他，使他住在乐园的对面，那显然就是西方。因此，我们寻求并努力返回我们古老的家园，向上主祈祷。此外，梅瑟的帐幕\BibleRef{肋 16:14}的幔子和赎罪盖是朝向东方的。犹太支派作为最珍贵的，也驻扎在东方。\BibleRef{户 2:3}在著名的撒罗满圣殿中，上主之门也是朝东放置的。此外，基督被钉在十字架上时，祂的脸朝向西方，我们也是如此朝拜，追随祂。当祂再次被接升天时，祂被带往东方，宗徒们如此朝拜祂，祂也将如此再来，正如他们看见祂升天时所行的\BibleRef{宗 1:11}；正如主亲自所说：\Quote{闪电从东方发出，直照到西方，人子的来临也要这样。}\BibleRef{玛 24:27}\par

因此，我们期待祂的来临，面向东方朝拜。但这是宗徒们未成文的传统。因为许多传授给我们的传统都是未成文的。\par

% structure: p0049 source=h2 target=subsection reason=relative-heading-rank; base=h2

\subsection{第十三章 论神圣而纯洁的主的奥迹。}

那位良善、全然良善且超越良善、充满良善的天主，因祂美善的极丰富，并未使祂自己（即祂的本性）单独为善，而无其他参与者；相反，为此祂首先创造了属神的精神力量，其次创造了可见、可感的宇宙，然后创造了具有精神与感觉本性的人。因此，祂所造的一切，就其存在而言，都分享了祂的良善。因为祂自身就是万有的存在，万有都在祂内\BibleRef{罗 11:36}，不仅因为祂将他们从无中带出，更因为祂的能力保存并维护了祂所造的一切，尤其是生物。因为无论在其存在，还是在其享有生命方面，他们都分享了祂的良善。但事实上，那些拥有理性的受造物，由于上述原因，也由于他们所拥有的理性本身，分享了更多的良善。因为他们似乎在某种方式上更亲近祂，尽管祂远超过他们。\par

然而，人既被赋予理性和自由意志，就获得了通过自身选择而与天主持续联合的能力，只要他坚守良善，即服从他的造物主。但是，当他违犯了造物主的诫命，陷入死亡与腐朽时，我们族类的创造者和造主，因祂的慈悲心肠，取了我们的模样，成为完全的人，只是没有罪，并与我们的本性结合\BibleRef{希 2:17}。因为祂曾赐予我们祂的肖像和祂的神，我们却没有保持它们，祂便亲自分担我们贫弱的本性，为洗净我们，使我们不朽，并再次立我们为祂天主性的分享者。\par

因为理当不仅我们本性的初果分享更高的美善，而且每个愿意的人也应如此，并应有一次新生，以及适合新生的新的滋养，从而达至圆满的程度。通过祂的诞生（即降生成人）、受洗、苦难和复活，祂将我们的本性从原祖父母的罪恶、死亡和腐朽中解救出来，成为复活的初果，并使自己成为道路、肖像和模范，使我们也能追随祂的足迹，通过义子名分成为祂本性所是的\BibleRef{罗 7:17}：天主的儿子和继承人，并与祂同为继承人。因此，祂赐予我们（如我所说）第二次新生，以便如同我们由亚当所生，是在他的肖像中，承受诅咒和腐朽，同样由祂所生，我们也在祂的肖像中，承受祂的不朽、祝福和荣耀。\par

既然这亚当是属神的，那么新生以及相应的食物也应是属神的；但因为我们有双重和复合的本性，所以新生也应是双重的，食物也应是复合的。因此，我们被赐予了由水和圣神而生的新生：我指的是神圣的圣洗；而食物就是生命之粮，即从天上降下的我们的主耶稣基督\BibleRef{若 6:48}。因为当祂为我们自愿接受死亡时，在祂交付自己的那一夜，祂与圣门徒和宗徒们订立了新约，并通过他们与所有信祂的人订立了新约。于是，在神圣而光荣的熙雍的楼上，祂与门徒们吃完了旧约的逾越节，完成了旧约，然后洗了门徒的脚，作为神圣圣洗的标记。接着，祂擘开饼，递给他们说：\Quote{你们拿去吃，这是我的身体，为你们所擘，为赦免罪过。}同样，祂拿起酒水之杯，递给他们说：\Quote{你们都喝这个，因为这是我的血，新约的血，为你们倾流，以赦免罪过。你们应这样做，来纪念我。因为你们每次吃这饼，喝这杯，就是宣告人子的死亡，承认祂的复活，直到祂再来。}\par

鉴于天主圣言是生活的、有效的\BibleRef{希 4:12}，主完成了祂所愿意的一切；祂说：有光！就有了光，祂说：有穹苍！就有了穹苍；诸天藉主的圣言而建立，其万军由祂口中的气息而成；天地、水、火、空气，以及其中一切的荣耀，实在说，最高贵的受造物——人，也是藉主的圣言而完成的；如果天主圣言凭祂自己的意愿降生成人，并用那位至圣永贞者的纯洁无玷之血，不藉种子而形成了祂的肉躯，那么祂难道不能使饼成为祂的身体，酒水成为祂的血吗？祂在起初说过：\BibleQuote{让大地生出青草}\BibleRef{创 1:11}，直到今日，雨降下时，大地便生出适宜的果实，因着天主的命令而受到催迫和增强。天主说：\Emph{这是我的身体}，和\Emph{这是我的血，你们应这样行，为纪念我}。所以，在祂的全能命令下，\Emph{直到祂来：}因为祂说直到祂来正是这个意思。圣神的荫庇能力藉着呼求成为这场新耕作的雨水。因为正如天主藉圣神的德能创造了祂所创造的一切，同样，现在圣神的德能实行那些超性的事，这些事除非藉着信德，是无法理解的。\Emph{这事怎能成就？}，童贞玛利亚说，\Emph{因为我不认识男人？}总领天使加俾额尔回答她说：\BibleQuote{圣神要临于你，至高者的能力要庇荫你。}\BibleRef{路 1:34}现在你问：饼如何成了基督的身体，酒水如何成了基督的血。我告诉你：\Quote{圣神亲临，并实行那些超越理智和思想的事。}\par

再者，饼和酒被采用：因为天主知道人的软弱：因为人通常厌弃不习惯的事：所以祂以惯常的宽容，通过熟悉的对象施行祂超性的工程。正如在圣洗圣事中，由于人的习惯用水沐浴、用油傅抹，祂将圣神的恩宠与油和水结合，使之成为重生之水；同样，由于人的习惯吃喝水和酒，祂将祂的天主性与这些结合，使之成为祂的身体和血，好使我们通过熟悉和自然的事上升到超性的事。\par

由童贞玛利亚所生的身体，真实地是与天主性结合的身体；并不是那被接升天的身体降下，而是饼和酒本身被变化为天主的身体和血。但若你问这事如何发生，你只需知道它是藉着圣神而成就的，就如主藉着圣神取了那存在于祂内的肉躯，并由天主之母所生。我们除了知道天主圣言是真实的、有效的、全能的以外，别无更多知识，但这事的方式是无法探究的。然而可以这样恰当地说：正如自然中，饼通过吃，酒和水通过喝，被变化为吃者和喝者的身体和血，并且不成为与以前不同的身体；同样，祭台上的饼、酒和水，藉着圣神的呼求和亲临，被超自然地变化为基督的身体和血，不是两个，而是同一个。\par

因此，对于以信德相称地领受的人，它是为赦免罪过、为永生，并为保护灵魂和身体；但对于没有信德而不相称地领受的人，它是为惩罚和刑罚，就如主的死亡对于相信的人成为生命和不朽，以享永福，而对于不相信的人和杀害主的人，则成为永远的惩罚和刑罚。\par

饼和酒不仅仅是基督身体和血的预像（千万不要这样想！），而是主本身的神化身体：因为主说过：\Quote{这是我的身体}，而不是\Quote{这是我身体的预像}；以及\Quote{我的血}，而不是\Quote{我血的预像}。在先前祂曾对犹太人说：\Emph{你们若不吃人子的肉，不喝祂的血，在你们内便没有生命。因为我的肉是真实的食物，我的血是真实的饮料。}又说：\BibleQuote{吃我的人，将要生活。}\BibleRef{若 6:51}\par

因此，让我们以敬畏之心、纯洁的良心和坚定的信德走近，它必将按我们所信的成就我们，毫不怀疑。让我们以灵魂和身体的全部纯洁朝拜它：因为它是双重的。让我们以热切的渴望走近，用交叉成十字形状的手接受被钉者的身体；让我们将眼、唇和额靠近，分受这神圣的炭火，好使我们内心的渴望之火因炭火的附加热量而完全焚尽我们的罪过，光照我们的内心，并使我们在分受这神圣之火时被点燃并被神化。依撒意亚看见了那炭火\BibleRef{依 6:6}。但炭火不是普通的木头，而是与火结合的木头；同样，共融的饼也不是普通的饼，而是与天主性结合的饼。但与天主性结合的身体并非一个本性，而是具有属于身体的一个本性和属于与之结合的天主性的另一个本性，因此这复合体不是一个本性，而是两个。\par

默基瑟德，至高天主的司祭，带着饼和酒迎接亚巴郎从击杀外邦人归来\BibleRef{创 14:18}。那餐桌预象了这一奥秘的餐桌，正如那司祭是基督——真正的最高司祭——的预象和形像。\EditorialNote{肋未纪第十四章}\Emph{你按照默基瑟德的品位，永为司祭。}这饼的预象是供饼。这确实是那纯洁无血的祭献，主曾藉先知说，从日出到日落，人向祂奉献这祭献\BibleRef{拉 1:11}。\par

基督的体血为滋养我们的灵魂和身体而设，不被消耗也不腐朽，不成为排泄物（愿天主不许！），反而为我们的存有和保存，是抵御一切伤害的护卫，是洁净一切污秽的净化：若有人接受劣金，他们以精炼之火净化它，免得将来我们与这世界一同被定罪。它们净化疾病和一切灾祸；依照神圣宗徒的话，\BibleRef{格前 11:31}：\BibleQuote{因为我们若先审断自己，就不至于受审断；但我们受主审判时，是被惩治，免得我们与这世界一同被定罪。}这也是他所说的：\BibleQuote{所以那不相称地吃主的饼、喝主的杯的人，就是干犯主的身、主的血，吃喝自己的罪案。}藉此得以净化，我们与基督的身体及其圣神联合，成为基督的身体。\par

这饼是未来之饼的初果，那饼是\OriginalTerm{日用}{\GreekText{ἐπιούσιος}}，即存在所必需的。因为\OriginalTerm{日用}{\GreekText{ἐπιούσιον}}一词意指未来，即那为将来世代而有的那一位，或那为我们所分享以保存我们本质的那一位。无论取哪一种意义，如此谈论主的身体都是合适的。因为主的肉是赋予生命的神，因它是由赋予生命的圣神所怀的。因为由神所生的，就是神。但我这样说并非要取消身体的本质，而是愿彰显其赋予生命和神圣的能力\BibleRef{若 6:63}。\par

但若有人称饼和酒为主的身体和血的预像，如同受天主启示的巴西略所做，他们这样说不是在祝圣之后，而是在祝圣之前，如此称呼奉献本身。\par

称为\Quote{分受}，因我们藉此分受耶稣的天主性。也称\Quote{共融}，这是真实的共融，因为我们藉此与基督共融，分享祂的肉和祂的天主性：是的，我们藉此共融并彼此结合。因为我们共享一个饼，都成为基督的一个身体和同一血，彼此互为肢体，与基督同为一体。\par

因此，我们要尽力谨慎，免得从异端者手中领受共融或给予他们共融；主说：\BibleQuote{不要把圣物给狗，也不要把你们的珍珠投在猪前}\BibleRef{玛 7:6}，免得我们分担他们的羞辱和定罪。因为若结合真实地与基督及彼此结合，我们肯定也自愿地与所有与我们共分的人结合。因为这结合是自愿成就的，并非违逆我们的倾向。正如神圣宗徒所说：\BibleQuote{因为我们众人共享一个饼，都是一个身体}\BibleRef{格前 10:17}。\par

此外，也称未来之事的预像，并非因为它们实质上不是基督的身体和血，而是如今我们藉着它们分受基督的天主性，而那时我们将单凭默观以心神分受。\par

% structure: p0067 source=h2 target=subsection reason=relative-heading-rank; base=h2

\subsection{第十四章 论我们主的族谱及论神圣天主之母}

关于神圣、备受赞颂的、永远童贞的天主之母玛利亚，我们在前几章中已有所论述，提出了最合适的要点，即严格而真实地，她是被称为天主之母。现在我们要补足空白。因为她被天主永恒预知的上智预定，藉圣神在先知们的各种预像和言论中被预示和宣告，在预定的时期从达味的根苗中生出，正如向他所作过的诺言。\BibleQuote{因为上主曾发誓，他向达味所起的誓是真实的，必不收回：你的后裔，我要立在你王位上。}又\BibleQuote{我曾经指着我的圣洁起誓，决不欺骗达味：他的后裔必永远存在，他的御座在我面前如同太阳；如同月亮永远确立，在天上作忠实的见证。}依撒意亚说：\BibleQuote{从叶瑟的树干要长出一根嫩枝，从他的根上要生出一个芽}\BibleRef{依 11:1}。\par

但若瑟出自达味支派，是由至圣的福音作者玛窦和路加明确证实的。但玛窦是从达味经由撒罗满追溯若瑟，而路加则经由纳堂；然而两位对圣童贞的起源都默不作声。\par

应当记得，希伯来人以及神圣圣经都没有记载女子族谱的习惯；律法也禁止一个支派从另一支派娶妻。因此，既然若瑟出自达味支派，且是一个义人（神圣福音如此证明），他不会违反律法迎娶圣童贞；除非她属同一支派，他绝不会娶她。因此，证明若瑟的出身就足够了。\par

还应当注意这一点，律法规定，若一个人没有子嗣而死，他的兄弟要娶死者之妻，为死者立嗣\BibleRef{申 25:5}。因此，所生的后代按本性属于第二个，即生他的人，但按法律属于死者。\par

生于达味的儿子纳堂一系：肋未生默耳基和潘特尔；潘特尔生巴儿潘特尔，如此称呼。这巴儿潘特尔生约雅金；约雅金生天主之母。在达味的儿子撒罗满一系中：玛堂娶妻，生雅各伯。玛堂死后，纳堂支派肋未的儿子默耳基——潘特尔的兄弟——娶了玛堂的妻，即雅各伯的母亲，生赫里。因此雅各伯和赫里成了同母异父的兄弟，雅各伯属于撒罗满支派，赫里属于纳堂支派。然后，纳堂支派的赫里无子而死，他的兄弟、撒罗满支派的雅各伯娶了他的妻子，为兄弟立嗣，生若瑟。因此，若瑟按本性是雅各伯的儿子，属撒罗满支派，但按法律是赫里的儿子，属纳堂支派。\par

若亚敬遂娶了可敬而可赞的亚纳为妻。但正如先前不孕的亚纳\BibleRef{撒上 1:2}藉祈祷和恩许生了撒慕尔，同样，这位亚纳也藉恳求和天主的恩许生了天主之母，好使她在这一点上不逊于著名的妇人。因此，恩宠（因为亚纳即此意）生了主母（因为她在成为造物主之母时，确实成了所有受造物的主母）。此外，若亚敬在称为\Emph{\OriginalTerm{羊门}{Probatica}}的房屋中诞生，并被带到圣殿养育。她犹如一棵结果的橄榄树，栽种在天主的殿中，藉圣神而成长，成为一切德行的寓所，使自己的心意远离一切世俗和肉欲的贪求，从而保持灵魂和身体都童贞，正适合她怀抱天主：因为祂是圣的，在圣者中安息。因此，她奋力追求圣德，并被称为适合至高天主的圣洁而奇妙的圣殿。\par

此外，由于我们救恩的仇敌一直注视着童贞女，正如依撒意亚的预言所说：\BibleQuote{看，一位贞女将怀孕生子，给祂起名叫\OriginalTerm{厄玛奴耳}{Emmanuel}，翻译出来就是\InnerQuote{天主与我们同在}}，为叫那以自己的智慧自夸的人被\BibleQuote{祂以巧计捕捉聪明人}所欺骗，司祭们便将童女许配给若瑟，如同将一本新书交给精通文字的人\BibleRef{依 29:11}；但这婚姻既是童贞女的保护，也是那注视童贞女者的迷惑。及至时期一满，主的使者被派往她那里，报告我们主成胎的喜讯。于是她怀了圣子——圣父的本体能力，\BibleQuote{不是由血气，也不是由肉欲，也不是由男欲}\BibleRef{若 1:13}，即并非藉交合与种子，而是由圣父的美意和圣神的合作。她服事造物主，使祂被造；服事塑造者，使祂被塑造；服事天主子和天主，使祂由她纯洁无玷的血肉成为肉躯并成为人，满足了第一位母亲的罪债。因为正如前者是由亚当无交合而成，同样，后者也生了新亚当，祂是按生产规律、超越生育本性而诞生的。\par

因为那由圣父所生而无母的，是由妇人所生而无父的合作。就祂由妇人所生而言，祂的诞生是符合生产规律的；就祂无父而言，祂的诞生超越生育本性；就祂在通常时间（因为祂在第九个月完满、第十个月开始时诞生）而言，祂的诞生符合生产规律；就祂无痛而言，祂超越生育规律。因为，正如快乐没有先于它，疼痛也没有随之而来，如先知所说：\BibleQuote{她未经历产痛就生了，疼痛未临到她身上，她就生了男儿}\BibleRef{依 66:7}。因此，由她所生的，是天主圣子降生成人，并非一位受天主感动的人，而是降生成人的天主；并非一位受傅油赋予能力的先知，而是由傅油者完全临在，以致傅油者成为人，受傅油者成为天主，不是藉本性的变化，而是藉位格的合一。因为傅油者和受傅油者是同一位，祂以天主身分自傅为人。那么，岂能没有一位生了降生成人的天主的天主之母呢？确实，那扮演造物主婢女和母亲角色的她，严格而真实地是天主之母、主母和一切受造物的元后。但正如那成胎者使她成胎的仍然保持童贞，同样，那诞生者也使她童贞完整无损，只穿过她而保持关闭\BibleRef{则 44:2}。成胎确实是通过听觉，而诞生则是通过通常的孩子来到的途径，虽然有些人虚构说祂是通过天主之母的肋旁诞生的。因为祂由这门而来，不损坏她的封印，并非不可能。\par

永远童贞者甚至在出生后仍然保持童贞，直到去世从未与男人交往。因为虽经记载：\BibleQuote{他没有认识她，直到她生了头胎儿子}\BibleRef{玛 1:25}，但请注意，那首生者就是头胎子，即使祂是独生子。因为\Quote{头胎子}一词意指祂是首先诞生的，但绝不暗示还有其他子女。而\Quote{直到}一词则指明指定时间的限度，但并不排除此后。因为主说：\BibleQuote{看，我天天与你们同在，直到今世的终结}，祂并非意指在此世代完成后祂将与我们分离。确实，神圣的宗徒说：\BibleQuote{这样，我们就要时常与主同在}\BibleRef{得前 4:17}，意指在普遍复活之后。\par

因为那怀了天主并因后来事件的经验而知道奇迹的她，怎么可能接受男人的拥抱呢？绝对不可！这样的念头不是一个纯洁的心智所能想，更不用说做这样的事了。\par

但这蒙福的妇人，被赐予超越本性的恩赐，在受难时刻遭受了她在生产时所免去的痛苦，因母性的同情而肠断，当她看见她由生育方式确认为天主的那位被当作凶犯处死时，她的心被利剑刺透，这正是这句经文的含义：\BibleQuote{至于你，要有一把利剑刺透你的心灵}\BibleRef{路 2:35}。但复活的喜乐转变了痛苦，宣告那在肉身中死去的天主。\par

% structure: p0079 source=h2 target=subsection reason=relative-heading-rank; base=h2

\subsection{第十五章 论对圣人的尊敬及其遗骸}

应当尊敬圣人，视之为基督的朋友、天主的子女和继承人，正如神学家兼圣史\Person{若望}所言：\BibleQuote{凡接受祂的，祂赐给他们成为天主子女的权能。}\BibleRef{若 1:12} \Quote{所以他们不再是仆人，而是子女；既然是子女，便是继承人，天主的继承人，与基督同为继承人}：主在圣福音中对祂的宗徒说：\BibleQuote{你们是我的朋友。}\BibleRef{若 15:14} \Quote{从此我不再称你们为仆人，因为仆人不知道主人所做的事。}此外，如果万物的创造者和主宰也被称为万王之王、万主之主\BibleRef{默 19:16}和万神之神，那么圣人也必然是神、主和王。因为天主是这些人的天主，也被称为他们的天主、主和王。祂对\Person{梅瑟}说：\BibleQuote{我是亚巴郎的天主，依撒格的天主，雅各伯的天主。}\BibleRef{出 3:6}天主并使\Person{梅瑟}成为法郎的神。我所说的神、王、主，并非指本性上的，而是指他们作为自身情欲的统治者和主宰，并保持按天主肖像所造的真理相似（因为王的肖像也称为王），且因自己的自由意志与天主结合，接纳祂居住在内心，并因恩宠藉分享而成为祂本性所是。那么，敬拜者、朋友和天主的子女难道不该受到尊敬吗？因为对最有思想的同仆表示尊敬，是对共同主人怀有好感的明证。\par

这些圣人是天主的宝库和纯洁的住所：\BibleQuote{我要住在他们中间，}天主说，\BibleQuote{在他们中间行走，我要作他们的天主。}圣经同样说，义人的灵魂在天主手中\BibleRef{智 3:1}，死亡不能抓住他们。因为死亡对圣人而言更像是睡眠，而非死亡。\Quote{他们在这世上劳苦，直到终点，}而且\Quote{在主眼中，圣人之死是宝贵的。}那么，有什么比在天主手中更宝贵的呢？因为天主是生命和光明，凡在天主手中的，就在生命和光明中。\par

此外，天主甚至以属灵方式居住在圣人的身体内，宗徒告诉我们说：\BibleQuote{你们不知道你们的身体是住在你们内的圣神的宫殿吗？}\BibleRef{格前 3:16} \BibleQuote{主就是那神}\BibleRef{格后 3:17}，\BibleQuote{若有人毁坏天主的宫殿，天主必要毁坏那人。}\BibleRef{格前 3:17}那么，我们当然应该尊敬天主的活宫殿、天主的活会幕。这些圣人活着时，曾坦然无惧地站在天主面前。\par

主基督使圣人的遗骸成为我们的救恩之泉，涌流多种祝福，并充满芬芳的油膏；任何人都不应怀疑这一点。因为如果水在旷野中从陡峭坚硬的岩石中涌出\BibleRef{出 17:6}，并从未成形的驴腮骨中涌出以解\Person{三松}之渴\BibleRef{民 15:17}，那么从殉道者的遗骸中涌出芬芳的油膏，难道不可信吗？绝不会，至少对于那些知道天主的权能和祂赐予圣人的光荣的人来说。\par

在法律上，凡触摸尸体的，就被视为不洁\BibleRef{户 19:11}，但这些圣人并非死人。因为从那本身就是生命和生命之源的主被列于死者中之时起，我们就不称那些怀着复活希望、信德而安眠的人为死人。否则，尸体怎能行奇事？怎能藉着他们驱逐魔鬼、消除疾病、治愈病人、使盲者复明、洁净癞病人、克服诱惑和患难？又怎能从光明之父\BibleRef{雅 1:17}那里，藉着他们，将一切美好的恩赐赐给那些以坚定信心祈祷的人？你岂不愿费尽心力寻求一个庇护者，好让他把你引荐给一位有死的君王，并替你说情？那么，那些作为全人类的庇护者、为我们向天主转求的人，难道不更应受尊敬吗？是的，实在应该尊敬他们，为他们以天主之名建立圣殿、献上初果、纪念他们、并在他们内获取属灵的喜乐，以便那些呼求我们的人所拥有的喜乐也成为我们的喜乐，免得我们在敬礼时反而触怒他们。因为敬礼天主的人，必因那些使天主受敬礼的事而喜悦；而天主的护卫者必因那些使天主发怒的事而愤怒。让我们信徒以圣咏、圣诗和属灵的歌曲\BibleRef{弗 5:19}、以痛悔和怜恤穷人，来敬礼圣人，因为天主也正是在这样的事上最受敬礼。让我们为他们建立纪念碑和可见的像，又让我们自己借着效法他们的德能，成为他们活的纪念碑和像。让我们尊敬那生养天主者，视她为严格而真实的天主之母。也让我们尊敬先知若翰，作为前驱和施洗者，作为宗徒和殉道者，因为主说：\Emph{\Quote{凡妇人所生的，没有一个比洗者若翰更大}}\BibleRef{玛 11:11}，他成了最先宣告天国的人。让我们尊敬宗徒们，作为主的弟兄，他们曾面对面看见主并为祂的苦难服务，\Emph{\Quote{天主所预知的人，也预定他们与祂儿子的肖像相同}}\BibleRef{罗 8:29}，\Emph{\Quote{第一是宗徒，第二是先知}}\BibleRef{格前 12:24}，\Emph{\Quote{第三是牧者和教师}}\BibleRef{弗 4:11}。也让我们尊敬从各阶层中选出的主的殉道者，作为基督的士兵，他们饮了祂的杯，并以祂带来生命的死亡之洗礼受了洗，成为祂苦难和光荣的分享者：为首的是斯德望，是基督的首位执事、宗徒和首位殉道者。也让我们尊敬我们圣洁的教父，那些拥有天主的苦修者，他们的奋斗是更长久、更艰辛的良心的战斗：\Emph{\Quote{他们披着绵羊皮、山羊皮，困苦、受折磨、遭患难；在旷野、山岭、山洞和地穴中飘流无定，这世界原配不上他们}}\BibleRef{希 11:37}。让我们尊敬那些在恩宠之前的先知、族长和义人，他们预知了主的来临。让我们细心审视这些人的生活，效法他们的信德、爱德、望德、热忱、生活方式，以及忍受苦难甚至流血的耐心，好使我们与他们一同分享荣耀的冠冕。\par

% structure: p0085 source=h2 target=subsection reason=relative-heading-rank; base=h2

\subsection{第十六章 论图像。}

但既然有人指责我们敬礼和尊敬我们的救主及其圣母的圣像，以及其余圣徒和基督仆人的圣像，就让他们记住：起初，天主按自己的肖像造了人\BibleRef{创 1:26}。我们彼此致敬，若不是因为我们按天主的肖像被造，又是基于什么理由呢？正如那位深谙神圣事物的巴西略所说，对图像的敬意归于原型。而原型是被成像者，由此衍生出肖像。为什么梅瑟的子民普遍尊敬那承载天上事物影像和样式，甚至整个受造界样式的会幕\BibleRef{出 33:10}？天主确实对梅瑟说：\Emph{\Quote{你要谨慎，按照在山上指示你的样式去做}}。那遮掩约柜的革鲁宾，不也是人手所造的吗\BibleRef{出 25:18}？再者，那著名的耶路撒冷圣殿，难道不是人手所作、凭借人的技巧建造的吗\EditorialNote{列王纪上第八章}？\par

再者，神圣的圣经责备那些崇拜雕刻偶像的人，也责备那些向魔鬼献祭的人。希腊人献祭，犹太人也献祭：但希腊人是向魔鬼献祭，而犹太人是向天主献祭。希腊人的祭品被拒绝和谴责，义人的祭品却极其蒙天主悦纳。因为诺厄献了祭，\Emph{\Quote{上主闻到了馨香}}\BibleRef{创 8:21}，接受了正确抉择和对祂美善意愿的芬芳。同样，希腊人的雕刻偶像，因为是神祇的像，所以被拒绝和禁止。\par

然而，除了这些，谁能摹仿那不可见、无形体、无界限、无形态的天主呢？因此，赋予神性以形态，乃是极致的愚昧和亵渎。故此，在旧约中，图像的使用并不普遍。但是，当天主在祂怜悯的怀抱中，为了我们的救恩而真人地成为人——并非像祂在\Person{亚巴郎}面前以人的形像显现，也非像祂藉先知所见——而是真正地成为人，且生活在地上，居于人间\BibleRef{巴 3:38}，行了奇迹，受苦，被钉十字架，复活，并被接升天之后，因这一切确曾发生且为人所见，便记录下来，以纪念和训诲我们这些当时不在世的人，好使我们虽未目睹，却能因听闻而相信，获得主的祝福。然而，并非人人识字，也非人人有时间阅读，因此教父们准许将这些事迹描绘在图像上，作为英勇行为的简明纪念。的确，我们常常在未思想主的苦难时，一看见基督被钉十字架的图像，祂救世的苦难便被记起，于是我们俯伏朝拜，不是朝拜物质，而是朝拜所描绘的对象：正如我们不朝拜制作福音书的材料，也不朝拜十字架的材料，而是朝拜它们所象征的实体。因为那象征主的十字架与不象征的十字架有何区别呢？同样，对于主之母也是如此。因为我们对她的尊敬，归于那由她降生成人的那一位。同样，圣人们的英勇事迹也激励我们勇敢，并效法和模仿他们的勇毅，光荣天主。因为正如我们所说，对最卓越的同仆所表示的尊敬，是对我们共同主母的善意的证明，且对图像的尊敬归于原型。然而，这是一项\OriginalTerm{未成文传统}{unwritten tradition}，正如朝向东方的敬礼、对十字架的敬礼，以及其他许多类似之事。\par

还有一个故事，讲述当\Person{阿布加尔}作\Place{厄德萨城}王时，他派了一位肖像画家去画主的肖像。由于画家因主面容所发之光而不能绘画，主便将一块布覆在祂神圣而赋予生命的面容上，将其容貌印在上面，并将此物送给\Person{阿布加尔}，以满足他的愿望。\par

此外，宗徒们传下了许多未成文的传统，外邦人的宗徒\Person{保禄}用以下的话告诉我们：\BibleQuote{所以，弟兄们，你们要站立稳固，持守我们所传授给你们的传统，无论是藉着言语，还是藉着书信。}\BibleRef{得后 2:15} 他又写信给格林多人说：\BibleQuote{我称赞你们，弟兄们，因为你们凡事记念我，并持守我传授给你们的传统。}\BibleRef{格前 11:2}\par

% structure: p0091 source=h2 target=subsection reason=relative-heading-rank; base=h2

\subsection{第十七章 论圣经}

旧约和新约所宣告的是同一位天主，祂在圣三中受赞美并得荣耀：\BibleQuote{我来不是为废除法律，而是为成全。}\BibleRef{玛 5:17} 因为祂亲自完成了我们的救恩，整个圣经和一切奥秘皆为此而存在。又说：\BibleQuote{你们查考圣经，因为你们以为其中有永生，正是这些为我作证。}\BibleRef{若 5:39} 宗徒也说：\BibleQuote{天主在古时，曾多次并以多种方式，藉着先知向我们的祖先说过话；在这末后的日子里，祂藉着祂的儿子向我们说话。}\BibleRef{希 1:1} 因此，法律和先知、福音作者和宗徒、牧者和教师，都是藉着圣神说话的。\par

全部圣经都是\BibleQuote{天主所默示的，并且确实是有益的。}\BibleRef{弟后 3:16} 因此，查考圣经是极美好且极有益于灵魂的工作。正如栽种在水渠旁的树，为神圣圣经所灌溉的灵魂也是如此，它按时结果实，即正统的信仰，并饰以常青的叶子，即中悦天主的作为。因为藉着圣经，我们受训导以行中悦天主的事，并获致无扰的默观。因为我们在其中找到对每一德行的劝勉和对一切恶行的劝阻。因此，如果我们热爱学习，我们也将博学多闻。因为藉着勤勉、辛劳和赐予者天主的恩宠，万事得以成就。\BibleQuote{因为凡求的，就得到；找的，就找到；敲的，就给他开门。}\BibleRef{路 11:10} 因此，让我们敲击那极美的圣经花园吧，它芳香甜美、繁花盛开，周围回响着属灵且由天主所感的鸟儿的各种歌声，抓住我们的心，安慰忧伤者，平息愤怒者，并充满永久的喜乐：它将我们的心思置于那神圣鸽子的金光灿烂的背上，其明亮的翅膀将我们带到那属灵葡萄园的农夫\BibleRef{玛 21:37}的独生子及继承人那里，并藉着祂将我们带到众光之父那里。\BibleRef{雅 1:17} 但让我们不要漫不经心地敲，而要热切且恒常地敲，免得我们因敲而疲倦。因为这样，门将为我们开启。如果我们读了一次或两次，却不明白所读的，让我们不要疲倦，而要持之以恒，多谈论，多询问。因为经上说：\BibleQuote{你要问你的父亲，}他说，\BibleQuote{他必指示你；问你的长老，他们必告诉你。}\BibleRef{申 32:7} 因为\BibleQuote{知识并非人人都有的。}\BibleRef{格前 8:7} 让我们从那花园的泉源中汲取涌流到永生的纯净井水。\BibleRef{若 4:14} 让我们在此尽情享受，不知餍足：因为圣经拥有无尽的恩宠。倘若我们能从外在来源采撷任何有益之物，也未尝不可。让我们成为经考验的兑换银钱者，积攒真实而纯正的金子，丢弃假冒的。让我们保存最美好的言辞，却将荒谬的神祇和怪异的神话丢给狗：因为藉着它们本身，我们可以强力地胜过它们。\par

此外，请注意，旧约共有二十二卷，与希伯来语的字母数相同。因为希伯来语有二十二个字母，其中有五个是重复的，因此共有二十七个。这些字母是：Caph、Mem、Nun、Pe、Sade。因此，按此方式书卷数为二十二，但因五个字母的双重性而成为二十七。因为《卢德传》附于《民长纪》，希伯来人将其算作一卷；\WorkTitle{列王纪}上、下算作一卷；\WorkTitle{列王纪}三、四算作一卷；\WorkTitle{编年记}上、下算作一卷；\WorkTitle{厄斯德拉}上、下算作一卷。这样，这些书卷被分为四部\WorkTitle{五书}和两部剩余，从而构成正典书目。其中五卷属于法律书，即\WorkTitle{创世纪}、\WorkTitle{出谷纪}、\WorkTitle{肋未纪}、\WorkTitle{户籍纪}、\WorkTitle{申命纪}。这法律的汇编构成第一部\WorkTitle{五书}。其次是第二部\WorkTitle{五书}，即所谓\OriginalTerm{圣卷}{Grapheia}，或有些人所称的\OriginalTerm{圣文集}{Hagiographa}，包括以下书卷：\WorkTitle{若苏厄书}、\WorkTitle{民长纪}附\WorkTitle{卢德传}、\WorkTitle{列王纪}上、下（合为一卷）、\WorkTitle{列王纪}三、四（合为一卷）、\WorkTitle{编年记}上、下（合为一卷）。这是第二部\WorkTitle{五书}。第三部\WorkTitle{五书}是诗歌书，即\WorkTitle{约伯传}、\WorkTitle{圣咏集}、\WorkTitle{箴言}、\WorkTitle{训道篇}、\WorkTitle{雅歌}。第四部\WorkTitle{五书}是先知书，即十二位小先知（合为一卷）、\WorkTitle{依撒意亚}、\WorkTitle{耶肋米亚}、\WorkTitle{厄则克耳}、\WorkTitle{达尼尔}。然后是被合为一卷的\WorkTitle{厄斯德拉}上、下和\WorkTitle{艾斯德尔传}。还有\OriginalTerm{全德之书}{Panaretus}，即\WorkTitle{智慧篇}，以及息辣之子耶稣的\WorkTitle{智慧篇}（即\WorkTitle{德训篇}），原文由息辣之父以希伯来文出版，后由其孙息辣之子耶稣译成希腊文。这些书卷是良善高贵的，但未被计数，也未放在约柜中。\par

新约包含四部福音：\WorkTitle{玛窦福音}、\WorkTitle{马尔谷福音}、\WorkTitle{路加福音}、\WorkTitle{若望福音}；圣史路加所著的\WorkTitle{宗徒大事录}；七封公函：雅各伯书一封、伯多禄书两封、若望书三封、犹达书一封；保禄宗徒的十四封书信；若望圣史的\WorkTitle{默示录}；由克莱孟所收集的\OriginalTerm{宗徒训条}{Canons}。\par

% structure: p0096 source=h2 target=subsection reason=relative-heading-rank; base=h2

\subsection{第十八章 论及关于基督所说的一切}

关于基督所说的话分为四种主要模式。有些话甚至适用于祂降生成人之前，有些适用于结合之中，有些适用于结合之后，有些适用于复活之后。再者，涉及降生成人之前的话语有六种模式：有些阐明性体的合一及与圣父本质的同一，例如\Quote{我与父原是一体}\BibleRef{若 10:30}，以及\Quote{看见我的，就是看见了父}，还有\Quote{祂本具有天主的形体}\BibleRef{斐 2:6}等。另一些阐明\OriginalTerm{自立}{subsistence}的圆满，例如\Quote{天主子}、\Quote{祂是神体之真像}\BibleRef{希 1:3}，以及\Quote{伟大的谋士、奇妙的谋士}\BibleRef{依 9:6}等。\par

再者，另一些阐明各\OriginalTerm{自立}{subsistence}\OriginalTerm{相互内住}{indwelling}，例如\Quote{我在父内，父在我内}\BibleRef{若 14:10}；以及不可分离的根基，例如圣言、智慧、能力、\OriginalTerm{光辉}{Effulgence}。因为圣言不可分离地存在于理性之内（我指的是本质的理性），同样，智慧和能力存在于有能力者之内，光辉存在于光明之内，皆由此而生。\par

还有一些话语显明祂源于圣父作为原因（\OriginalTerm{原因意义上的}{causal sense}），例如\Quote{父比我大}\BibleRef{若 14:28}。因为祂从父获得存在及所有一切：祂的存在是藉生育而非受造，如\Quote{我从父出来，来到了世界上}\BibleRef{若 16:28}，以及\Quote{我因父而生活}。但祂所有的一切并非藉恩赐或教导，而是在原因意义上，如\Quote{子不能由自己做什么，只有看见父做什么，才能做什么}。因为若没有父，子也不存在。因为子出于父，在父内，与父同在，并非在父之后。同样，祂所做的一切也是出于父并与父同在。因为在父、子、圣神中，有一个同一的（并非相似而是同一的）意愿、能力和行动。\par

此外，另一些话则表明父的善意藉其能力而实现，并非藉着工具或仆人，而是藉其本质的、\OriginalTerm{位格的}{hypostatic}圣言、智慧和能力，因为在父和子中只观察到同一行动，例如\Quote{万物是藉着祂造成的}\BibleRef{若 1:3}，以及\Quote{祂派遣了自己的话，治好了他们}，还有\Quote{为叫他们相信是你派遣了我}\BibleRef{若 17:21}。\par

另一些话具有预言意义，其中有些是未来时态：例如\Quote{祂要公开来临}，以及匝加利亚的\Quote{看，你的君王到你这里来}\BibleRef{匝 9:9}，还有米该亚的\Quote{看，上主离开自己的地方，降来踏地的高处}\BibleRef{米 1:3}。但另一些虽属未来，却以过去时态表达，例如\Quote{这是我们的天主：为此，祂显现于地上，与人同住}\BibleRef{巴 3:38}，以及\Quote{上主在造化之初创造了祂，为祂的工程}\BibleRef{箴 8:22}，还有\Quote{因此天主，你的天主，以喜乐之油傅你，胜过你的同伴}等。\par

那么，那些指称合一之前的事，即使在合一之后也适用于祂；但那些指称合一之后的事，在合一之前则完全不适用，除非如我们所说，是先知性的说法。那些指称合一时期的事有三种方式。因为当我们的论述涉及较高层面时，我们谈论肉身的被神化、圣言的被摄取和超越性的高举，等等，显明肉身从与至高的天主圣言的合一及本性结合中所获得的富饶。而当我们的论述涉及较低层面时，我们谈论天主圣言的降生成人、祂成为人、祂的自我空虚、祂的贫穷、祂的谦卑。因为这些事情和类似的事是通过祂与人性结合而归于圣言和天主的。当我们再次同时兼顾两个方面时，我们谈论合一、相通、傅油、本性结合、形态相似等。前两种方式其理由在于这第三种方式。因为通过合一，我们明了每一方从与另一方的亲密结合与渗透中得到了什么。因为通过位格合一，肉身被说成是被神化的、成为天主、与圣言同等为天主\BibleRef{依 48:12}：这不是说两种本性被转化为一个复合的本性（因为对立的本性特质不可能同时存在于一个本性之中），而是说两种本性在位格上合一，并互相渗透，毫无混淆或转变。而且，渗透并非来自肉身，而是来自神性：因为肉身不可能渗透神性；但神性本性一旦渗透肉身，也赐予肉身同样不可言喻的渗透能力；而这正是我们所谓的合一。\par

也要注意，在属于合一时期的第一种和第二种方式中，可以观察到互易性。因为当我们谈论肉身时，我们使用被神化、圣言的被摄取、超越性的高举和傅油这些术语。因为这些源自神性，但却是相对于肉身而被观察到的。而当我们谈论圣言时，我们使用自我空虚、降生成人、成为人、谦卑等术语：这些，如我们所说，是通过肉身而归于圣言和天主的。因为祂亲自甘愿承受了这些事。\par

在指称合一之后的事中有三种方式。第一种声明祂的神性本性，例如：\Quote{我在父内，父也在我内}\BibleRef{若 14:1}，以及\Quote{我与父原是一体}；以及所有那些在祂摄取人性之前肯定于祂的事，这些在祂摄取人性之后也肯定于祂，除了祂并没有摄取肉身及其本性属性这一点。\par

第二种声明祂的人性本性，例如：\Quote{如今你们却想杀死我，一个告诉你们真理的人}，以及\Quote{人子也必须照样被举起来}，等等。\par

此外，关于基督救主按人情方式所说所写的话，无论是涉及言语还是行动，有六种方式。因为其中有些是按照降生成人自然地做成或说出的；例如，祂由童贞女诞生、祂在年龄上的成长与进步、祂的饥渴、疲倦、恐惧、睡眠、被钉、死亡以及所有这一类自然而无罪的苦难。因为在这其中，都有神性与人性的混合，尽管它们实际上被认为属于身体，神性并未经受其中任何一样，而是通过它们为我们谋取了救恩。\par

另一些则属于归因的性质，例如基督的问题：\Quote{你们把拉匝禄放在哪里？}\BibleRef{若 11:34}，祂跑向无花果树，祂的退避，即祂的后退，祂的祈祷，以及祂\Quote{装作还要往前行}\BibleRef{路 24:28}。因为无论作为天主还是作为人，祂都不需要这些或类似的事，只是因为祂的形态是人的形态，符合需要和适宜。例如，祈祷是为了显示祂并非与天主对立，因为祂尊敬父，视父为祂自己的原因；提问并非出于无知，而是为了显示祂既是天主也是真人；后退是为了教导我们不要急躁，不要放弃自己。\par

还有一些是以联合与关系的方式说的，例如：\Quote{我的天主，我的天主，祢为什么舍弃了我？}\BibleRef{玛 27:46}，以及\Quote{祂曾使那不知罪的，替我们成了罪}\BibleRef{格后 5:21}，\Quote{为我们成了咒骂}\BibleRef{迦 3:13}；还有，\Quote{那时，子也要自己屈服于那使万有屈服于祂的}\BibleRef{格前 15:28}。因为无论作为天主还是作为人，祂从未被父舍弃，也没有成为罪或咒骂，也不需要屈服于父。因为作为天主，祂与父同等，既不与父对立，也不屈服于父；作为天主，祂从未在任何时候不顺服于祂的生养者，以至于需要使祂屈服。因此，祂借用我们的人格，将自己与我们等同，说了这些话。因为我们因不信与不顺从而被罪的锁链和咒骂所束缚，因而被舍弃。\par

另一些则是因思想上的区分而说的。因为如果你在实际上不可分的事物中于思想上加以区分，将肉身从圣言中割裂，那么\OriginalTerm{仆人}{servant}和\OriginalTerm{无知}{ignorant}这些词就用在祂身上，因为祂确实有一个服从的和无知的本性，除非与天主圣言位格合一，祂的肉身本是奴仆性的和无知的。但由于与天主圣言的位格合一，它既非奴仆性也非无知。同样，祂也称父为祂的天主。\par

\Quote{另有其他句子旨在向我们启示祂并坚固我们的信德，就如：} \BibleQuote{父啊，现在求祢以我与祢在未有世界以前所同享的荣耀，荣耀我。}\BibleRef{若 17:5} \Quote{以及宗徒所说的：} \BibleQuote{按圣德的神，因从死者中复活，被宣示为具有大能的天主子。}\BibleRef{罗 1:4} \Quote{因为借着奇迹、复活和圣神的降临，祂是天主子的事实向世界显明并得到确认。此外：} \BibleQuote{孩子渐渐长大，充满智慧和恩宠。}\BibleRef{路 2:40}\par

\Quote{另一些则涉及祂认同犹太人的个人生活，将自己列在犹太人中间，如祂对撒玛黎雅妇人所说：} \BibleQuote{你们朝拜你们所不认识的；我们朝拜我们所认识的，因为救恩出自犹太人。}\BibleRef{若 4:22}\par

\Quote{第三种方式是宣示一个位格并显明双重本性。例如：} \Quote{我因父而生活；那吃我的，也要因我而生活。} \Quote{以及：} \Quote{我往父那里去，你们不再见到我了。} \Quote{以及：} \BibleQuote{他们不会把光荣的主钉在十字架上。}\BibleRef{格前 2:8} \Quote{以及：} \BibleQuote{没有人上过天，除了那从天降下、仍在天上的人子。}\BibleRef{若 3:13}\par

\Quote{关于复活后时期的肯定陈述，有一些适合于天主，如：} \BibleQuote{因父、及子、及圣神之名给他们授洗。}\BibleRef{玛 28:19} \Quote{因为这里\InnerQuote{子}明显被用作天主；还有：} \Quote{看！我同你们天天在一起，直到今世的终结。} \Quote{等等。因为祂作为天主与我们同在。另一些适合于人的，如：} \Quote{他们抱住祂的脚} \Quote{以及} \Quote{他们在那里要看见我} \Quote{等等。}\par

\Quote{此外，关于复活后时期适合人的陈述有不同的方式。有些确实实际发生，但并非按本性，而是按救恩计划，为证实那受苦的身体确实复活了；例如伤痕、复活后的吃喝。另一些实际且自然地发生，如不费力地从此地移到彼地，穿过关闭的门户。另一些带有假装的性质，例如} \BibleQuote{祂装作还要往前行。}\BibleRef{路 24:28} \Quote{另一些适合双重本性，如：} \BibleQuote{我升到我的父和你们的父那里，我的天主和你们的天主那里。}\BibleRef{若 20:17} \Quote{以及} \Quote{光荣的王要进来。} \Quote{以及} \BibleQuote{祂坐在至高者尊威的右边。}\BibleRef{希 1:3} \Quote{最后，还有一些应理解为祂把自己放在与我们同等的地位，以纯粹思想中的分离方式，如} \BibleQuote{我的天主和你们的天主。}\BibleRef{若 20:17}\par

\Quote{因此，崇高的言辞必须归于神性的本性，它超越苦难和肉体；卑微的言辞必须归于人性；共同的言辞必须归于复合体，即唯一基督，祂既是天主又是人。并且应该理解，两者都属于同一位耶稣基督，我们的主。因为如果我们知道各自的特质，并认识到两者都由同一位所完成，我们就拥有真实的信德，不会走入歧途。从这一切中，我们认识到联合的本性之间的区别，以及正如最虔诚的\Person{济利禄}所说，它们在神性和人性的自然性质上并非同一。然而，只有一位圣子和基督和主：既然祂是唯一，祂也只有一个位格，位格的合一绝不会因认识本性的差异而分裂成部分。}\par

% structure: p0116 source=h2 target=subsection reason=relative-heading-rank; base=h2

\subsection{天主不是邪恶的原因}

\Quote{应当注意，在\WorkTitle{圣经}中，习惯将天主的许可说成是天主的运作，就如宗徒在\BibleBook{罗马}{书}中所说：} \Quote{陶匠难道没有权能，从同一团泥中，造一个器皿作贵重的，另一个作卑贱的吗？} \Quote{因此，是祂自己造了这个或那个。因为祂独自是万物的创造者；然而，不是祂自己造了贵重或卑贱之物，而是每个人自己的自由选择。这从同一位宗徒在\BibleBook{弟茂德}{后书}中所说的显明：} \BibleQuote{在大户人家，不但有金器银器，也有木器瓦器；有些是贵重的，有些是卑贱的。所以人若自洁，离开这些，他必成为贵重的器皿，圣洁的，合主人用的，预备行各种善事。}\BibleRef{弟后 2:20} \Quote{显然，这洁净必须是自愿的，因为他说：\InnerQuote{人若自洁}。随之而来的对句回应：} \Quote{人若不自洁，他必成为卑贱的器皿，不合主人使用，只配被打碎。} \Quote{因此，我们所引用的这段话以及} \BibleQuote{天主把众人都禁锢在不信之中}\BibleRef{罗 11:32} \Quote{和} \Quote{天主赐给他们沉睡的神，眼睛看不见，耳朵听不见} \Quote{，这一切都必须理解为不是天主亲自运作，而是天主允许，既因为自由意志，也因为善不会强制。}\par

\Quote{所以，祂的许可在\WorkTitle{圣经}中通常被称为祂的运作和工程。不仅如此，即使当祂说\InnerQuote{天主创造邪恶之事}，以及} \BibleQuote{城中若有大祸，不是上主所为？}\BibleRef{亚 3:6} \Quote{，祂并非通过这些言语指主是邪恶的原因，而是\InnerQuote{邪恶}一词有两种用法，两种含义。因为有时它指本性上是恶的，这相反德行和天主的旨意；有时它指对我们的感觉是恶和压迫的，即痛苦和灾难。这些看似是恶，因为使人痛苦，但实际上却是善。因为对于那些明白的人，它们成了悔改和救恩的使者。\WorkTitle{圣经}说这些恶是天主所造的。}\par

此外，也当留意，这些（恶事）的起因在我们：因为非自愿的恶是自愿的恶的产物。\par

也当承认，在圣经中，有些本应视为结果的事，常以原因的意义被陈述，如：\BibleQuote{我得罪了祢，惟独得罪了祢，在祢眼前行了这恶，以致祢在说话时显为公义，在审判时显为得胜}。因为犯罪者并非为了叫天主得胜而犯罪，天主也不需要我们的罪来藉此显为胜利者。因为祂远超一切，甚至对无罪者也赢得胜利的奖赏，祂是造物主，不可理解，非受造，拥有本然的而非外加的荣耀。但事实是：当我们犯罪时，天主在向我们发怒时并非不义；当祂赦免悔改者时，祂藉此显为战胜我们邪恶的胜利者。但我们犯罪并非为此，而是结果如此。好比一人坐着做工，一位朋友站在旁边，那人便说：我的朋友来，是为叫我那天不做工。其实朋友在场并非为叫那人不做工，而是结果如此。因为忙于接待朋友，他就没有做工。这些事也被称为结果，因为事情如此演变。况且，天主并不愿意只有祂自己是公义的，而愿意所有的人都尽可能相似祂。\par

% structure: p0121 source=h2 target=subsection reason=relative-heading-rank; base=h2

\subsection{第二十章 并没有两个国度}

并没有两个国度，一善一恶，我们从以下可以看出。因为善与恶彼此对立，互相毁灭，不能共存于彼此或与彼此共存。因此，它们各自在自己的部分中属于整体，并且首先它们不仅被整体所限，也被整体的部分所限。\par

其次我问：是谁为每一个指定其位置？他们不会声称彼此已达成友好协议或和好。因为当恶与善和平、和好时，它就不是恶；而当善与恶友好相处时，它也不是善。但如果为这两者各自划出活动范围者是与它们不同的存在，那么祂更应该是天主。\par

确实，二者必居其一：要么它们彼此接触并互相毁灭，要么存在某个中间地带，既无善也无恶，像隔墙一样将两者分开。那么就不再有两个国度，而是三个。\par

再者，以下情况也是必然的：要么它们和平相处（这与恶完全不相容，因为和平的不是恶），要么它们争斗（这与善不相容，因为争斗的不是完全善），要么恶争斗而善不报复，却被恶毁灭，要么它们永远处于烦恼和痛苦中（这并非善的标志）。因此，只有一个国度，远离一切邪恶。\par

但若如此，他们说，邪恶从何而来？因为恶绝不可能源自善。那么我们的回答是：恶无非是善的缺失，是从自然之物堕落为非自然之物：因为没有任何邪恶是自然的。因为天主所造的一切，在受造时都\BibleQuote{非常好}\BibleRef{创 1:31}。因此，如果它们保持被造时的原样，就非常好；但当它们自愿偏离自然之物而转向非自然之物时，就滑入了邪恶。\par

因此，按本性，万物都是造物主的仆役，服从祂。所以，当祂的任何受造物自愿反叛，变得不服从其造物主时，它就为自己引入了邪恶。因为邪恶不是任何本质，也不是本质的属性，而是一种\OriginalTerm{依附体}{accidens}，即从自然之物自愿偏向非自然之物，这就是罪。\par

那么，罪从何而来？它是魔鬼自由意志的发明。魔鬼是恶的吗？就他被造而言，他不是恶的，而是善的。因为他被其造物主造为明亮且非常光明的天使，作为理性者被赋予自由意志。但他自愿偏离了本然之德，进入了邪恶的黑暗，远离了天主（唯独祂是善，能赐予生命和光明）。因为一切善物都从祂获得善性，而任何事物（在意志上，而非在空间上）离祂越远，就越陷入邪恶。\par

% structure: p0129 source=h2 target=subsection reason=relative-heading-rank; base=h2

\subsection{第二十一章 天主预知那些将犯罪而不悔改的人，为何仍造他们}

天主出于美善，从无中创造了万有，并预知将来存在的一切。因此，如果它们将来不存在，它们将来既不为恶，也不会被预知。因为知识关乎存在的事物，预知关乎将来必定存在的事物。因为单纯的存在在先，然后是善或恶的存在。但如果那些因天主的良善而将在未来存在者，因他们自择将变恶而被阻止存在，那么恶就胜过了天主的良善。因此，天主使祂的一切受造皆为善，但各物凭自己的选择成为善或恶。所以，虽然主说：\BibleQuote{那人倒不如没有生下来为好}\BibleRef{谷 14:21}，祂说这话并非谴责祂的创造，而是谴责祂的受造物因自己的选择和自己疏忽而获得的恶。因为人判断的疏忽，使造物主的恩惠对他毫无益处。好比有人从君王那里获得了财富和权柄，却凌驾于恩主之上；当恩主击败他之后，若见他坚持夺权到底，便会按他应得的惩罚他。\par

% structure: p0131 source=h2 target=subsection reason=relative-heading-rank; base=h2

\subsection{第二十二章 论天主的法律和罪的法律}

天主是善的，且超越善，祂的旨意亦然。因为天主所愿的，就是善。再者，那教导这点的诫命，就是法律，使我们持守它，能在光中行走\BibleRef{若一 1:7}：违反这诫命就是罪，而罪之所以继续存在，是由于魔鬼的攻击以及我们不受约束、自愿接纳它。\BibleRef{罗 7:23}这也被称为法律\BibleRef{罗 7:25}。\par

因此，天主的法律安顿在我们的理智中，将它吸引向自身，并刺痛我们的良心。而我们的良心，也被称为我们理智的法律。再者，恶者的攻击，即罪的律法，安顿在我们肉身的肢体中，通过它向我们发动攻击。因为我们曾经一次自愿违反天主的法律并接纳恶者的攻击，我们就为它开了门，将自己卖给了罪。因此，我们的身体很容易被驱向它。所以，储存在我们身体中的罪的滋味和知觉，即身体的贪欲和快感，就是我们肉身的肢体中的律法。\par

因此，我理智的法律，即良心，同情天主的法律，即诫命，并以之为其旨意。但是罪的律法，即通过我们肢体中的律法所作的攻击，或通过身体和灵魂非理性部分的贪欲、倾向和动作，与我的理智的法律（即良心）相对抗，并因混合而俘虏我（即使我以天主的法律为我的旨意，并爱慕它，而不以罪为我的旨意）；并且通过享乐的软弱和身体及灵魂非理性部分的贪欲，如我所说，它误导我，诱使我成为罪的奴仆。但是\BibleQuote{法律若做不到的事，因它为肉身所软弱，天主就派遣自己的儿子，带着罪恶肉身的形状}（因为祂取了肉身，但没有罪）\BibleQuote{在肉身中定罪了罪，使法律的要求，成就在我们这些不随从肉身，而随从圣神的人身上}\BibleRef{罗 8:3}。因为\BibleQuote{圣神扶助我们的软弱}，并赐予能力给我们理智的法律，以对抗我们肢体中的律法。因为经上说：\BibleQuote{我们不知道我们该怎样祈求，但圣神亲自以无可言喻的叹息代我们转求}，它亲自教导我们该求什么。因此，除非藉着忍耐和祈祷，否则不可能遵行主的诫命。\par

% structure: p0135 source=h2 target=subsection reason=relative-heading-rank; base=h2

\subsection{第23章 驳犹太人论安息日问题}

第七日被称为安息日，意即安息。因为在这一天，天主歇了祂所做的一切工作\BibleRef{创 2:2}，正如神圣的圣经所说：因此日数上升到七，然后又循环回到第一日。这对犹太人而言是宝贵的数字，天主曾命令应尊崇它，且不是随意地，而是以施加最重的惩罚来规定违反者。祂这样规定并非简单随意，而是出于某些原因，被有灵性和明悟的人以奥秘的方式理解。\par

事实上，就我所知的无知而言，先从低等和较粗糙的事物说起：天主知道以色列子民的迟钝、他们肉欲的爱和凡事倾向于物质的习性，就制定了这条法律：首先，为了使\BibleQuote{仆人及牲畜得到休息}\BibleRef{申 5:14}，正如经上所记，因为\BibleQuote{义人顾惜他牲畜的性命}\BibleRef{箴 12:10}。其次，为了使他们在摆脱物质事物的纷扰后，能聚集到天主面前，用整个第七日以圣咏、圣歌和属灵歌曲，并研究神圣的圣经，安息于天主。因为在没有法律、没有天主启示的圣经时，安息日并未奉献给天主。但是当天主启示的圣经由梅瑟颁赐后，安息日就被奉献给天主，目的是使那些不将全部生命奉献给天主、不使自己的欲望臣服于主人如同臣服于父亲、反而像愚笨仆人的人，在那一天多多谈论有关操练的事，并从他们的生命中抽取一小部分——实在是微不足道的一部分——用于事奉天主，而且是出于对威胁违命者的惩戒和惩罚的恐惧。因为\BibleQuote{法律不是为义人设立的，乃是为不义的人设立的}\BibleRef{弟前 1:9}。诚然，梅瑟是第一个与天主一起禁食四十天，又再四十天，因此无疑地在安息日使自己忍饥挨饿，尽管法律禁止在安息日自苦。但如果他们反驳说这发生在法律之前，那么对于特斯彼人厄里亚，他靠一顿饭走了四十天的路程\BibleRef{列上 19:8}，他们又能说什么呢？因为他这样在安息日自苦，不仅忍饥挨饿，还走了四十天的路程，破坏了安息日；然而赐予法律的天主并没有对他发怒，反而在曷勒布山向他显现，作为对他德行的赏报。关于达尼尔，他们又能说什么呢？难道他不是三周没有进食吗\BibleRef{达 10:2}？再者，难道所有以色列人不是在安息日给孩子行割损礼，如果出生后第八天正好是安息日的话\BibleRef{创 17:12}？难道他们不遵守法律所规定的大斋，如果它碰巧落在安息日\BibleRef{肋 16:31}？再者，难道司祭和肋未人做会幕的工作时不是亵渎了安息日\BibleRef{玛 12:5}，却被认为无罪吗？是的，如果一头牛在安息日掉进坑里，把它拉上来的人是无罪的，而疏忽不救的人却被定罪。难道所有以色列人不是抬着天主的约柜环绕耶里哥城墙七天，其中当然包括了安息日\EditorialNote{若苏厄书第三章}。\par

因此，正如我所说，为了确保闲暇来敬拜天主，以便使仆人和牲畜都能将极小的一部分奉献给祂并得安息，安息日的遵守是为那些仍属血肉、尚在稚气、\Emph{受世界蒙学所束缚} \BibleRef{迦 4:3}、且无法理解超越身体和文字的事物的人而设立的。\Emph{但是时期一满，天主就派遣了自己的独生子，生于女人，生于法律之下，为赎出那些在法律之下的人，使我们获得义子的地位。}\Emph{凡接受祂的，祂赐给他们权能，成为天主的子女，即那些信祂名字的人} \BibleRef{若 1:12}。\Emph{所以我们不再是仆人，而是子女} \BibleRef{迦 4:7}：不再在法律之下，而在恩宠之下；我们不再因恐惧而部分地事奉天主，而是必须将自己整个生命奉献给祂，并使那仆人——我意思是忿怒和情欲——停止犯罪，并命令它投身于事奉天主，常将我们的一切愿望导向天主，并用我们的忿怒武装起来反对天主的敌人；同样，我们也阻止那牲畜——即身体——从事于罪的奴役，并催促它尽力协助神圣的诫命。\par

这些就是基督的属灵法律所命令我们的事，遵守这些的人就超越了梅瑟的法律。\Emph{因为那完美的来到时，那局部的就必被废除} \BibleRef{格前 13:10}；当法律的遮盖——即那幔子——因救主的被钉而裂开，圣神以火舌照耀时，文字就被废除，有形的事物就终止，奴役的法律就被成全，自由的法律就被赐予我们。是的，我们将庆祝人性的完美安息，即复活后的那一天，主耶稣——生命之作者和我们的救主——将带领我们进入应许给那些在圣神内事奉天主的人的产业，这产业祂自己作为我们的前驱在从死者中复活后进入了，并且在那里，天门为祂敞开，祂以身体的形式坐在圣父的右边，那些遵守属灵法律的人也将去到那里。\par

因此，属于我们这些按圣神而非按文字行走的人的，是完全放弃血肉之事，是属灵的敬拜和与天主的共融。因为割礼就是放弃血肉之乐以及一切多余和不必要的东西。因为包皮无非是性器上多余的皮肤。的确，一切不是来自天主、不在天主内的愉悦，对愉悦而言都是多余的；包皮正是这多余之事的预象。安息日则是停止犯罪；因此这两件事合而为一，当它们被属灵的人遵守时，两者一同并不造成任何违反法律的事。\par

此外，请注意，数字七表示整个现世时间，正如至智者撒罗满所说：\Quote{当分给七份，甚至八份} \BibleRef{训 11:2}。而达味，那位神圣的咏歌手，在创作第八篇圣咏时，歌颂了复活后未来的更新。既然法律命令第七天应从血肉之事中休息并致力于属灵之事，这对那真正的以色列人——有心看见天主的人来说——是一个神秘的指示：他应透过一切时间将自己奉献给天主，并超越血肉之事。\par

% structure: p0142 source=h2 target=subsection reason=relative-heading-rank; base=h2

\subsection{第二十四章 论童贞}

血肉之人滥用童贞，而追求享乐之人则提出以下经文作为证据：\Emph{凡不在以色列中留下子嗣的，是可诅咒的。}但我们——因那由童贞玛利亚降生成人的天主圣言而得信心——回答说，童贞自古以来就由上植入人的本性中。因为人是由处女地形成的。厄娃仅由亚当一人受造。在乐园中，童贞盛行。的确，神圣的圣经说：\Emph{亚当和厄娃二人都赤身露体，并不羞耻。} \BibleRef{创 2:23}但在他们违命之后，他们知道自己赤身露体，因羞耻而为自己缝制了围裙。当违命之后，亚当听到：\Emph{你既是灰土，你还要归于灰土。}当死亡因违命而进入世界时，\Emph{亚当认识了自己的妻子厄娃，她怀孕生了儿子。} \BibleRef{创 4:1}因此，为防止种族因死亡而消耗和毁灭，婚姻被创造出来，以便人类通过生育子女得以保存。\par

但他们也许会问：那么，\Quote{男人和女人，} \BibleRef{创 1:27}以及\Quote{你们要生育繁殖？}有什么意义呢？作为回答，我们要说：\Quote{你们要生育繁殖}并不完全指通过婚姻结合而繁殖。因为天主本可以以不同的方式繁殖人类，如果他们始终遵守命令而不违背的话。但是，天主在一切存在之前就知道一切，因祂的预知知道他们将来会陷入违命并应受死亡，祂预先就创造了\Quote{男人和女人，}并命令他们\Quote{你们要生育繁殖。}让我们继续前行，看看童贞的光荣：这也包括贞洁。\par

\Person{诺厄}奉命进入方舟，受托保存世界的种子时，主对他说：\Quote{进去吧}，主说，\Quote{你和你的儿子、你的妻子、你儿子的妻子}。他使他们与妻子分开，为叫他们能纯洁地逃脱洪水及全世界的沉没。然而洪水过后，主又说：\BibleQuote{走出方舟，你和你的儿子、你的妻子、你儿子的妻子}。\BibleRef{创 8:16} 看，婚姻再次被赐予，为使人繁衍。其次，\Person{厄里亚}，这位喷火的车夫和天上的暂居者，并未选择独身，但难道他的美德不是由他超人的升天所证实的吗？\BibleRef{列下 2:11} 谁关闭了天空？谁复活了死者？谁分开了\Place{约旦河}？难道不是守贞的\Person{厄里亚}吗？他的门徒\Person{厄里叟}，在证明同等的美德后，岂不祈求并继承了圣神恩宠的双倍福分吗？那三个青年如何？他们岂非因实践贞洁而变得比火更强，他们的身体因贞洁而能抵抗火焰吗？\BibleRef{达 3:20} \Person{达尼尔}的身体岂非因贞洁而如此坚硬，以致野兽的牙齿无法咬入？天主在愿以色列人看见祂时，岂不命令他们洁净身体？司铎们岂不洁净自己，然后才走近圣殿的至圣所并奉献祭品？法律岂不称贞洁为伟大的许愿？\par

因此，法律的诫命要在更属神的意义上理解。因为有一种属神的种子，是通过对天主的爱和敬畏在属神的子宫中怀胎，产生救恩的精神。在这个意义上，必须理解这句经文：\Quote{那在\Place{熙雍}有种子、在\Place{耶路撒冷}有后代的人是有福的。}因为这句话的意思是，即使他是一个淫乱者、一个酒鬼、一个拜偶像者，只要他在\Place{熙雍}有种子、在\Place{耶路撒冷}有后代，他仍然是有福的吗？任何有理智的人都不会这样说。\par

贞洁是天使的生活准则，是一切无形本体的属性。我们说这话并非贬低婚姻：绝无此意！（因为我们知道主曾以祂的临在祝福了婚姻\BibleRef{若 2:1}，我们也知道他说过：\BibleQuote{婚姻是可敬的，床第是无玷的}\BibleRef{希 13:4}），但我们知道，贞洁终究比婚姻更好。因为在美德中，如同在恶习中一样，有高低等级之分。我们知道，始祖之后的所有凡人都出自婚姻。因为始祖是贞洁的产物，而非婚姻的产物。然而独身，正如我们所说，是天使的仿效。因此，贞洁比婚姻更尊贵，正如天使高于人。但我为何说到天使？\Person{基督}自己就是贞洁的荣耀，祂不仅是\Person{圣父}所独生，无起始、无发轫、无结合，而且为了我们，由\Person{童贞玛利亚}无结合地降生成人，取了我们的形像，并在自身中显示了真实完美的贞洁。因此，尽管祂并未以法律命令我们这样做（因为正如祂所说：\BibleQuote{这话不是人人都能领受的}\BibleRef{玛 19:11}），但实际上祂以此教导我们，并赐给我们力量。因为现今贞洁在人中繁盛，这是众人都清楚的事。\par

法律所命令的生儿育女确实是好的，婚姻因淫乱也是好的，因为它消除了这些\BibleRef{格前 7:2}，并以合法的结合阻止欲望的疯狂点燃非法行为。婚姻对于那些没有自制力的人是好的，但贞洁更好，它增加灵魂的果实，并向天主献上合时的祈祷之果。\BibleQuote{婚姻是可敬的，床第是无玷的，但淫乱者和奸淫者，天主必要审判。}\BibleRef{希 13:4}\par

% structure: p0149 source=h2 target=subsection reason=relative-heading-rank; base=h2

\subsection{论割损}

割损是在法律之前、祝福之后、许诺之后赐给\Person{亚巴郎}的，作为标记，将他及其后裔和家族与他所居住的外邦人分开。\BibleRef{创 17:10}这是明显的，因为当以色列人独自在旷野度过四十年，与任何其他民族无往来时，所有在旷野出生的人都未受割损；但当\Person{若苏厄}带领他们渡过\Place{约旦}后，他们就受了割损，并且颁布了第二次割损的法律。因为在\Person{亚巴郎}时代，割损的法律被赐下，但在旷野的四十年中，它被中止了。再次，天主在渡过\Place{约旦}后第二次将割损的法律赐给\Person{若苏厄}，正如\WorkTitle{若苏厄书}中记载的：\BibleQuote{那时主对\Person{若苏厄}说：\InnerQuote{你要用尖石做石刀，召集以色列子民，第二次给他们施行割损。}}\BibleRef{苏 5:2}稍后：\BibleQuote{因为以色列子民在\Place{巴塔里斯}旷野中走了四十二年，直到所有从埃及出来能作战的人，都因不听从主的声音而未受割损灭了，主曾向他们起誓，不让他们看见主向他们的祖先所起誓要赐给他们的那流奶流蜜的美地。他们的子孙，就是主在他们之后兴起的，\Person{若苏厄}给他们行了割损；因为他们未受割损，是因为在路上没有给他们行割损。}\BibleRef{苏 5:6}所以，割损是一个标记，将\Place{以色列}与他们所居住的外邦人分开。\par

此外，它也是圣洗的预像。因为正如割损并非切除身体上有用的部分，而只是无用的多余部分；同样，藉着圣洗，我们从罪中被割损，而罪显然可以说是欲望中多余的部分，并非有益的欲望。因为任何人完全没有任何欲望、从未体验过快感的滋味，是绝不可能的。但是，快感中无用的部分，也就是无用的欲望和快感，就是罪；圣洗将我们从罪中割损出来，并给我们一个标记，就是将宝贵的十字架印在额上，不是为了将我们与外邦人分开（因为万民都领受了圣洗，并以十字架的标记受了印记），而是为了在每一民族中区分信者与不信者。因此，当真理被揭示时，割损便成了无意义的预像和阴影。所以，割损如今是多余的，且与圣洗相悖。\BibleQuote{凡是受割损的，必须遵守全部法律。}\BibleRef{迦 5:3}再者，主受了割损，为要成全法律；祂成全了全部法律，并遵守了安息日，为要成全并确立法律。\BibleRef{玛 5:17}而且，在祂受洗以后，圣神以鸽子的形状降在祂身上、显现给人类，从那时起，就开始宣讲属灵的敬拜、生活方式和天国。\par

% structure: p0152 source=h2 target=subsection reason=relative-heading-rank; base=h2

\subsection{第二十六章 论假基督。}

应该知道，假基督必定要来。因此，凡不承认天主子降生成人、原是完美天主、且成为完美的人（先前是天主）的，便是假基督。\BibleRef{若一 2:22}但在一个独特而特殊的意义上，那在末世来临的被称为假基督。首先，必须将福音传给万民，正如主所说\BibleRef{玛 24:14}，然后他要来驳斥不虔敬的犹太人。因为主对他们说：\BibleQuote{我因我父的名而来，你们却不接纳我；如果别人因自己的名而来，你们反而接纳他。}\BibleRef{若 5:43}宗徒说：\BibleQuote{因为他们没有接受爱慕真理的心，为得救恩；为此，天主给他们派遣一种强烈的错谬，叫他们相信谎言，为使一切不信真理而喜爱不义的人，都被定罪。}犹太人果然没有接受主耶稣基督——天主子、天主，却要接受那自称天主的骗子。因为那骗子将冒用天主的名字，天使教导达尼尔说这些话：\BibleQuote{他不顾他祖先的天主。}\BibleRef{达 11:37}宗徒又说：\BibleQuote{不要让人用任何方式欺骗你们，因为那日子以前，必先有叛离之事，并出现那无法无天的人，即丧亡之子：他敌对并高举自己在各种称为神或受人崇拜者以上，以致坐在天主的殿中，宣布自己是天主；}\BibleRef{得后 2:3}他说\Quote{坐在天主的殿中}，不是指我们的圣殿，而是指古代的犹太圣殿。因为他不是来到我们这里，而是来到犹太人那里；不是为了基督或基督的事，因此他被称为假基督。\par

因此，首先必须将福音传给万民。\BibleRef{玛 25:14}\BibleQuote{然后那恶者将要显露出来，就是那照着撒殚的运作而出现的人，具有一切异能、奇迹和欺骗的奇事，并以各种不义的诈术，在那些丧亡的人身上，因为主将用祂口中的气息消灭他，并以祂来临的光辉摧毁他。}魔鬼自己并不会像主降生成人那样成为人。断乎不可！但他是作为淫乱的产物而成为人，并接受了撒殚的全部能量。因为天主预知他将要做出的奇异选择，就允许魔鬼住在他内。\par

因此，正如我们所说，他是淫乱的产物，秘密地长大，突然之间崛起反叛并夺取统治。在他的统治（或更确切地说，暴政）之初，他假装圣洁。但当他成为主人后，便迫害天主的教会，并显露出他全部的邪恶。但他将\BibleQuote{以奇迹和欺骗的奇事}\BibleRef{得后 2:9}（虚假而不真实的奇迹）来到，并欺骗和引诱那些心思不稳、建立在动摇基础上的人离开永生天主，甚至使选民（如果可能）也要被迷惑。\BibleRef{玛 24:24}\par

但是，哈诺客和特斯彼人厄里亚将被派遣，他们要将父亲的心转向儿女，也就是说，将犹太人的会堂转向我们的主耶稣基督和宗徒的宣讲；他们将被假基督毁灭。主将从天上降临，正如圣宗徒们看见祂升天时一样，是完美的天主和完美的人，带着荣耀和权能，用祂口中的气息消灭那无法无天的人、丧亡之子。\BibleRef{宗 1:11}因此，没有人应当期待主从地上来，而是从天上，正如祂自己所确定的。\BibleRef{得后 2:8}\par

% structure: p0157 source=h2 target=subsection reason=relative-heading-rank; base=h2

\subsection{第二十七章 论复活。}

我们也信死者的复活。因为确确实实将有死者的复活，而我们所说的复活是指身体的复活。\BibleRef{格前 15:35}因为复活是那已坠落之物的第二次状态。灵魂是不死不灭的，那么它们又怎能再次复活呢？因为如果他们将死亡定义为灵魂与身体的分开，那么复活无疑是灵魂与身体的重新结合，是那已遭受分解和坠落之生灵的第二次状态。因此，正是这个可朽坏的、易于分解的身体，将要起来成为不朽坏的。因为那在起初用尘土造它的天主，并不缺乏能力，当它再次分解并回到它所由来的尘土中时，按照造物主判决的翻转，再次将它复活起来。\par

因为如果没有复活，就让我们吃吃喝喝吧：让我们追求享乐和欢愉的生活。如果没有复活，我们与无理性的禽兽有什么区别？如果没有复活，就让我们以为田野的野兽是幸福的，因为它们过着无忧无虑的生活。如果没有复活，也就没有天主，没有眷顾，而是万物自行驱动和承载。因为请看，我们看见最正义的人在现世忍受饥饿和不义，得不到任何帮助，而罪人和不义之人却富有钱财和一切享乐。有理智的人谁会认为这是正义审判或明智眷顾的作为？因此，必须有复活，必须有复活。因为天主义，祂是那些耐心服从祂的人的赏报者。因此，如果只是灵魂独自参与德行的竞赛，那么也就只是灵魂独自获得冠冕。如果只是灵魂独自沉溺于快乐，那么也就只是灵魂独自受到公正的惩罚。但是，既然灵魂并非脱离身体而追求德行或恶行，那么两者将一同得到应有的报应。\par

不仅如此，圣经也见证身体将会复活。天主在洪水后确实对梅瑟说：\Quote{如青菜一样，我把一切赐给了你们。但是，凡有生命带血的肉，你们不可吃。并且，我要追讨你们生命的血，向一切野兽追讨，向人的弟兄追讨人的生命。凡流人血的，他的血也要被人流，因为人是按天主的肖像造的}。祂将如何向每只野兽追讨人的血？除非是因为死人的身体将要复活？因为野兽不会为人而死。\par

再次，对梅瑟说：\Quote{我是亚巴郎的天主，依撒格的天主，雅各伯的天主：天主不是死人的天主}（即那些死了且不再存在的人），\Quote{而是活人的天主}，他们的灵魂确实活在天主手中\BibleRef{智 3:1}，但他们的身体将通过复活再次获得生命。达味，神圣的先祖，对天主说：\Quote{祢收回他们的气息，他们就死亡，归于尘土}。请看他是如何谈论身体的。然后他补充道：\Quote{祢发出祢的神，它们便被创造，祢使地面更新}。\par

此外，依撒意亚说：\Quote{死人要复活，在坟墓里的人要苏醒。}\BibleRef{依 26:18}显然，灵魂并不躺在坟墓里，而是身体躺在那里。\par

再者，真福厄则克耳说：\Quote{当我正说预言时，忽然有响声，骨头彼此连接，各归其位：我观看，见骨头上生了筋，肉也长起来，皮肤在上面遮盖着。}\BibleRef{则 37:7}后来他教导说，当命令发出时，气息如何归来。\par

神圣的达尼尔也说：\Quote{那时，守护你百姓的伟大首领米额尔必起来：必有一段灾难的时期，自从有国以来直到那时从未有过。那时，你百姓中凡写在书上的必得拯救。许多睡在尘土中的人要醒来：有些得到永生，有些受到羞辱和永远憎恶。智慧者将发光如穹苍的光辉，义人中的多数将发光如星辰，直到永远和永远。}\Quote{许多睡在尘土中的人要醒来}这句话清楚地表明将有身体的复活。因为没有人会说灵魂睡在尘土中。\par

此外，主在圣福音中也明确承认身体的复活。\Quote{因为在坟墓里的人，祂说，要听见祂的声音而出来：行善的复活得生命，作恶的复活受审判。}\BibleRef{若 5:28}有理智的人绝不会说灵魂在坟墓里。\par

但是，主不仅以言语，也以行为揭示了身体的复活。首先，祂复活了拉匝禄，尽管他已经死了四天，并且发臭了。\BibleRef{若 11:39}因为祂不是复活没有身体的灵魂，而是身体连同灵魂一起复活；而且不是另一个身体，正是那个已经腐烂的身体。因为如果不是通过其特有的标记，死人的复活如何能被认识或相信？但事实上，正是为了彰显祂本性的神性，并坚定对我们自己和祂自己的复活的信仰，祂复活了最终会再次死去的拉匝禄。主自己成为了完美复活、不再受死亡支配的初果。因此，神圣的宗徒保禄也说：\Quote{如果死人不会复活，基督也就没有复活。如果基督没有复活，我们的信仰就是虚妄的，我们仍在罪中。}\BibleRef{格前 15:16}又说：\Quote{如今基督已从死人中复活，成为睡了之人初熟的果子}\Quote{，是死者中的首生者}\BibleRef{哥 1:18}；又：\Quote{因为我们若信耶稣死了，也复活了，照样，那些在耶稣内睡了的人，天主也必把他们与耶稣一同带来。}\BibleRef{得前 4:14}\Quote{照样}，他说，\Quote{正如基督复活了。}此外，主的复活是不朽的身体与灵魂的结合（因为正是这两者分离了），这是显然的：因为祂说：\Quote{你们拆毁这圣殿，三天内我要把它重建起来。}\BibleRef{若 2:19}而圣福音是可信的见证，证明祂说的是自己的身体。\Quote{你们摸摸我，看看}，主对自己的门徒说，因为他们以为自己看见了一个鬼魂：\Quote{是我自己，我没有改变}\BibleRef{路 24:37}\Quote{：因为鬼神没有肉和骨头，如同你们看见我有的。}说了这话，祂就把手和肋旁给他们看，并伸出来让多默触摸。\BibleRef{若 20:27}这难道不足以建立对身体复活的信仰吗？\par

再者，神圣的宗徒说：\BibleQuote{因为这必朽坏的必须穿上不朽坏的，这必死的必须穿上不死的}。\BibleRef{格前 15:35}　又说：\BibleQuote{它播种在朽坏中，复活在不朽坏中；播种在软弱中，复活在强能中；播种在羞辱中，复活在光荣中；播种的是属血气的身体}（即粗糙且可朽的），\BibleQuote{复活的是属神的身体} \BibleRef{格前 15:42}，就像我们主复活后的身体，穿过关闭的门，不知疲倦，不需食物、睡眠或饮料。\BibleQuote{他们要}，主说，\BibleQuote{像天使} \BibleRef{谷 12:25}；将不再有婚姻或生儿育女。确然，神圣的宗徒说：\BibleQuote{我们的家乡在天上，从那里我们也等待救主，主耶稣基督，祂要改变我们卑贱的身体，使之相似祂光荣的身体} \BibleRef{斐 3:20}：这并不是指改变成另一种形状（天主不许！），而是指从朽坏变成不朽坏。\par

但有人会说：\BibleQuote{死人怎样复活呢？}哦，何等不信！哦，何等愚妄！祂凭自己的旨意将泥土变成身体，命令那微小的种子在母腹中生长，最终成为这变化多端的身体器官，难道祂不是更凭自己的旨意使那已经并正在分解的复活起来吗？\BibleQuote{他们带着什么身体来呢？愚昧人哪} \BibleRef{格前 15:35}，如果你的刚硬不让你相信天主的话语，至少相信祂的作为。\BibleQuote{你所种的，若不先死，就不能得生；并且你所种的，不是那将来要成的身体，不过是子粒，即如麦子，或是别的谷物。但天主随自己的心意给它一个身体，并给各种子各自己的身体} \BibleRef{格前 15:35}。因此，看哪，种子是如何像葬在坟墓里一样埋在犁沟中。是谁赐给它们根、茎、叶、穗和最纤细的芒？难道不是宇宙的创造者？难道不是那筹划万有的命令？因此，这样相信吧：死人的复活必因天主的意愿和号令而成就。因为祂的能力足以与祂的旨意相配。\par

我们将因此复活，灵魂再次与身体结合，身体如今变得不朽坏，脱去了朽坏，我们将站在基督可畏的审判宝座前；那魔鬼及其恶魔，以及属魔鬼的人，即假基督、不虔敬者和罪人，将被投入永火：不是物质火像我们的火，而是天主所知道的火。但那行善的人将如同太阳发出光辉，与天使一同进入永生，与我们的主耶稣基督同在，永远看见祂、在祂面前，从祂那里获得不断的喜乐，同圣父及圣神赞美祂，直到无穷无尽的世世代代。阿们。\par

\SourceRecord{Translated by E.W. Watson and L. Pullan.；Nicene and Post-Nicene Fathers, Second Series；Vol. 9.；Edited by Philip Schaff and Henry Wace.；Buffalo, NY: Christian Literature Publishing Co.,；1899.；<http://www.newadvent.org/fathers/33044.htm>.}
